\documentclass[11pt,a4paper]{article}

% -------------------------------------------------------
% Page layout and typography
% -------------------------------------------------------
\usepackage[margin=2.5cm]{geometry}
\usepackage[T1]{fontenc}
\usepackage[utf8]{inputenc}
\usepackage{lmodern}
\usepackage{microtype}
\usepackage{setspace}
\onehalfspacing

% -------------------------------------------------------
% Tables, figures and maths
% -------------------------------------------------------
\usepackage{graphicx}
\usepackage{float}
\usepackage{booktabs}
\usepackage{tabularx}
\usepackage{longtable}
\usepackage{xltabular}
\setlength{\LTcapwidth}{\textwidth}
\usepackage{array}
\usepackage{amsmath}
\usepackage{siunitx}
\usepackage{gensymb}

% -------------------------------------------------------
% TikZ figures
% -------------------------------------------------------
\usepackage{tikz}
\usetikzlibrary{
    arrows.meta,
    positioning,
    shapes.geometric,
    shapes.misc,
    calc,
    fit,
    backgrounds
}

% -------------------------------------------------------
% References and links
% -------------------------------------------------------
\usepackage[round,authoryear]{natbib}
\usepackage[hidelinks]{hyperref}
\usepackage[nameinlink,noabbrev]{cleveref}

% -------------------------------------------------------
% Metadata
% -------------------------------------------------------
\title{
What Makes Mountain Wildfire Different?\\
Evidence from the European Alps and Other Ranges
}

\author{
% DRAFT author list: order, institutes and e-mail to be confirmed by the author team
Alexandre Dunant$^{1,*}$, Alice Crespi$^{1}$, Michael de Rachewiltz$^{1}$, Christian Hoffmann$^{1}$,\\
Marco Mina$^{1}$, Massimiliano Pittore$^{1}$, Uta Schirpke$^{1}$, Stefano Terzi$^{1}$,\\
Jenny Ufer$^{1}$, Marc Zebisch$^{1}$ \\
\small $^{1}$Eurac Research, Bolzano, Italy \\
\small $^{*}$Corresponding author: [e-mail to be added]
}

\date{}

% -------------------------------------------------------
% Begin document
% -------------------------------------------------------
\begin{document}

\maketitle


\begin{abstract}
\noindent
Wildfire has been treated as a Mediterranean and lowland problem, and the European Alps have sat at the margin of European fire science. That position is no longer tenable: across European mountains the elevation of large fires has risen by roughly 72\,m per decade since 2000, danger is projected to rise across the arc under every emissions pathway, and land abandonment has assembled fuel continuity on steep ground that was pasture a century ago. Setting the Alpine case, with comparative evidence from other European and non-European mountain ranges, against the non-mountain baseline in each of hazard, climate, water, forest, risk and governance, we argue that mountain wildfire is not a scaled-down Mediterranean problem but a structurally different one. The fire-size distribution is heavy-tailed, around 80\% of fires burning under a hectare, and fires act on forests managed to hold rock, snow and sediment above settlement and infrastructure, so consequence is decoupled from area burnt. Consequence follows several pathways: loss of protective structure, a geomorphic cascade in which fire lowers the rainfall threshold for erosion, debris flow and rockfall over a window lasting from three years to, at subalpine elevations, effectively permanent, and an ecological loss whose sign differs between closed-forest and open-habitat species and whose faunal term is unmeasured. Because the asset at risk is protective function rather than timber, the response is substantially a forestry problem, and management capacity is a first-order determinant rather than a residual after climate. We show why the linear prevention instruments of lowland practice transfer awkwardly to terrain where a fuel break is a platform for ground suppression that crews cannot reach, and why an apparatus producing hazard surfaces without exposure or vulnerability terms cannot yet produce risk. A research agenda follows from what the observational record can support.
\end{abstract}

\noindent\textbf{Keywords:}
mountain wildfire; European Alps; mountain and lowland fire regimes; fire regime change; protection forest; fuel management; drought and water availability; post-fire hazards; wildfire governance and economics; risk assessment


% =======================================================
% =======================================================
\section{Introduction}
\label{sec:intro}

\subsection{An emerging fire frontier in temperate mountains}

European fire science has been shaped by where burnt area is concentrated. Because the Mediterranean basin accounts for most of the continent's burning, it has also supplied most of the continent's fire statistics, operational danger rating and prevention doctrine. Where fire occurs is now changing. Burnt area has been declining in the Mediterranean while increasing in mid- and high-latitude regions where fuel is abundant and plant communities have little adaptation to fire \citep{miller_compounding_2025}. The shift is vertical as well as latitudinal. Across eight European mountain regions, the median elevation of fires burning at least 30\,ha above 800\,m rose from 910\,m in 2000--2014 to 1004\,m in 2015--2025, and the ninetieth percentile from 1353 to 1460\,m. This upslope shift of approximately 72\,m per decade reflects a change after 2015 rather than a gradual trend, and burnt area in these regions was most strongly predicted by vapour pressure deficit \citep{beloiu_climatedriven_2026}.

The pattern is not confined to mountains. In the Netherlands, a densely populated temperate country with no fire tradition, 3{,}667 landscape fires were recorded between 2017 and 2022, averaging 611 a year and around 405\,ha annually. Mean fire size was 1.5\,ha, 88\% of events burned less than a hectare, and almost all were human-ignited \citep{stoof_temperate_2024}. This profile is close to the Alpine one, and so is the institutional setting: fire services specialise in urban firefighting, and neighbouring countries hold no systematic wildfire statistics at all.

Projected increases for the Alps are large relative to this baseline. Mean hazard probability across the Alpine Space cooperation area, which spans the Alps and their surrounding lowlands, is projected to increase by 15.5\% under SSP1-2.6 and by 35.4\% under SSP5-8.5 by 2070--2100. The largest relative increases are in the northern and eastern sectors, not in the already fire-affected south-west \citep{gerberding_mapping_2025}. For adjacent central Europe, the forced fire-danger signal is projected to emerge before 2050 \citep{miller_climate_2024}. Both projections suggest that fire may become a routine hazard within the planning horizon of the institutions that would have to absorb it (Section~\ref{sec:climate}).

Climate alone does not fix how much of this increase is realised. One assessment couples two fire models, SPITFIRE and the Burnt Area Simulator for Europe, to the fire-enabled dynamic global vegetation model LPJmL. In these simulations, intensified fire weather alone is projected to raise annual burned area by around 39\% under SSP1-2.6 and by nearly 192\% under SSP3-7.0 by the end of the century. Continued improvement in management capacity would reduce burned area by 72\% to 92\% relative to scenarios without it. Net increases would still remain in roughly 55\% of Europe's fire-prone regions. Under the higher pathway, the hotspots are concentrated in foothills and medium-elevation terrain, including the Alps, Pyrenees, Carpathians and Massif Central \citep{billing_future_2026}. Management capacity is therefore a primary determinant of Alpine outcomes, not a residual after the climate signal has been counted.

Mitigation also shapes the outcome. By the accounting of the European Forest Fire Information System, burnt area across Europe in 2025 was the largest since its records began in 2006, with the largest contributions from the Iberian Peninsula \citep{burton_state_2026}. The same assessment, still a preprint, finds that emissions mitigation has already reduced projected end-of-century burnt area in extreme events in its focal regions, which include north-western Iberia. Relative to a worst case now considered implausible, the reduction is 4\% to 23\% under a high pathway, with a further 5\% to 36\% under a Paris-consistent pathway. Alpine fire danger therefore depends on mitigation as well as on adaptation, although it is rarely presented this way.

Fire is thus arriving in terrain where the science, the management system and the exposed population have little experience of it. Alpine forestry identified this gap more than a decade ago. Reviewing the risks facing Alpine Space forests, the MANFRED consortium observed that, because of a history of active suppression, fire risk had not received full attention in the Alps. It suggested that future climate conditions might well require an improved alerting system \citep{jandl_management_2013}. The climate signal has since emerged, but attention to fire has followed more slowly.

\subsection{Why mountain fire is a distinct problem}

Tools developed for the Mediterranean do not transfer cleanly to mountains, and the problems begin with the atmosphere. These tools assume a continuous fire-weather field that can be interpolated from valley-floor stations. That assumption breaks down where upper winds interact with complex relief to produce discontinuous and locally amplified surface conditions \citep{sharples_overview_2009}. Evidence imported from one Tasmanian fire shows wind profiles that enhance surface winds becoming several times more frequent as terrain grows more rugged. On that rugged terrain, fire simulations driven by surface winds alone were outperformed by those that included upper-air winds (Section~\ref{sec:climate}) \citep{ozaki_impact_2022}.

Relief also changes how fire behaves. In controlled experiments, rate of spread accelerates non-linearly with upslope angle, and the dominant preheating mechanism switches from radiation to convection as the slope steepens \citep{hayajneh_wind_2026}; Section~\ref{sec:hazard} gives the thresholds. The effect is measurable at landscape scale. Across 21 Swiss pyroregions, the length of upslope paths within forested area, together with the aggregation of that area, explains between 56\% and 65\% of the variance in large-fire statistics without any weather term \citep{conedera_linking_2024}.

The larger difference, however, lies in what happens after the fire is out, because in mountains the fire is rarely the end of the event. Much Alpine forest is designated for protection against rockfall, avalanche and shallow landslide, and in South Tyrol about 58\% of the forest area is classed as protection forest \citep{hillebrand_contrasting_2023}. Fire can remove that function for several years and, depending on forest type and fire intensity, for decades \citep{melzner_wildfire-induced_2025, melzner_brief_2019}. The asset at risk is typically not the stand that burns but the settlement, road or railway below it. These differences recur throughout the review. Together they explain why burnt area, a statistic that serves the Mediterranean reasonably well, says little about the consequences of fire in mountains.

\subsection{Objective, argument and the structure of this review}

This review asks what makes wildfire in mountains different from wildfire in the lowland and Mediterranean settings that dominate European fire science. The European Alps are the focal case, and evidence from other ranges shows whether a difference belongs to mountains in general or to the Alps alone (Section~\ref{sec:scope}). Our aim is to assemble the evidence needed to treat mountain wildfire as a systemic risk rather than as a forest-loss problem.

The Alpine fire community already has a consensus statement. The EUSALP white paper of February 2020 set out the state of knowledge for policy makers and established the arc-wide statistical baseline. It identified the loss of protective function as the central concern and proposed a costed integrated fire management framework \citep{muller_forest_2020}. We take that position as our starting point and ask what the evidence has established since. Part of our argument extends the 2020 position and part complicates it. We keep the two apart, and name the finding responsible wherever the evidence complicates the 2020 position, because an explicit account of what is inherited is more useful than implied originality.

Comparing the evidence with the 2020 statement serves a further purpose, because the statement is a policy white paper rather than a contribution to the scientific record. Its evidence base is not indexed, peer reviewed or available for citation and appraisal in the ordinary way, so the community that must act on its positions cannot readily test or extend them. A review can supply part of what is missing by setting that position against the peer-reviewed literature that has accumulated since and stating for each element how securely it is now supported.

The extension builds on framings set out in 2020. The white paper argued from arc-wide statistics that the Alpine fire regime is structurally distinct, not merely smaller than the regimes on which European fire science was built. Later work has quantified this in the vertical dimension, through a measured upslope migration of fire activity in European mountains \citep{beloiu_climatedriven_2026}, and in the operational dimension, through terrain-conditioned estimates of containment difficulty \citep{cardil_assessing_2025}. The white paper also recognised that the response is substantially a forestry problem, because what fire destroys in the Alps is a managed protective function. The material for acting on this has been added since. We therefore treat silviculture, fuel treatment and post-fire restoration as central subjects, drawing on monetised valuations of the protective service \citep{accastello_asforesee_2019, moos_climate_2021}, a controlled Alpine prescribed-burning trial \citep{gamba_prescribed_2025} and a conceptual account, grounded in burn-severity mapping of recent Alpine fires, of how the duration of protection loss depends on fire intensity and forest type \citep{melzner_wildfire-induced_2025}.

The complications form the core of the review, and the most far-reaching concerns consequence. In the Alps, consequence is decoupled from burnt area. It runs through several pathways at once, namely direct loss of protective structure, ecological and hydrological change, and a geomorphic cascade in which fire creates a susceptibility that a later hydro-meteorological event realises. The geomorphic pathway is the least represented in current assessment, but it is best treated as one term in the loss function rather than as the whole of it. Part of this pathway also falls outside the return-period logic in which mountain hazard is conventionally expressed. In one burnt catchment in the Susa Valley, a debris flow was triggered at a radar-estimated intensity of about 24\,mm\,h$^{-1}$, below the 30\,mm\,h$^{-1}$ that basins of its type typically require, although radar reliability is limited in such relief \citep{laporta_postwildfire_2025}.

Because burnt area does not capture consequence, neither do the assessment tools built on it. These tools carry hazard some way towards risk but stop short of it. The exposure and vulnerability terms that would convert a burn probability into an expected loss exist as instruments but have not been assembled for the Alpine arc. The Alpine data record is also too thin to justify the conventional remedy of added complexity. We use that constraint, rather than an appeal to greater model sophistication, to organise a research agenda.

We organise the material by the domains in which the evidence is produced rather than by the stages of a fire. Section~\ref{sec:hazard} treats wildfire as a natural hazard, covering initiation, spread, hazard modelling, and the datasets and Earth observation on which all of it depends. Section~\ref{sec:climate} covers climatic and meteorological drivers, including the wind systems, lightning climatology and projections specific to complex terrain. Section~\ref{sec:water} follows the water pathway, from drought and fuel moisture through soil alteration and the geomorphic cascade to water supply, snow and ice. Section~\ref{sec:forest} treats recent forest change, ecology and management. Section~\ref{sec:risk} assembles the exposure and vulnerability terms that convert hazard into risk, and Section~\ref{sec:governance} treats governance, policy and economics. Each of these sections closes with a comparison between the mountain case and the lowland and Mediterranean baseline, collected in Table~\ref{tab:contrast}. We close by considering what the observational record can support.

% =======================================================
\section{Review Design and Scope}
\label{sec:scope}

\subsection{Scope and inclusion rule}

The review therefore compares mountain fire with the settings on which European fire science was built. The European Alps are the focal case, because they have the densest and most recent evidence and because the practitioner and policy community is organised at that scale. The claim being tested, however, is about mountains, so we built the corpus in two tiers rather than as one core with a set of exceptions.

Tier one is the Alpine arc as delimited by the Alpine Convention, covering Austria, France, Germany, Italy, Liechtenstein, Monaco, Slovenia and Switzerland. Harmonised Alpine fire statistics were first assembled within this perimeter \citep{valese_wildfires_2011, muller_forest_2020}, arc-scale hazard mapping has been carried out within it \citep{gerberding_mapping_2025}, and fire practitioners across the related EUSALP perimeter have been surveyed jointly \citep{hoang_stakeholders_2026}.

Tier two is mountain fire systems in general, European and otherwise, admitted because they share the processes under review, not as exceptions. Results from the Pyrenees, Carpathians, Apennines, Scandes, Rocky Mountains or Sierra Nevada enter wherever they bear on a process the review treats, with or without Alpine evidence on that process. Where both are available, we compare them. Agreement across ranges indicates that a feature belongs to mountains generally, and disagreement shows what is specific to the Alps. Table~\ref{tab:ranges} describes the systems drawn on and what each contributes.

\begin{table}[H]
\centering
\footnotesize
\setlength{\tabcolsep}{4pt}
\renewcommand{\arraystretch}{1.1}
\caption{Mountain systems drawn on in this review, with the regime feature that makes each informative and the contribution it makes to the mountain-versus-lowland comparison.}
\label{tab:ranges}
\begin{tabularx}{\textwidth}{@{}>{\raggedright\arraybackslash}p{2.0cm} >{\raggedright\arraybackslash}p{3.3cm} >{\raggedright\arraybackslash}X >{\raggedright\arraybackslash}p{3.1cm}@{}}
\toprule
\textbf{Range} & \textbf{Regime feature} & \textbf{Contribution to the argument} & \textbf{Main references} \\
\midrule
European Alps & Bimodal, spring anthropogenic and summer lightning seasons; most fires under a hectare; protective forest above fixed assets & Focal case: consequence decoupled from burnt area, the geomorphic cascade, protective-function loss, and the thinness of the observational record & \citealp{muller_forest_2020, conedera_characterizing_2018, moos_climate_2021, melzner_wildfire-induced_2025} \\
Pyrenees (Aran valley) & Winter-dominant fire season, Atlantic rather than Mediterranean, transhumance history, deep organic soils & A Pyrenean valley that has operationalised managed fire in a formally approved strategic plan, so the Alpine prescribed-fire question has a European mountain precedent & \citealp{oliveres_path_2025} \\
Carpathians & Sharp spring peak in March, cropland the leading susceptibility predictor, lower elevations more fire-prone & Shows that the elevation relationship is not fixed across mountains, and that a spring anthropogenic peak recurs outside the Alps & \citealp{manczinger_fire_2025} \\
Apennines & Long land-use history recorded palaeoecologically; monospecific beech replacing mixed fir & Independent European mountain evidence that present forest composition, and therefore the fuel bed, is a land-use artefact & \citealp{moralesmolino_longterm_2021} \\
Dinarides and Tatras & Long-term fire and land-use history on karst (Dinarides); windthrow and salvage logging (Tatras) & Fire history on a karst massif, and the salvage question after large windthrow, a non-fire analogue & \citealp{cagliero_effect_2022, seben_impact_2024} \\
Scandes and Fennoscandia & Very small fires, sparse satellite detection & Quantifies the gap between satellite and national fire records in northern Europe (Norway, 2016--2019), without separating mountain terrain & \citealp{bakke_datadriven_2023} \\
Rocky Mountains & Frequent post-fire debris flow with operational warning & The most developed post-fire cascade science, supplying thresholds and volume models against which Alpine observations can be placed (tested at one fire in Glenwood Canyon, Colorado); terrain control of persistent fire refugia (Canadian Rockies) & \citealp{rengers_evaluating_2024, rogeau_effect_2018} \\
Sierra Nevada and Cascades & Long fire-weather records; measured post-fire snow response & The fire and snow coupling measured directly, which the Alps lack & \citealp{micheletty_application_2014, koshkin_wildfire_2022} \\
Southern Appalachians & Recent fire outbreak after long suppression; temperate broadleaf & Soil and microbial endpoints, and a bound on how far the suppression-legacy narrative transfers & \citealp{reilly_drivers_2022, snyder_taxonomic_2026} \\
Australian Alps and Tasmania & Alpine treeline vegetation; rugged terrain; three-jurisdiction parks & Terrain control on surface winds, alpine ecology after fire, and standing cross-border coordination & \citealp{sharples_overview_2009, ozaki_impact_2022, wilson_distribution_2022, lockwood_biodiversity_2014} \\
East Asian mountains & Boreal and subtropical mountain forest & Slope effects on sub-surface smouldering, and diurnal valley circulation resolved by coupled modelling & \citealp{shan_influence_2024, wang_dynamic_2026} \\
\bottomrule
\end{tabularx}
\end{table}

Drawing on other ranges requires care about what transfers. We treat mechanisms as transferable where the physics is general but do not transfer coefficients, and we mark the distinction where each result is used. A result from outside the Alps therefore establishes that a process operates, bounds its magnitude, or shows whether an Alpine observation is exceptional. Alpine parameters are estimated from Alpine data alone.

The breadth is justified because none of the processes involved is confined to the Alps. Terrain reorganises atmospheric flow and therefore fire weather, and slope and wind interact non-linearly to set the rate of spread. Fire removes the vegetation and soil structure that hold sediment on steep ground, and the assets exposed in mountains lie below the fire rather than within it. A third tier, used as context rather than cited as evidence, covers mountain processes without a fire component and European fire without a mountain component.

\subsection{The six-step review framework and evidence weighting}

We assembled and reviewed the literature in six steps (Figure~\ref{fig:flow}). Appendix~\ref{app:protocol} gives the numbers at each step and the reasons for excluding publications.

\textbf{Step 1, starting documents.} We started from the 2020 EUSALP white paper on forest fires in the Alps \citep{muller_forest_2020}, the most recent synthesis for the whole Alpine arc, whose conclusions this review sets out to extend and test (Section~\ref{sec:intro}). Alongside it we used a mapping of forest fire hazard across the Alps under climate change scenarios, co-authored at Eurac Research \citep{gerberding_mapping_2025}, which adds the projected, spatially explicit view of hazard that the white paper does not cover.

\textbf{Step 2, keyword search.} We then searched Google Scholar by hand, first with the terms wildfire, mountain and Alps, and then with wildfire and mountains to include ranges outside the Alps. We kept publications that fell within one of the disciplines represented at Eurac Research, which are the disciplines of Sections~\ref{sec:hazard} to~\ref{sec:governance}, so that a specialist could later assess each of them. Together with the two starting documents, this gave an initial set of 49 publications.

\textbf{Step 3, citation chaining.} We expanded this set through its citations in both directions. Working backward, we went through the reference lists of the white paper and of the initial publications. Working forward, we looked up later publications that cite the most central Alpine ones.

\textbf{Step 4, targeted additions.} While drafting the sections, we searched specifically for publications on topics that the literature gathered so far covered poorly.

\textbf{Step 5, screening.} We screened every publication against an inclusion rule with four levels: fire evidence from the Alps, fire evidence from other mountain systems, context, and exclusion. Each decision was recorded with its reason. Every source cited in the review was read and summarised in a structured evidence note.

\textbf{Step 6, discipline review.} Each thematic section is reviewed by the co-authors from Eurac Research who work in that discipline. Climatologists review the climate section and forest scientists the forest section, while specialists in drought and water, natural hazards and risk, and economics and governance review the others. They check that their field is covered, that each source is read correctly and that the text gives it appropriate weight. Publications they add return to Step~3, so that their references are followed and screened in the same way as the rest of the literature.

\begin{figure}[H]
\centering
\begin{tikzpicture}[
  scale=0.8, transform shape,
  font=\small,
  stp/.style={draw, rounded corners=2pt, align=left, text width=6.4cm, inner sep=5pt, fill=black!6},
  res/.style={draw, rounded corners=2pt, align=left, text width=6.2cm, inner sep=5pt},
  arr/.style={-{Latex[length=2mm]}, thick},
  lnk/.style={-{Latex[length=2mm]}, thin}
]
\node[stp] (s1) {\textbf{Step 1. Starting documents}\\ EUSALP white paper on forest fires in the Alps \citep{muller_forest_2020}\\ Mapping of Alpine forest fire hazard under climate change \citep{gerberding_mapping_2025}};
\node[stp, below=5mm of s1] (s2) {\textbf{Step 2. Keyword search}\\ Google Scholar: \emph{wildfire} + \emph{mountain} + \emph{Alps}, then \emph{wildfire} + \emph{mountains}; publications kept where a Eurac Research discipline can assess them};
\node[stp, below=5mm of s2] (s3) {\textbf{Step 3. Citation chaining}\\ Backward through the reference lists of the white paper and the initial set; forward to later publications citing the core Alpine ones};
\node[stp, below=5mm of s3] (s4) {\textbf{Step 4. Targeted additions}\\ Specific searches on topics found to be thinly covered while drafting};
\node[stp, below=5mm of s4] (s5) {\textbf{Step 5. Screening}\\ Four-level inclusion rule (Alps, other mountain systems, context, exclusion), with the reason for each decision recorded};
\node[stp, below=5mm of s5, dashed] (s6) {\textbf{Step 6. Discipline review}\\ Each thematic section reviewed by the Eurac Research co-authors working in that discipline};

\node[res, right=8mm of s1] (o1) {Arc-wide synthesis of the Alpine fire regime and its management, with 120 references; projected fire hazard for the whole Alpine arc};
\node[res, right=8mm of s2] (o2) {Initial set of 49 publications};
\node[res, right=8mm of s3] (o3) {3{,}375 publications from reference lists and 319 citing publications, passed on to screening};
\node[res, right=8mm of s4] (o4) {Further publications on topics in Sections 3--8 that were thinly covered};
\node[res, right=8mm of s5] (o5) {3{,}228 publications screened; 896 retained (290 Alps, 194 other mountain systems, 412 context)\\ 150 sources cited, each summarised in an evidence note};
\node[res, right=8mm of s6, dashed] (o6) {Climate (Sec.~4), drought and water (Sec.~5), forest (Sec.~6), natural hazards and risk (Secs.~3 and 7), economics and governance (Sec.~8)};

\foreach \i/\j in {s1/s2, s2/s3, s3/s4, s4/s5, s5/s6} \draw[arr] (\i) -- (\j);
\foreach \i/\j in {s1/o1, s2/o2, s3/o3, s4/o4, s5/o5, s6/o6} \draw[lnk] (\i) -- (\j);
\draw[arr, dashed] (s6.west) -- ++(-6mm,0) |- node[pos=0.25, rotate=90, above, font=\footnotesize] {publications added by reviewers} (s3.west);
\node[font=\footnotesize\bfseries, anchor=south] at ([yshift=1mm]s1.north) {Step};
\node[font=\footnotesize\bfseries, anchor=south] at ([yshift=1mm]o1.center |- s1.north) {Result};
\end{tikzpicture}
\caption{The six steps through which the literature was assembled and reviewed. The left column describes each step and the right column its result. Step~6, outlined dashed, is the review by discipline that the manuscript is now undergoing; publications added by the reviewers return to Step~3 (dashed arrow).}
\label{fig:flow}
\end{figure}

Our synthesis is thematic rather than meta-analytic, because the evidence does not support pooling. Alpine fire studies differ in their definitions of severity, burnt-area thresholds, reporting periods and the administrative basis of their fire records, so their effect sizes are not comparable. The same obstacle has been documented for the neighbouring literature on land abandonment and disturbance, where study sites cluster strongly and coverage across disturbance types is uneven \citep{mantero_influence_2020}.

Much of the Alpine evidence consists of individual documented events, so we weight it by an explicit convention. We use case evidence to show that a process operates, to describe its mechanism and to bound its magnitude in one setting, and we reserve regional parameters for multi-site evidence. Wherever a claim rests on fewer than three independent cases, we state their number. A mechanism established outside Europe is marked as imported and used for mechanism only. Where multiple independent Alpine studies support a claim, we say so. A single-case magnitude should therefore be read as evidence that the process can reach that value, not as a central estimate. Several of the quantities cited in Sections~\ref{sec:water} and~\ref{sec:forest} are of this kind.

% =======================================================
\section{Wildfire as a Natural Hazard}
\label{sec:hazard}

\subsection{Initiation: the Alpine ignition regime and its modelling}

An Alpine ignition model must reproduce the shape of the fire regime more than its magnitude. Harmonised statistics for 2006--2017 place annual fire numbers between 37 in Slovenia and 1043 in Italy, and annual burnt area between 64\,ha in Austria and 9984\,ha in Italy \citep{muller_forest_2020}. Normalised by forest area, burnt area spans two orders of magnitude, from 0.22 to 44.13\,ha per 10{,}000\,ha of forest per year. Roughly 80\% of Alpine fires burn less than one hectare, and median fire size falls below 0.1\,ha in Austria, Germany and Switzerland. The distribution is nonetheless strongly heavy-tailed. In the French and Italian Alps the mean fire covers about a hectare, yet the 10\% of fires exceeding ten hectares account for some 85\% of total burnt area. Both parts of this statistic matter, and either one alone gives a misleading picture. The many small fires set what monitoring must detect and what suppression routinely handles. The few large ones account for most of the burnt area and, in the Alps, are the events that can remove protective function across a whole slope. For comparison, the mean patch size of all detected forest disturbance across Europe over three decades is 1.09\,ha \citep{senf_mapping_2021}, so Alpine fire sits at the fine end of an already fine-grained disturbance regime.

Seasonality is bimodal, and we consider this the most consequential structural feature of the regime. A spring season of fast-spreading surface fires at low elevation peaks in March and April, and a summer season peaks in July and August. Winter fires are also documented \citep{valese_wildfires_2011}, particularly in the foehn regions north and south of the Alps, where a thermal belt forms above the fog line \citep{muller_forest_2020, schunk_forest_2013}. The pattern extends beyond the Alpine arc. In the Czech Republic the same bimodality appears, with an April peak attributed to forestry operations resuming after winter and a July--August peak attributed to heat, drought and tourism, with a threshold response above 18\,\degree C \citep{bercak_combination_2024}. Because the two modes have different controls, there is no reason to expect them to change in step.

The causes of ignition are mixed in a way that has no Mediterranean analogue. Human causes account for approximately 90\% of Alpine fires and lightning for around 10\% \citep{muller_forest_2020}. The lightning share reaches up to 30\% in some regions and seasons and up to half of summer fires in Austria and parts of Switzerland, with lightning fires concentrated on the southern side of the Alps and at higher elevation \citep{muller_forest_2020, valese_wildfires_2011, conedera_lightninginduced_2006}. Anthropogenic ignition follows the access network rather than population alone. At national scale in Austria, the main predictors are railway, forest-road and hiking-trail density, together with agricultural and forestry activity, and they resolve to municipality-level danger classes that can be used to target prevention \citep{arndt_modeling_2013}. The relationship with access is not linear, however. In the Valais between 1956 and 2006, fire occurrence rose slightly with road density up to about 30\,m per hectare and then fell sharply, probably because the densest networks mark urbanised or intensively farmed ground with little fuel and rapid suppression \citep{zumbrunnen_human_2012}. Relief complicates the proxy further. In steep terrain, roads follow geomorphological constraints and settlements cluster on valley floors and terraces. Distance from either therefore increases with altitude while temperature decreases, which makes any attribution of fire occurrence to such anthropogenic proxies uncertain \citep{bekar_crossregional_2020}.

The two seasonal modes can also be separated statistically, and different variables drive each. In an inner-Alpine valley, summer ignition responds to maximum temperature and to precipitation of the wettest month, while winter forest fire responds to precipitation and elevation. Proximity to infrastructure matters in winter but not in summer. Negligence in agro-silvicultural activity accounts for around two-thirds of identified ignitions, and the lightning share rose over the study period \citep{vacchiano_modeling_2018}. A single annual model fitted to this system would therefore average across processes with different drivers acting in different seasons.

The spring anthropogenic peak is not an Alpine peculiarity, but comparison with the neighbouring Carpathians shows where the resemblance ends. Susceptibility modelling across the Carpathians covered some 286{,}295\,km$^2$ in seven countries and 5{,}173 fire incidents between 2010 and 2020 \citep{manczinger_fire_2025}. It reproduces the sharp spring maximum, there peaking in March, and identifies cropland and pasture as the leading predictors, ahead of any climatic term. The elevation relationship, however, is inverted. In the Carpathians lower elevations are more fire-prone and steeper slopes are comparatively protected, which the authors attribute to easier human access on flat ground. In the Alps, by contrast, activity is migrating upslope. Two ranges with the same anthropogenic spring season therefore differ in where that season concentrates, so elevation should not be treated as a general mountain fire predictor.

Year-to-year departures from this seasonal structure are large. Between 1993 and 2013, Austria recorded three consecutive anomalous seasons in different parts of the calendar \citep{muller_anomalies_2015}. Autumn 2011 had 47 fires and 27\,ha against a mean of 11 and 3. Spring 2012 had 184 fires, the highest spring count in the series, against a mean of 57. Summer 2013 had 107\,ha burnt against a mean of 22. Across the arc, the larger anomalies occur in the eastern Alpine countries, Austria, Slovenia and northeastern Italy, rather than in the west. With interannual variance this high, a short record cannot separate trend from fluctuation, and the same study reports that roughly half of Austrian fire records lack a burnt area at all.

Despite this interannual variability, the regime divides into coherent regions \citep{conedera_characterizing_2018}. Southern, Northern and Maritime clusters separate the high fire density of the southern slope from the substantially lower burnt-area fraction of the north. The clusters also differ in drought frequency and duration, vegetation type, mean elevation, valley orientation, population density and educational level. Because the classification rests on only ten years of record, it serves as a baseline for detecting future regime shifts and anomalous seasons rather than as a description of a stable state. The same study concludes that, on present evidence, fire does not rival storms and bark beetles for damage in the Alpine region, although future regimes may evolve under climatic and socio-economic change. This framing is more measured than much of the recent literature, and we adopt it in this review.

Of the modelling domains treated here, ignition has matured fastest, because ignition points are the only Alpine fire observation available in reasonable numbers. Skill is high where models are fitted locally, and parsimony is not an obstacle. In Trentino--South Tyrol, a small set of dynamic meteorological predictors combined with static proxies for fuel, topography and human influence proved sufficient \citep{moreno_spacetime_2025}. In the Aosta Valley, static climatic, topographic and land-use predictors reached an area under the curve of 0.90 to 0.95 \citep{vacchiano_modeling_2018}. Thresholds and danger classes, rather than raw probability surfaces, have been the practical outputs \citep{moreno_spacetime_2025, arndt_modeling_2013}. What such models lose is transferability. Pooling five European regions, none of them Alpine, into a single ignition model lowered skill at every site, by up to 0.10 in area under the curve relative to site-specific fitting \citep{gelabert_assessing_2025}. This directly measures the cost of treating heterogeneous fire environments as one. The cost, however, buys a robustness that local fitting cannot provide. Occurrence models fitted separately for southern Switzerland, Carinthia, Valencia and south-western Turkey transferred well only between regions with similar environmental conditions, whereas a model trained on all four performed better than the regional models when transferred \citep{bekar_crossregional_2020}. The pooled model also drew more heavily on climate, with climatic predictors contributing about 68\% of its explanatory power against a regional mean of 49\%. Pooling therefore trades local skill for stability beyond present conditions, which is what a climate scenario requires. In the same study, resolution mattered as much as pooling. A 100\,m grid outperformed a 1\,km grid in every region, and in Carinthia the contribution of forest variables rose from about 2\% to over 40\% once the small-scale forest mosaic could be resolved.

\subsection{Spread: fire behaviour and spread modelling}

For spread, the Alpine observational record is far less developed than the process literature. Observational evidence for anomalous, terrain-driven fire behaviour in the Alps currently rests on a single documented case. During the 2018 Hallstatt fire in Austria, vertical rock walls, anabatic winds and patchy vegetation made the fire jump upward and left a spotty burn pattern \citep{melzner_brief_2019}. The result was spatially variable intensity and correspondingly heterogeneous changes in soil and rock mass. This variability matters because it carries through into burn severity and from there into post-fire hazard. Terrain-driven complexity in fire behaviour is therefore a risk-assessment problem and not only an operational one.

The process literature, by contrast, supplies both mechanism and magnitude. Wind effects on spread are strongly non-linear \citep{hayajneh_wind_2026}. Preheating shifts from radiation-led to convection-led near 1.5\,m\,s$^{-1}$, and in controlled fuels rate of spread increases by an order of magnitude by 3\,m\,s$^{-1}$. Discontinuous fuels show discrete ignition regimes, with no ignition below about 1\,m\,s$^{-1}$, direct-contact ignition between 1.5 and 2.5\,m\,s$^{-1}$, and spotting above 3\,m\,s$^{-1}$. Alpine valley winds routinely cross these thresholds within a diurnal cycle, so the regime transitions they imply are ordinary conditions rather than exotic ones.

Coupled fire-atmosphere modelling can now resolve explicitly what the Alpine literature documents case by case, as shown for a European event outside the Alps. A simulation of the 2017 Pedr\'og\~ao Grande fire at 80\,m resolution reproduced updraughts above 15\,m\,s$^{-1}$ and a convective column reaching roughly 10\,km \citep{couto_numerical_2024}. The plume top rose from below 4\,km to over 9\,km, and a localised downward outflow of similar magnitude developed ahead of the fire front. Changes in wind direction triggered acceleration and merger of fire fronts. Alpine valleys generate the terrain-forced circulations that such models are built to resolve, yet no equivalent coupled simulation has been published for an Alpine fire. This is the methodological gap between descriptive and mechanistic work on Alpine fire meteorology.

Spread modelling, unlike ignition, is the least mature Alpine domain, and the instruments in use rest on incompatible foundations. The only systematic Alpine evaluation compared two simulators on three Lombardy fires between 2017 and 2022, at Tirano, Campolaro and Sonico \citep{malanchini_modeling_2025}. FlamMap is a deterministic simulator that propagates a surface and crown spread solution across a gridded landscape, while PROPAGATOR is a stochastic cellular automaton. Both reproduced the observed propagation, but they diverged substantially on fireline intensity and rate of spread. This divergence follows from their differing assumptions rather than from calibration. The authors conclude that FlamMap suits planning because it resolves fine-scale interaction between fuel and topography, while PROPAGATOR suits operations because it is fast and accepts time-varying wind. No Alpine instrument currently does both. The disagreement on intensity matters because intensity determines severity and therefore the whole post-fire cascade of Section~\ref{sec:water}. Experimental evidence from laboratory benches, wind tunnels and field burns, none of it Alpine, suggests why such disagreement is to be expected \citep{hayajneh_wind_2026}. Operational spread models apply wind and slope as separate correction factors and so miss the non-linear interaction between them. On laboratory benches, standard formulations under-predicted spread on steep slopes with wide fronts, and such models cannot represent eruptive acceleration in confined geometries without ad hoc parameterisation. We therefore expect them to perform worst in the steep terrain where Alpine fires burn.

The same synthesis of experimental and modelling work shows where the separable treatment stops working \citep{hayajneh_wind_2026}. On slopes, preheating becomes convection-led beyond roughly 30\degree{}, as it does above the wind threshold given above. In pine-needle beds, rate of spread rises about two and a half times between 20\degree{} and 30\degree{}, and the distance over which convective heating acts expands from some 7\,cm at 20\degree{} to more than 130\,cm at 30\degree{}. Eruptive transition appears where slope exceeds 25\degree{} to 30\degree{} in a channel of aspect ratio at least 0.2. There, the convective influence distance grows from about 8\,cm at 25\degree{} to nearly 3\,m at 40\degree{}. Because empirical and analogue formulations prescribe the coupling of wind and slope rather than resolving it, they miss these transitions and degrade outside their calibration envelope, which is defined by steep slopes, extreme winds and discontinuous fuels.

Terrain also determines where fire does not go, and this side of the relationship has been quantified outside Europe. Across 912{,}000\,ha of the Canadian Rocky Mountains, 557 mapped patches of persistent fire refugia covered 27.7\% of the forested area \citep{rogeau_effect_2018}. They were predicted from terrain and fuel, with elevation the leading variable at 23.7\% relative influence. Elevation was followed by the proportion of non-fuel in the surrounding landscape at 20.3\%, solar radiation as an aspect proxy at 14.6\% and topographic position at 10.1\%. Skill depends strongly on scale. The area under the curve rises from 0.71 for refugia of all sizes to 0.84 for patches under 10\,ha, so terrain predicts small unburnt patches far better than large ones, and the authors note that topography may matter less during large fire events. This scale dependence is favourable in a range whose fire-size distribution is dominated by small events, and the same geometric argument supports terrain-based prioritisation in the Alps.

The terrain circulation that coupled models resolve is itself a driver of spread. A high-resolution simulation of a 2023 fire in the Greater Khingan Mountains reproduced the observed burnt area above 800\,m with a spatial similarity of 0.82, while overestimating it in the valleys \citep{wang_dynamic_2026}. It attributes the fastest phase of the fire to nocturnal down-valley winds of 3 to 5\,m\,s$^{-1}$. During that phase, maximum spread rate rose from 0.3 to 1.1\,m\,s$^{-1}$ and peak heat flux from 11 to 159\,kW\,m$^{-2}$, before both fell back as daytime up-valley winds returned. The setting is Chinese boreal rather than Alpine, and the coefficients do not transfer. The mechanism, however, is the diurnal valley circulation that operates in every Alpine valley. It implies that the hours of fastest Alpine spread may be the hours when crews are least able to work.

Slope also governs a mode of spread that perimeter-based evaluation does not capture. Alpine lightning fires are reported as harder to extinguish because they burn below the surface. Controlled smouldering experiments show that slope also accelerates this mode of spread \citep{shan_influence_2024}. Between level ground and a twenty-degree slope, peak temperature rose from 769 to 948\degree{}C and vertical and horizontal spread rates rose by roughly 75 to 90\%, with the slope effect on temperature exceeding its effect on rate. The experiments are Chinese boreal rather than Alpine, so the coefficients do not transfer, but the direction of the effect does. It bears on the mop-up phase, which takes most of the effort in an Alpine incident.

Severity itself appears only partly predictable from terrain. Across 113 large Italian fires, pre-fire land cover was the strongest predictor of burn severity, with conifer stands highest, followed by elevation, while slope contributed weakly \citep{malandra_burn_2022}. Distance to roads and railways entered positively, indicating higher uncontrolled severity in remote terrain. Fuel type therefore leads and access modulates, while slope, so dominant for spread, is secondary for severity.

\subsection{Fire danger and hazard modelling}

In Europe, the Canadian Fire Weather Index is preferred in the Mediterranean countries, and when it is used outside Canada its danger classes are redefined by local calibration \citep{zacharakis_environmental_2023}. The same review covered 63 danger systems and indices from 16 countries, few of which represent terrain at all beyond coarse slope classes or pressure-level adjustments. The Italian operational system is the clearest exception, and it contradicts the view that European danger rating ignores fuel and terrain. RISICO computes dead fuel moisture content, rate of spread and fireline intensity hourly on a one-kilometre grid \citep{perello_fuelaware_2025}. Its fuel term comes from a susceptibility classification that splits the territory into twelve classes, combining vegetation type with a susceptibility level in which topography is already embedded. The underlying susceptibility model was validated against 102{,}904 wildfires over sixteen years and compared with twelve established indices. It reaches an area under the curve of 0.89 at European level, and larger fires concentrate in its highest susceptibility quartile. Its authors also state its limits. Live fuel moisture is held at a fixed worst case rather than varying, and because of the one-kilometre grid the indices were evaluated only against fires larger than 70\,ha, which is above the size of almost every Alpine fire.

Because these instruments are rarely discussed together, Table~\ref{tab:models} lists those that appear in the Alpine and European evidence reviewed here, with the constraint each meets in steep terrain.

{\footnotesize
\setlength{\tabcolsep}{4pt}
\renewcommand{\arraystretch}{1.05}
\begin{xltabular}{\textwidth}{@{}>{\raggedright\arraybackslash}p{2.2cm} >{\raggedright\arraybackslash}p{2.1cm} >{\raggedright\arraybackslash}X >{\raggedright\arraybackslash}X >{\raggedright\arraybackslash}p{2.5cm}@{}}
\caption{Instruments used in Alpine and European fire hazard and post-fire assessment in the literature reviewed, with the constraint that mountain terrain places on each.}\label{tab:models}\\
\toprule
\textbf{Instrument} & \textbf{Family} & \textbf{Represents} & \textbf{Constraint in mountains} & \textbf{References} \\
\midrule
\endfirsthead
\multicolumn{5}{@{}l}{\emph{Table~\ref{tab:models}, continued}}\\
\toprule
\textbf{Instrument} & \textbf{Family} & \textbf{Represents} & \textbf{Constraint in mountains} & \textbf{References} \\
\midrule
\endhead
\bottomrule
\endlastfoot
Canadian Fire Weather Index & Meteorological danger index & Fuel moisture and spread potential from weather alone & Terrain does not enter the index; class boundaries require local recalibration, and in Austria no single index performs best in every ecoregion & \citealp{zacharakis_environmental_2023, arpaci_selecting_2013} \\
Austrian integrated danger system & Multi-layer danger assessment & Daily fire weather with human-ignition hazard, lightning hazard, fuel type and topography & Improves prediction of whether a fire starts, not of how large it becomes & \citealp{muller_towards_2020} \\
RISICO & Fuel-aware operational danger system & Hourly dead fuel moisture, rate of spread and fireline intensity on a 1\,km grid, with a topography-informed susceptibility classification & Live fuel moisture fixed at worst case; because of the 1\,km grid, evaluated only against fires larger than 70\,ha & \citealp{perello_fuelaware_2025} \\
FlamMap & Deterministic gridded spread simulator & Surface spread across a fuel and terrain landscape under constant conditions & Resolves fuel and topography finely, but cannot represent time-varying wind or suppression and needs a detailed fuel map, so suits planning rather than operations & \citealp{malanchini_modeling_2025} \\
PROPAGATOR & Stochastic cellular automaton & Probabilistic propagation with time-varying wind & Fast and operationally flexible, but diverges from deterministic solutions on intensity and rate of spread & \citealp{malanchini_modeling_2025} \\
LandClim & Dynamic landscape vegetation model with fire & Fire as an emergent property feeding back on vegetation & Applied to future climate scenarios and to Alpine vegetation history rather than to hazard assessment; the coupling has not been carried into operational hazard work & \citealp{schumacher_relative_2006, colombaroli_species_2010} \\
LPJmL with SPITFIRE and BASE & Fire-enabled dynamic global vegetation model & Continental burnt area under climate and socio-economic pathways & Resolution and process representation are continental; management capacity enters as a scenario rather than as terrain & \citealp{billing_future_2026} \\
BASE (Burnt Area Simulator for Europe) & Empirical continental burnt-area model & European burnt area from climate and human variables & Explains 32--36\% of deviance; does not reproduce fire occurrence in temperate and boreal Europe; trained on satellite burnt area that misses many small fires & \citealp{forrest_understanding_2024} \\
Post-fire debris-flow likelihood model & Empirical logistic model & Debris-flow probability from burn severity on steep gradients, soil erodibility and short-duration rainfall & Transfer tested in one Greek post-fire area of 21 basins with the original coefficients; regional refitting recommended & \citealp{diakakis_exploring_2023} \\
Initial Attack Assessment index & Composite terrain and behaviour index & Difficulty of initial attack from simulated fire behaviour and a terrain sub-index of accessibility, penetrability, line construction and slope & Calibrated on Californian fires and demonstrated on six fires elsewhere; not yet applied in the Alps & \citealp{cardil_assessing_2025} \\
Rockfor.net, RockyFor3D and root-bundle slope stability models & Protective-function models & Rockfall arrest and shallow-landslide probability as functions of stand structure and root reinforcement & Require stand inventory and root data that exist for few burnt Alpine stands & \citealp{maringer_temporal_2016, moos_climate_2021, gehring_shallow_2019} \\
dNBR, RdNBR and RBR with Composite Burn Index & Spectral severity indices & Fire-induced change to vegetation and soil from optical imagery & Seasonal timing, illumination and phenology bias retrieval; severity decouples from soil hydrological condition & \citealp{morresi_mapping_2022, malandra_burn_2022, negri_wildfire_2026} \\
\end{xltabular}
}

The size of the recalibration problem has been measured directly in an Alpine country. Twenty-two fire weather indices and two raw meteorological variables were tested against 1{,}331 Austrian fire records from 1993 to 2007 \citep{arpaci_selecting_2013}. The analysis covered 26 stations spanning 22 ecoregions from colline to alpine zones, and no single index performed best across regions or seasons. In summer, mean daily temperature and the Canadian Build-Up Index were selected most often, with the Build-Up Index dominating once the analysis was restricted to fires above 1{,}000\,m$^2$. In winter, M68dwd, a central European index of a different lineage, was selected at 14 of 25 stations for all fires and at 17 for fires above 100\,m$^2$. Index choice in this range therefore depends on season and ecoregion at the same time, which goes further than the usual observation that indices need local class boundaries. The same analysis reports that 42\% of the fire records lacked a size, which bears directly on the observational problem discussed in Section~\ref{sec:data}.

Alpine practice has responded to this dependence through local recalibration in some places and through integration in others. In Canton Ticino, a locally fitted danger index was developed with cantonal funding \citep{de_angelis_modelling_2015}. The niche-modelling approach developed there found that meteorological variables alone matched or outperformed the full Fire Weather Index, but that the resulting index was not transferable to other regions. The study also concluded that model selection must be revisited when legislation or socio-economic conditions change the pyrological environment. The same work showed that usable indices can be built from small numbers of fires, which matters in regions with sparse records.

Austria chose integration instead. Its system combines daily fire weather with human-ignition hazard, lightning hazard, a high-resolution fuel type map and topography \citep{muller_towards_2020}. In validation, some integrated variants outperformed conventional weather indices for predicting ignitions and performed better on false alarms, but correlation with burnt area remained low across all variants. Adding physically motivated layers improved prediction of \emph{whether} a fire starts without improving prediction of \emph{how large} it becomes.

Taken together, these results show where the gap lies, and it is not where it is usually assumed to be. The gap is not mainly one of variables, since a locally fitted weather-only model can perform well. It is a lack of local calibration data and of information on dynamic fuel condition. Operational danger systems also almost entirely omit terrain, even though terrain demonstrably controls the surface wind field \citep{ozaki_impact_2022, sharples_overview_2009}. On the fuel side the gap is literal in Austria, where no fuel models for predicting fire behaviour exist and behaviour modelling relies mainly on North American models. The Swiss models that do exist were developed in national parks and carry higher fine-fuel loads than the managed Austrian forests measured since \citep{neumann_fuel_2022}.

The Alpine literature already contains a way beyond static fuel representation, but it has not been carried forward. Landscape simulation for the Swiss Alps with the dynamic landscape vegetation model LandClim treated fire as an emergent property that fed back on vegetation \citep{schumacher_relative_2006, colombaroli_species_2010}. Palaeoecological reconstruction has also been combined with dynamic simulation to constrain species response to fire and temperature \citep{colombaroli_species_2010}. More recently, forest-dynamics models linked to protection indices \citep{hillebrand_contrasting_2023} or coupled to rockfall simulation \citep{moos_climate_2021} have quantified changes, often losses, in protective function under climate scenarios. On the evidence assembled here, reconnecting this tradition to fire hazard assessment is the most valuable methodological opportunity in Alpine fire science.

\subsection{Data, records and Earth observation}
\label{sec:data}

The records on which Alpine fire science depends differ in resolution by orders of magnitude, and we describe them from the coarsest to the finest. Continental burnt-area products derived from coarse optical sensors map burnt pixels at 250\,m and are aggregated to grids of roughly 9\,km for European fire modelling \citep{forrest_understanding_2024}. The European Forest Fire Information System maps individual perimeters and is the source of harmonised European statistics, but its mapping is effective only from about 30\,ha upward \citep{beloiu_climatedriven_2026}. Medium-resolution optical archives support stand-scale disturbance mapping and have been used to classify fire, wind and bark beetle events across European forests since 1985 \citep{miguel_unravelling_2026}, as discussed in Section~\ref{sec:forest}. National fire databases record events far below any satellite threshold, including an Austrian national register holding several thousand historical and current forest fires \citep{papathomakohle_towards_2025}.

The mismatch between these records is the main methodological problem in Alpine fire science. Mapped fires account for approximately 95\% of the burnt area that European countries report, but for only a fraction of the number of fires \citep{jrc_advance_2026}. A gap of the same order has been quantified in another European mountain setting. A danger model for Fennoscandia was built from 444{,}030 monthly grid observations containing 1{,}439 satellite-detected fire occurrences, that is 0.3\% of the record \citep{bakke_datadriven_2023}. For 2017 and 2018, Norwegian national fire records held about 800 fire data points where the satellite product held 24, and a model trained on the national data performed considerably better. Satellite products systematically underestimate burnt area wherever small fires dominate, as they do in mountains. In an Alpine country the result is a stark asymmetry: four fires were mapped in Austria in 2025, against a national record running to the low hundreds annually \citep{jrc_advance_2026, papathomakohle_towards_2025}. Statements about the European fire regime derived from satellite mapping are therefore statements about area. Statements about the Alpine regime from the same source describe only its largest events.

Earth observation also supplies the severity products on which post-fire assessment depends, and here the Alpine case differs most from standard practice. Remotely sensed severity maps are commonly used to locate where the effects of fire on soil are likely to be greatest \citep{negri_wildfire_2026}. In the Italian analysis described above, severity estimated in this way was related successfully to land cover, elevation, population density and distance to roads and railways \citep{malandra_burn_2022}. But the proxy fails where it matters most for mountains. Three years after a single Alpine mixed-severity fire, the burnt soils most water-repellent in laboratory tests lay mostly in low-severity plots, while in the field the unburnt soils were the most repellent \citep{negri_wildfire_2026}. Severity mapped from satellite and field data was thus decoupled from the soil condition it is meant to represent. As a result, the geomorphic cascade of Section~\ref{sec:water} cannot currently be initialised from the product that operational practice would normally use.

Terrain narrows the Alpine observational window further, and the bias that results from imperfect seasonal timing has been measured in the Alps \citep{morresi_mapping_2022}. Steep, deeply incised relief produces illumination contrasts and cast shadow that vary with season and aspect, and seasonal snow and persistent cloud occupy part of the usable window. A difference computed between two dates therefore confounds the change caused by fire with the change caused by solar geometry and plant phenology. The study mapped ten fires in the mountains of Piedmont following the drought of autumn 2017. Severity maps built from phenologically matched reflectance composites, rather than from a single pair of cloud-free images, raised overall classification accuracy by 4.4 to 5.2 percentage points, reaching 83.7\% at best. The gain fell to 2.8 to 3.2 points when the paired-image indices were calibrated. They also increased the extent classified as high severity by roughly 19\%, because suboptimal seasonal timing reduces the contrast between healthy and burnt surfaces and so understates what burned. An Alpine severity map produced without controlling for seasonal timing therefore understates high severity. The error is large enough to matter wherever severity class is used as an input rather than as a description.

National registers also hold variables that no satellite product contains, including the number of action forces, helicopters and brigades committed to each event \citep{muller_forest_2020}. They are therefore the only source from which Alpine suppression effectiveness could be estimated, but they are not currently used for that purpose. The same report notes that detection technology faces its own version of the terrain problem. Automated camera systems that measurably reduce burnt area where they operate were designed for towers overlooking flat ground, and their transfer to complex relief is unresolved. Drones are limited both by sensitivity to wind and by dense canopy.

Validation practice in Alpine and adjacent fire modelling has improved and is now reasonably rigorous by the standards of the field. National susceptibility mapping has combined spatial cross-validation with an out-of-sample operational test against a severe subsequent season \citep{trucchia_defining_2022}. Alpine ignition modelling has used multi-perspective validation that explicitly accounts for spatial and temporal structure \citep{moreno_spacetime_2025}. Negative results have been reported alongside positive ones, as in the Austrian finding that integration improves ignition prediction but not burnt-area correlation \citep{muller_towards_2020}. The remaining problem is statistical power rather than rigour. In a region with few large fires, the events that dominate risk are the events on which models can least be tested.

Seasonal stratification is one of the few design choices that demonstrably helps, because it avoids averaging across the two seasonal modes described above. National susceptibility mapping for Italy accordingly produced separate winter and summer surfaces \citep{trucchia_defining_2022}. In south-west China, analysis of mountain wildfire found that different drivers dominate in fire and non-fire seasons, and that including the seasonal dimension improved understanding of lagged effects between drivers and occurrence \citep{wang_seasonal_2023}. The same logic applies more strongly to the Alps. In an inner-Alpine valley, the seasonal contrast is not only between active and quiet periods but between climate-driven summer fires and anthropogenic winter fires \citep{vacchiano_modeling_2018}, and different mechanisms also appear to drive the two fire seasons in the Czech Republic \citep{bercak_combination_2024}.

The most fundamental constraint is the observational threshold itself. Its magnitude has been measured directly, although not in the Alps. For the Netherlands, the European reference dataset captured 0.2\% of recorded wildfires and 45\% of burnt area \citep{stoof_temperate_2024}. There is no reason to expect better performance for an Alpine fire population that is smaller still and observed through terrain that obstructs detection. The Alpine fire-size distribution described above sits almost entirely below the thresholds at which the defining analyses operate. These are the burnt-area mapping threshold given at the start of this subsection \citep{beloiu_climatedriven_2026} and the 1\,ha minimum fire size adopted by a recent five-region European ignition model \citep{gelabert_assessing_2025}, while continental burnt-area products are coarser still. The consequences run in opposite directions and are easily confused. Detected trends in large Alpine fires are trends in the tail of the distribution and may not represent the regime. Conversely, an absence of detected change may simply reflect that most of the regime is invisible to the detector. In either case, Alpine regime detection requires purpose-built records at sub-hectare resolution rather than the reuse of continental products.

Across the modelling literature reviewed above, the models in current use are narrower than the process they represent, and consistently so in one direction. Weather-only danger indices omit vegetation, topography, lightning and human influence \citep{muller_towards_2020}. They are also tuned to the region in which they were developed and do not transfer directly \citep{de_angelis_modelling_2015}, while a pooled ignition model lost skill when evaluated at individual sites \citep{gelabert_assessing_2025}. Static fuel maps cannot represent a bimodal seasonal fuel bed, nor one that changes over decades as Alpine forest cover and wood volume have \citep{bebi_changes_2017} and as the treeline has in one Central Alpine valley \citep{frei_evidence_2023}. The Alpine-scale scenario model does not consider extreme events and the disturbance-driven fuel loading they favour, nor land-use change driven by social, economic and governance shifts \citep{gerberding_mapping_2025}. Continental burnt-area models are trained on satellite burnt-area products that miss many small fires \citep{forrest_understanding_2024}, and single-hazard susceptibility mapping cannot, by construction, represent the cascade \citep{trucchia_defining_2022}. Each omission removes a process that adds to hazard rather than reducing it, so the resulting estimates should be read as lower bounds.

The obvious remedy would be to widen the models, but the Alpine record cannot support this. The harmonised regional record covers only the ten years from 2000 to 2009 \citep{conedera_characterizing_2018}, and most Alpine fires fall below the observational thresholds discussed above. The analysis that establishes upslope migration excludes fires below the mapping threshold, while conceding that most historical Alpine fires are smaller than that \citep{beloiu_climatedriven_2026}. Post-fire geohydrological observation amounts to a first inventory with no long-term monitoring \citep{melzner_wildfire-induced_2025}. Documented post-fire debris flows in the western Italian Alps number two \citep{tiranti_wildfires_2021}, and systematic Alpine spread-model evaluation rests on three events \citep{malanchini_modeling_2025}. The attribution study of the record 2025 European fire season excludes the Alps and central Europe altogether \citep{keeping_attribution_2026}. So does the annual global assessment for the same season, whose three focal regions are Canadian, Iberian and Chilean \citep{burton_state_2026}. Adding parameters to fit a record this thin would produce models that describe the observations they were fitted to and little else. This constraint, more than any limitation of method, defines what an Alpine modelling programme can realistically attempt.

The Alpine literature nonetheless shows a way forward, which is to widen scope through better-chosen structure rather than through more free parameters. A single upslope path-length metric, computed from a digital elevation model and a forest map, accounts for 56 to 65\% of the variation in Swiss large-fire statistics \citep{conedera_linking_2024}. A parsimonious additive model with a representative sampling design and threshold-based output achieves high skill on Alpine ignition \citep{moreno_spacetime_2025}. In Ticino, a locally fitted, low-dimensional index matched the transferred Fire Weather Index for two fire regimes and outperformed it for summer lightning fires \citep{de_angelis_modelling_2015}. Even for post-fire debris flow, a preliminary study of two events has proposed modifying catchment runoff parameters inside an existing warning system, rather than building a coupled model \citep{tiranti_wildfires_2021}.

On this reading, transferability appears selective rather than uniformly poor, and we consider the distinction a practical one. Empirically tuned weather indices do not transfer, but physically grounded process models often do. The empirical post-fire debris-flow likelihood model used operationally in the western United States is driven by the proportion of moderate and high soil burn severity on gradients steeper than 23\degree{}, mean differenced normalised burn ratio and soil erodibility, each scaled by peak fifteen-minute rainfall, and it has been tested outside its region of origin \citep{diakakis_exploring_2023}. Running on European soil datasets, it reproduced observed debris flows in a Mediterranean post-fire setting with an area under the ROC curve of 0.913. Similarly, the Initial Attack Assessment index combines simulated fire behaviour with a terrain sub-index of accessibility, penetrability, line-construction difficulty and slope \citep{cardil_assessing_2025}. It was calibrated on nearly 27{,}000 Californian fires and has so far been demonstrated on only six fires elsewhere, one of them in Spain \citep{cardil_assessing_2025}. The operational rule that the Alpine evidence supports is therefore to import the mechanism and recalibrate the coefficients.

\subsection{What differs in mountains}

Here, and in the closing subsections of the sections that follow, we compare the mountain case with the lowland and Mediterranean baseline rather than describing the Alps alone. Where evidence from other ranges agrees with the Alpine evidence, we treat the difference as a property of mountains. Where it disagrees, we identify what is specific to the Alps.

Outside mountains, hazard assessment usually ends with burnt area. Continental models take it as their output \citep{forrest_understanding_2024}, and in the general lowland case it also serves as the measure of loss. Ignition is predominantly human; in the Netherlands, for example, informal cause assessments attribute only 0.3\% of recorded fires to lightning \citep{stoof_temperate_2024}. Among human factors, distance to the wildland-urban interface and to roads predicts it best \citep{gelabert_assessing_2025}. Spread models treat slope as a separable correction to the effect of wind, which is defensible on gentle slopes \citep{hayajneh_wind_2026}. Danger is rated with a weather index, often the Fire Weather Index with locally calibrated classes \citep{zacharakis_environmental_2023}. In the Alps the differences begin with ignition.

Alpine ignition is mixed rather than almost entirely human. Lightning accounts for a minority of fires overall but for a large share of summer ignitions (Section~\ref{sec:hazard}) \citep{valese_wildfires_2011}. Lightning frequency also ranks among the strongest predictors in Alpine-scale hazard mapping \citep{gerberding_mapping_2025}, and lightning fires start at higher elevation \citep{conedera_lightninginduced_2006, valese_wildfires_2011} and on steeper relief \citep{conedera_lightninginduced_2006} than human ones. The regime is bimodal, with a spring season of surface fires at low elevation and a summer season, and the two respond to different drivers \citep{valese_wildfires_2011, vacchiano_modeling_2018}. A model fitted to annual aggregates is therefore fitted to two processes at once. The spring peak also appears in the Czech Republic and in the Carpathians \citep{bercak_combination_2024, manczinger_fire_2025}, so it is not specific to the Alps. In the summer-dominated Mediterranean regime, a comparable March peak has been reported for Portugal, but it contributes little to annual burnt area \citep{bercak_combination_2024}.

The differences grow once a fire is burning. Experiments on laboratory benches, in wind tunnels and in non-Alpine field burns show that on steep slopes wind and slope interact non-linearly rather than adding up, and on a laboratory bench standard operational formulations under-predicted spread on steep slopes with wide fronts \citep{hayajneh_wind_2026}. We infer that models built on separable correction factors are likely to underestimate spread in the geometry where Alpine fires actually run. The only systematic Alpine comparison of spread models shows the practical consequence, since two simulators reproduced the observed propagation but disagreed on fireline intensity and rate of spread \citep{malanchini_modeling_2025}. The wind field also differs from the interpolated synoptic field. Relief reorganises the surface atmosphere at scales below the resolution of the products that danger systems ingest, a mechanism established mainly outside Europe \citep{sharples_overview_2009, ozaki_impact_2022}. Nocturnal down-valley winds can drive the fastest phase of spread, as documented in a Chinese boreal range rather than in the Alps \citep{wang_dynamic_2026}.

Danger rating inherits these problems and adds a calibration problem of its own. Indices are tuned to the region in which they were developed and do not transfer directly \citep{de_angelis_modelling_2015}, and pooled ignition models lose skill at individual sites \citep{gelabert_assessing_2025}. In Austria, no single index performs best across ecoregions or seasons \citep{arpaci_selecting_2013}. Locally fitted indices work and can be built from few fires \citep{de_angelis_modelling_2015}, but the records available to fit them are short \citep{conedera_characterizing_2018} and, in Austria, often lack fire size \citep{arpaci_selecting_2013}. The fuel-aware Italian system has been evaluated only against fires far larger than the typical Alpine one \citep{perello_fuelaware_2025}. Locally fitted models therefore cannot be tested on the large, rare events that dominate risk.

Not every difference holds across ranges. Elevation is not a universal mountain fire predictor. Of the large fires recorded above 800\,m in the Alps, 42\% burned above 1{,}400\,m, although the regime type to which the Alps belong has so far shown only limited upslope expansion \citep{beloiu_climatedriven_2026}. In the Carpathians, by contrast, fire concentrates at low elevation, where people reach the land most easily \citep{manczinger_fire_2025}. What generalises is that terrain organises where fires occur, not the direction in which it does so.

Underlying all of this is a measurement problem: detection thresholds sit above the size of most Alpine fires \citep{muller_forest_2020, beloiu_climatedriven_2026}. Satellite mapping captures most of the burnt area but only a fraction of fires \citep{jrc_advance_2026}. Temperate lowlands dominated by small fires share this problem \citep{stoof_temperate_2024}. In mountains, however, steep relief adds cast shadow and illumination contrasts, and snow and cloud shorten the window for well-timed images \citep{morresi_mapping_2022}. Across ten fires in the western Italian Alps, severity maps that did not control for seasonal timing understated the extent of high severity \citep{morresi_mapping_2022}. Trends detected from satellite records therefore describe large Alpine fires only, and the lack of a detected trend says little about the regime as a whole.

% =======================================================
\section{Climate, Weather and Projections}
\label{sec:climate}

\subsection{Climatic controls on Alpine fire activity}

Alpine fire climate varies with elevation, between the moist northern and the insubric southern slopes, and between the dry continental inner valleys and the wetter outer ranges. Elevation sets the main constraint at present, through snow cover and temperature-limited fuel drying, and warming is relaxing exactly this constraint. Warming appears to be elevation-dependent in some mountain ranges, with the snow- and ice-albedo feedback a leading candidate mechanism shared with Arctic amplification, although whether such gradients are general remains uncertain \citep{rangwala_climate_2012}. In one Italian Alpine district, warming and drying have already cut ski provision in the lower valleys, and shorter snow cover is projected to put high-altitude municipalities most at risk by 2021--2050 \citep{pede_mountain_2022}.

Fuel moisture, rather than temperature, appears to be the direct link between climate and fire. Across European climate regions, multi-month drought indices, soil moisture deficit and vapour pressure deficit are the strongest and most consistent drivers of fire activity. Their interactions are region- and season-specific. Productivity deficits matter in southern Europe in summer, whereas in northern Europe summer preconditions combine moisture deficit with productivity \emph{surpluses}, so fuel production and fuel drying compound rather than trade off \citep{miller_compounding_2025}. Fuel moisture also shapes ignition. In a five-site European model of human-caused ignition, the annual anomaly in dead fine fuel moisture content ranked first in the pooled model and at the Swedish and central European sites, but lower at the three Mediterranean sites \citep{gelabert_assessing_2025}. Vapour pressure deficit is likewise the strongest predictor of burnt area in European mountains \citep{beloiu_climatedriven_2026}.

Drought is not a single quantity, and the Alpine spring and summer fire seasons plausibly respond to different forms of it. Analysis of extreme fires in south-eastern France separates press drought from hot drought \citep{ruffault_extreme_2018}. In a press drought, a long-lasting water deficit from reduced precipitation and high evapotranspiration intensifies wind-driven fires. In a hot drought, a comparable deficit combines with a summer heatwave to intensify heat-induced fires. Alpine spring fires, carried by cured herbaceous fuel under foehn, match press drought in mechanism. The emerging high-elevation fires, which across eight European mountain regions occur under extreme atmospheric aridity above 1400\,m \citep{beloiu_climatedriven_2026}, plausibly match hot drought. A danger system calibrated on one will not represent the other, and a projection that reports a single aggregate trend hides which of the two is changing.

Alpine-specific trend evidence comes from 25 long-term stations in six Alpine countries, grouped into climatically homogeneous regions. Mean annual fire danger, taken as the median of seven indices, rose clearly and significantly south and west of the Alps. North of the Alps and in the inner valleys, a high proportion of trends were not statistically significant. At the 95th percentile, however, almost every station showed a trend towards higher index values \citep{wastl_recent_2012}. Warming is therefore raising the extreme tail of Alpine fire danger even where the mean has not changed detectably, a pattern that monitoring designed around averages is likely to miss.

Rising danger is not the same as a changing fire regime, and the Alps have much evidence for the former and little for the latter. European fire danger showed no clear trend between 1960 and 1999 but increased significantly in southern and eastern Europe from 1980 onward. At the national level in Greece, Spain and Finland, correlations between the mean seasonal Fire Weather Index and burnt area were of order 0.6 \citep{venalainen_temporal_2014}. Regime change is harder to detect in the Alps. Records are short and administratively fragmented, fires are small enough that recording thresholds dominate the statistics, and the strong seasonal and topographic structure lets opposing changes cancel out in regional aggregates. The lack of a detected signal is therefore weak evidence that nothing has changed.

A long-term record from the driest Alpine valleys supports this caution. The fire history of Valais, a dry continental valley of the Swiss Alps, has been reconstructed as the joint product of climatic, vegetation and anthropogenic factors, with regime shifts across the twentieth century in both Valais and Ticino \citep{zumbrunnen_reconstructing_2010}. Relating some 900 of those fires to human and climatic proxies shows the controls changing with the socio-economic regime. Between 1904 and 1955, population density was the strongest predictor, followed by wood harvest, and both forest pasturing and wood harvest were associated with fewer fires. Between 1956 and 2006, after the post-war turn from an agricultural to an industrial and service economy, temperature became the leading predictor. The authors attribute this to fire activity concentrating at lower elevations as farming and forestry withdrew from middle and high ground \citep{zumbrunnen_human_2012}. Both models explain only about a third of the variance, so the change is better read as a shift of emphasis than as one driver replacing another. Alpine fire regimes have thus shifted before under human forcing, and even the Valais temperature signal emerged through land use. This cautions against attributing any present-day shift to climate by default.

\subsection{Mountain meteorology: wind systems, foehn and the thermal belt}

Synoptic control on Alpine fire danger has been characterised directly by combining a sixty-year weather-type classification with station data, danger indices and observed fires across the Alps. Elevated danger occurred mainly under persistent anticyclonic situations or in the lee of the range under katabatic foehn, and the southern Alpine region was the most fire-prone. Differences between weather types narrowed for indices representing deeper soil layers, reflecting their longer drying response times \citep{wastl_largescale_2013}. Foehn is therefore one of the two synoptic configurations that generate Alpine fire danger, not an incidental local phenomenon.

Below the synoptic scale, Alpine wind systems are, in our reading, what turns small ignitions into fast fires. Foehn can remove tens of percentage points of relative humidity within hours and is associated with the late-winter and spring fires of inner and southern valleys. Diurnal thermal circulations reverse between anabatic flow by day and katabatic flow by night, imposing a daily cycle in spread direction with no lowland analogue. Relief also channels the synoptic flow along valley axes and through cols. Steep Alpine topography controls each of these regimes, and through them fire type, intensity and burn severity \citep{melzner_wildfire-induced_2025}.

Relief also breaks up the fire-weather field. Where upper winds meet complex terrain, surface conditions become discontinuous and locally amplified, so a field interpolated from valley-floor stations loses validity \citep{sharples_overview_2009}. In one Tasmanian fire, an imported case divided into three burnt sectors, atmospheric profiles capable of generating enhanced surface winds occurred in 6.5\% of time steps over the least rugged sector but in 40.2\% over the most rugged. Fire simulations driven by surface winds alone performed well on the gentlest terrain but were outperformed on rugged terrain by simulations that included upper-air winds \citep{ozaki_impact_2022}.

Temperature inversions add a less obvious structure, and the thermal belt is a recognised feature of the Alpine fire calendar. In foehn regions north and south of the range, winter fires occur because inversion places a sunny, dry belt above the fog line \citep{muller_forest_2020}. During the autumn 2011 drought in the Bavarian Alps, fire danger and fire incidents concentrated at mid-slope elevations in the thermal belt, while valley stations registered only moderate danger. Under inversion, the diurnal cycles of the controlling meteorological variables differ between valley floor and slope, and once-daily station observation cannot represent that difference \citep{schunk_forest_2013}. The same study recommends additional high-elevation stations, or high-resolution meteorological models feeding indices that accept at least hourly input. This remedy has not been generally implemented in the Alpine danger-rating systems reviewed in Section~\ref{sec:hazard}.

\subsection{Lightning, convection and the summer ignition mode}

Lightning frequency is one of the three strongest predictors of fire hazard in Alpine-scale statistical mapping, alongside mean annual temperature and precipitation seasonality \citep{gerberding_mapping_2025}. No Mediterranean hazard model would rank an ignition source so highly.

Multi-decadal records describe the lightning regime across Valais, the Grisons, Valle d'Aosta and Ticino. Lightning fires there occur from May to October, with 78\% of events in July and August. They ignite at higher elevation and on steeper relief than anthropogenic fires, and are harder to extinguish because they burn underground. In hot, dry summers such as 2003 they became markedly more frequent, with significantly greater burnt area and higher suppression costs \citep{conedera_lightninginduced_2006}. Lightning ignition is itself structured by terrain, through elevation-dependent cloud-to-ground discharge frequency, altitudinal fuel belts and strong vertical gradients in precipitation, as described for mountain forests in southern Siberia \citep{baranovskiy_deterministic_2021}. An Austrian integrated danger-assessment prototype reflects this by classifying lightning fire hazard jointly by elevation, slope, aspect and vegetation type, derived from the locations of past lightning fires \citep{muller_towards_2020}. Whether ignition potential and fuel receptiveness will rise together is the compounding question. In Croatia, which the Dinaric Alps cross, lightning ignited only 3.4\% of fires over six years, yet these burned on average twice the area of all fires, and some burned beneath transmission lines. Summer fire weather there is projected to reach very high or extreme danger classes by the end of the century, with more thunderstorms considered possible \citep{bozicek_impact_2023}.

\subsection{Climate projections and what they omit}

The Alpine scenario literature spans statistical hazard mapping and fire-weather projection, and the limitations the two share say more than their individual results. Hazard mapping under shared socio-economic pathways yields the increases reported in Section~\ref{sec:intro}, with the largest relative changes in the northern and eastern sectors \citep{gerberding_mapping_2025}. Fire-weather projection for the heterogeneous landscapes of adjacent central Europe, using a regional large ensemble under RCP8.5, places the emergence of the forced fire-danger signal before 2050 in every subregion examined. It also shifts present-day hundred-year fire weather to a ten-year recurrence by 2090 \citep{miller_climate_2024}.

These shared limitations are, in our view, structural rather than incidental. Neither approach represents the disturbance-driven fuel loading by which drought, storm and bark beetle outbreak increase flammable material, or the socio-economic and governance change that affects land use. Both state openly that local application carries far more uncertainty than the regional trend \citep{gerberding_mapping_2025}. Continental fire models show a related limit: the Burnt Area Simulator for Europe explains 32--36\% of deviance overall and does not reproduce fire occurrence in temperate and boreal Europe \citep{forrest_understanding_2024}. It does not evaluate its skill in mountains separately, which is itself a gap. A projection driven by climate normals describes the climatic envelope, not the fire regime.

Anticipation faces a further limit, although the evidence for it is indirect. A global analysis found that thirteen climate teleconnections together modulate around half of global burnt area. These relationships are markedly weaker in Europe than on other continents. The European associations are also spread across several modes, including the North Atlantic Oscillation and the East Atlantic pattern in temperate Europe and Pacific and Indian Ocean modes in Mediterranean and temperate regions \citep{cardil_climate_2023}. Seasonal fire forecasting from remote climate modes will therefore deliver less in the Alps than where a single mode dominates, which puts more weight on monitoring the land surface and the state of the fuel.

\subsection{What differs in mountains}

Outside mountains, fire climatology can treat fire weather as a spatially smooth field. Station observations are interpolated or assimilated into gridded fields, drought is expressed by multi-month indices computed on those fields \citep{miller_compounding_2025}, the fire season is a single window bounded by summer heat, and projection moves the same envelope forward. In the Alps each of these assumptions weakens.

Relief breaks up the weather field. Upper winds interacting with complex terrain produce discontinuous, locally amplified surface conditions \citep{sharples_overview_2009}, and profiles that generate enhanced surface winds are several times more frequent over rugged than over gentle terrain \citep{ozaki_impact_2022}. Foehn raises fire danger and fire occurrence in the lee of the Alps \citep{wastl_largescale_2013}, and foehn, mountain and valley winds, all shaped by steep topography, influence fire type, intensity and burn severity \citep{melzner_wildfire-induced_2025}. Under inversion, danger can peak in a mid-slope thermal belt while valley stations register only moderate values \citep{schunk_forest_2013}. A danger index computed from valley observations can therefore misrepresent conditions on the slope where the fire burns.

The range also holds more than one fire climatology. Danger trends differ between the south and west of the Alps and the north and inner valleys \citep{wastl_recent_2012}, and fires respond to different forms of drought. In Mediterranean France, wind-driven summer fires followed a long press drought and heat-driven fires a hot drought \citep{ruffault_extreme_2018}, a distinction not yet tested in the Alps, and across European mountains fires above 1,400\,m burn under drier atmospheric conditions than those below \citep{beloiu_climatedriven_2026}. Fire climates can thus differ within tens of kilometres. Continental analyses average across this variability, and the continental burnt-area model for Europe does not report its skill for mountain terrain separately \citep{forrest_understanding_2024}.

Snow gives the mountain fire season a second boundary that lowlands lack. At elevation, snow cover as well as moisture sets the length of the drying window. Snow-cover duration is projected to shorten in the upper valleys of one Italian Alpine district \citep{pede_mountain_2022}, and on the Tibetan Plateau, outside Europe, the snow season has shortened fastest at high elevation \citep{rangwala_climate_2012}. In one Swiss Alpine valley, most new trees over four decades infilled below the former treeline \citep{frei_evidence_2023}, while fire activity across European mountains has moved upslope \citep{beloiu_climatedriven_2026}. That snow loss and fuel infilling coincide in the same elevation belts is our inference, not a measured overlap. The Alpine fire season can therefore lengthen through a cryospheric control with no counterpart in Mediterranean or lowland fire calendars.

Terrain also structures ignition. In the lowland Netherlands, informal cause assessments attribute almost all fires to people and only 0.3\% to lightning \citep{stoof_temperate_2024}. In the Alps, lightning is a leading predictor of fire hazard \citep{gerberding_mapping_2025}, and lightning ignitions occur at higher elevation and on steeper ground than human ignitions (Section~\ref{sec:climate}) \citep{conedera_lightninginduced_2006}. An Austrian danger-assessment prototype accordingly classifies lightning hazard by elevation, slope, aspect and vegetation type, derived from past lightning-fire locations \citep{muller_towards_2020}.

Projection inherits these differences. Scenario studies for the Alps and adjacent central Europe agree that fire danger will rise, but they resolve the climatic envelope at regional scale and omit disturbance-driven fuel loading and land-use change \citep{gerberding_mapping_2025, miller_climate_2024}. Because terrain modifies that envelope locally, a regional projection says much less about a given slope or valley than about the range as a whole.

% =======================================================
\section{Water, Drought and Water Use}
\label{sec:water}

\subsection{Drought, fuel moisture and the physiological chain}

Drought reaches fire through a chain of separable links. The fire literature has mostly studied the link closest to fire, fuel moisture, while the forest health literature has studied the rest. Fuel moisture is well established as a driver but poorly modelled. A global evaluation of 1{,}603 measurements from 153 experiments in 21 countries found vapour-pressure-deficit models performing competitively over the full moisture range, with mean absolute errors near 4\% and coefficients of determination around 0.4. Machine-learning alternatives improved on them only marginally. Below 10\% moisture, however, no approach exceeded a coefficient of determination of 0.16, and this is the range in which spread responds exponentially to dryness \citep{rodrigues_vpdbased_2024}. Even so, fuel moisture and its atmospheric drivers remain the strongest predictors of ignition and burnt area (Section~\ref{sec:climate}).

Further back in the chain, drought acts physiologically and turns living trees into standing dead fuel. In Scots pine at the dry margins of its range, which include the inner Alpine valleys, carbon starvation may be the main route to drought-induced mortality, although evidence that reserves are actually exhausted is contested. Reduced assimilation under water stress would leave fewer resources for defensive compounds, so that pests and pathogens could amplify the effect of drought \citep{giuggiola_droughtinduced_2010}. A sapling experiment outside the Alps isolates the interaction. In Norway spruce, drought and root-rot infection compete for the same carbon pool, which must fund both compartmentalising the pathogen and maintaining turgor, while the infection reduces functional sapwood and so adds to water stress. Needle starch concentration predicted which saplings died \citep{gomezgallego_interaction_2022}. Scots pine showed no such interaction, and because the two species occupy different Alpine positions, they would be expected to fail by different routes.

Whether such decline proceeds gradually or abruptly is not settled, although the comparative evidence favours abruptness. In central Chile, more than 90\% of a Mediterranean-type forest browned within fewer than a hundred days, under conditions unprecedented across all five Mediterranean-type ecosystems. Long-term recovery is uncertain given a continuing decline in regional water balance \citep{miranda_widespread_2023}. The setting is not Alpine and the magnitude does not transfer, but the mode does: drought-driven decline can be synchronous and regional rather than incremental, which would turn a gradual fuel-loading trend into a step change.

The chain is completed when dead and weakened trees become fuel. Alpine hazard projections acknowledge that drought, heatwaves and storms favour bark beetle outbreak and windthrow and so increase flammable material, but they leave this pathway out \citep{gerberding_mapping_2025}. Since disturbance is a normal part of Alpine forest dynamics rather than an exception \citep{bebi_changes_2017}, and European forest disturbance has become more frequent over three decades while its severity has fallen \citep{senf_mapping_2021}, the omission is unlikely to be small.

Because the same stress drives both mortality and fuel loading, intervention against drought stress and intervention against fuel are largely the same operation in inner Alpine valleys. Section~\ref{sec:forest} develops this alignment, which management planning does not currently exploit.

\subsection{Soil alteration and the limits of severity as a proxy}

The post-fire hazard chain assumes that a severity map can stand in for soil condition. Fire alters mountain soils by removing vegetation and litter, depositing ash, creating or destroying water repellency and physically weathering rock \citep{santi_wildfire_2022}. For the Alps, the questions are how long these changes persist and whether the burn severity products used in risk assessment can reveal them. Present evidence suggests that they cannot.

Three years after a mixed-severity fire in an Alpine Scots pine forest, burnt plots still showed greater bare soil exposure and a pyrogenic surface horizon with elevated pH and pyrogenic carbon. Their sub-surface horizons were depleted in organic carbon and in sand-sized aggregates. These differences involved properties that normally recover quickly as well as properties that do not. The magnitude of change was not proportional to mapped severity class. In laboratory tests, the strongest water repellency occurred mainly in low-severity plots, associated with abundant uncharred organic matter and probably intermediate heating rather than intense heating, although field tests found unburnt soils the most repellent \citep{negri_wildfire_2026}. Remote severity mapping, even when calibrated with field measurement, is therefore a poor predictor of the soil condition that governs runoff generation. This challenges the standard post-fire workflow, in which a severity map serves as the spatial proxy for hydrological and geomorphic change. It also undermines the common fuel-management assumption that a low-intensity burn is hydrologically benign.

That so basic an inference has gone largely untested reflects a wider neglect of mountain soils. A review of the soils of the Australian Alps records that above-ground ecosystems received far more attention than the soils beneath them, although high-elevation humus and peat soils matter particularly for ecology, hydrology and carbon storage, and peat soils are particularly vulnerable \citep{wilson_distribution_2022}. The European Alpine literature shows the same imbalance, which matters because the Alpine cascade runs through soil condition and not only through canopy loss.

Alpine natural-hazard consequences of fire were catalogued before the recent inventories, and the list goes well beyond debris flow: new avalanche-prone slopes, mudslides, increased runoff, rockfall and soil erosion. Steep subalpine slopes with a thin humus layer over limestone were identified as the setting where complete soil loss and failed regeneration are most likely \citep{muller_forest_2020}. The following decade added quantification rather than recognition.

\subsection{Erosion, debris flow and sub-threshold triggering}

Post-fire debris flow carries the highest damage potential in Alpine valleys, and the small Alpine record is now quantitatively informative. Two documented events in the Western Italian Alps, at Pallanzeno in 2005 and Bussoleno in 2018, both followed brief, localised, moderate-intensity rainfall over catchments burnt in the preceding months. Both produced deposits unusually rich in fine particles and burnt wood, consistent with a centimetres-thick layer of thermally reworked material produced by the fire and progressively redistributed by intervening rainfall \citep{tiranti_wildfires_2021}.

Numerical modelling of the Bussoleno event makes the mechanism quantitative, although the sample remains a single catchment. An October 2017 fire burnt over 95\% of a 1.37\,km$^2$ catchment. After a spring 40\% wetter than normal, a radar-estimated rainfall of 12\,mm in 30 minutes (an intensity of 24\,mm\,h$^{-1}$, a five-year return period) mobilised roughly 20{,}000\,m$^3$ of sediment and woody debris some eight months after the fire. The flow destroyed two buildings and forced the evacuation of around 150 residents of the alluvial fan. In situ tests found the shear strength of burnt soils reduced by a factor of about three. The triggering intensity lay below the threshold of more than 30\,mm\,h$^{-1}$ typical of unburnt basins of that type, although complex orography limits the reliability of the radar estimate \citep{laporta_postwildfire_2025}. Fire therefore lowers the rainfall needed to trigger a flow, rather than only increasing the response to a given rainfall, and so moves the hazard into a return period that occurs several times per decade.

The Alpine value is not anomalous when compared with the range where this science is most developed. In Glenwood Canyon, Colorado, after a fire that burnt 13{,}000\,ha, the fire-wide triggering threshold in the first post-fire year was 25.9\,mm\,h$^{-1}$, close to the Bussoleno value. Of the observed debris flows, 89\% were triggered above the modelled threshold. In the second year the operational threshold was 33.7\,mm\,h$^{-1}$, and eight storms exceeded it without producing a single debris flow. The authors attribute this to vegetation recovery, from 24\% to 57\%, and to sediment depletion, and recalculating the threshold from imagery of the recovering vegetation raised it to 40\,mm\,h$^{-1}$ \citep{rengers_evaluating_2024}. The threshold is therefore a recovering quantity rather than a constant, so for Alpine application the length of the low-threshold window matters more than its depth. The same study found that the standard volume model overestimated observed erosion by a median factor of 4.4. A threshold imported with confidence should therefore not be accompanied by a volume estimate imported with the same confidence.

Thermal spalling, a sediment source specific to rocky burnt terrain, has been quantified in a European mountain setting. A fire burnt 131.64\,km$^2$ of Sardinia at temperatures above 800\,\degree{}C, with 63.1\% of the area at moderate-high or high severity. Spalling at 69 rock exposures removed between 0.07\% and 74.57\% of the exposed surface, with a median of about 20\%, in fragments from 3\,mm to 13\,cm. In the following two years, three rainstorms triggered 33 debris flows across 21 sub-basins, and spall fragments made up between 1.3\% and 36.6\% of their deposits. Freshly spalled angular cobbles formed 10 to 30\% of all cobbles in the deposits. Spalling did not differ between high and moderate-high severity but was sparse at lower severity \citep{pala_sediment_2025}. Fire in rocky terrain therefore produces transportable sediment directly, rather than only exposing what was already there.

Sediment supply, rather than runoff alone, appears to set event magnitude, although this demonstration is imported rather than Alpine. Lidar differencing of the 2018 Montecito flows in California showed 550{,}000\,m$^3$ of previously stored channel sediment progressively entrained and partly redeposited. Scour concentrated where coarse colluvium had accumulated and persisted beyond the mountain front. This implies that flow volume is controlled by the sediment available for scour, and hence by catchment history \citep{morell_sediment_2021}. In the Alps, post-fire debris-flow magnitude would then depend on the antecedent state of the channel network as much as on the fire.

\subsection{Water supply, water quality, snow and ice}

The cascade continues beyond the fan apex into the water system that the range supplies. Alpine catchments provide drinking water, hydropower and irrigation to populations far larger than those living in the mountains, and the processes that carry sediment onto a road also deliver it, with its chemistry, into storage.

The size of the water-quality signal has not been measured in the Alps. For a Mediterranean supply reservoir, a modelled scenario burning around 30\% of the contributing area raised median suspended solids by a factor of 4.7 and maximum nitrate and phosphate concentrations by a factor of four. Dilution within the reservoir removed roughly half the nitrate and nine tenths of the phosphate before the dam wall. Both nutrients stayed within drinking-water limits, although suspended solids exceeded the recommended maximum value \citep{basso_potential_2021}. The mobilisation, and the dependence of the outcome on storage volume, are general enough to transfer to the Alps. The buffering is contingent, however, since the same study notes that drought could weaken dilution, and drought is when the antecedent conditions for fire are present. In a review of 56 Greek studies of fire effects on ecosystem services, every recorded entry for hydrological cycle regulation was negative, as was every entry for cultural services. Across all services, effects were most pronounced in the first two post-fire years \citep{misseyanni_impact_2025}. Water regulation also features among the services assessed as at risk from wildfire in Mediterranean forest landscapes \citep{lecinadiaz_assessing_2021}.

Snow currently limits high-elevation Alpine fire, and the coupling between fire and snow runs both ways. The effect of snow on fire is the better documented direction, although largely by inference. Mountain warming reduces winter and spring snowpack \citep{rangwala_climate_2012}, snow cover duration is projected to shorten in the upper valleys of one Piedmont zone \citep{pede_mountain_2022}, and in one Swiss valley forest has advanced and infilled near the treeline \citep{frei_evidence_2023}. A shorter snow season should therefore lengthen the drying window where fuel is now accumulating. The effect of fire on the cryosphere has no Alpine literature at all, and comparative evidence suggests that this is a research gap rather than an absent process.

Fire alters snowpacks through several mechanisms at once. Canopy loss changes interception and the radiation balance at the snow surface, burnt debris and black carbon lower albedo, and the resulting changes to accumulation and ablation persist for years while the canopy recovers \citep{koshkin_wildfire_2022}. The radiative term has been measured directly on ice. On a well-instrumented glacier in the Canadian Rockies, upwind wildfire activity affected melt both through albedo reduction from deposited soot and through the effect of smoke on the overlying atmosphere. Smoky days were warmer than clear days, with a sustained reduction in surface shortwave irradiance of 103\,W\,m$^{-2}$ at peak \citep{aubrywake_fire_2022}.

The net effect on snow has been measured in a mountain range, and it runs in an unexpected direction. Over eleven water years in the northern Sierra Nevada, snow-covered area in a single burnt basin increased significantly after the fire compared with an unburnt control, regardless of slope aspect or burn severity. Days above a pre-fire high snow-cover threshold rose by 81\%, yet basin-average melt-out came earlier than in the control \citep{micheletty_application_2014}. Canopy removal therefore puts more snow on the ground and removes it sooner, with opposite effects on the two quantities that matter for an Alpine catchment, storage and timing. Nothing would prevent either mechanism from operating in the Alps, where snow and ice carry economic weight far beyond their hydrological role \citep{pede_mountain_2022}. The lack of Alpine work leaves a term missing from the water-resource accounting of a range whose runoff supports downstream agriculture, hydropower and navigation. It also leaves unexamined a feedback in which fire shortens the snow season that currently limits fire.

The drought that drives fire also strains the Alpine water system, and the two stresses compound. During the 2022--2023 snow drought, glacier melt in two large Italian Alpine catchments rose to 148\% above its decadal mean. Its share of summer streamflow doubled to tripled, reaching 75\% and 83\% at peak against medians of 31.5\% and 22.6\%. Streamflow nonetheless stayed below the long-term average, so the glacial buffer only partly compensated for the precipitation deficit \citep{leone_20222023_2026}. On one glacier in the Canadian Rockies, soot deposited from upwind fires increased melt by up to about 10\% in later seasons, although smoke itself partly offset this by reducing net radiation \citep{aubrywake_fire_2022}. Together the two results suggest a loop. The drought that raises fire danger also exhausts the snowpack and shifts water supply onto glaciers, and soot from fire may then accelerate their melt, although this last link is imported from a single Canadian glacier. No Alpine study has attempted to quantify this loop, and we identify it as a research priority rather than an established result.

\subsection{What differs in mountains}

Outside mountains, the post-fire water problem is well characterised and comparatively benign. Vegetation loss and soil sealing raise runoff and sheet erosion. The response peaks in the first one to three years and declines as cover returns, and in a review of Greek studies, fire effects on ecosystem services in general were most pronounced in the first two post-fire years \citep{misseyanni_impact_2025}. Sediment reaches channels as suspended load, and where the catchment supplies water, reservoir storage buffers the chemical signal \citep{basso_potential_2021}. Steep slopes change each part of this picture.

On steep slopes the process changes, not only its size. Fire alters soils and weathers rock \citep{santi_wildfire_2022}, and runoff on steep ground mobilises coarse material, producing debris flow rather than sheet wash. The fire itself supplies part of that material, as a thermally reworked surface layer \citep{tiranti_wildfires_2021} and as spalled rock fragments \citep{pala_sediment_2025}. The damaging event is a mass movement that runs out onto an alluvial fan, which at Bussoleno is where the destroyed buildings stood \citep{laporta_postwildfire_2025}. Severity maps do not reliably indicate the soil condition that starts this process \citep{negri_wildfire_2026}, and Alpine documentation of post-fire geohydrological processes is still at the stage of a first inventory \citep{melzner_wildfire-induced_2025}.

Because this material is already primed, fire lowers the triggering threshold instead of only amplifying the response. In the one Alpine catchment where this has been examined, the burnt basin responded to a radar-estimated rainfall intensity below that required by unburnt basins of its type \citep{laporta_postwildfire_2025}. The first-year threshold in Colorado, a 15-minute intensity of 25.9\,mm\,h$^{-1}$, was of similar magnitude, although that study made no comparison with unburnt basins \citep{rengers_evaluating_2024}. The hazard thus moves from the tail of the rainfall distribution to intensities that recur several times per decade. At that point, the return-period logic used to design mountain protection works no longer describes the post-fire state.

How long the threshold stays low depends on biology as much as on geomorphology, because it rises as vegetation recovers \citep{rengers_evaluating_2024}. The general figure of about three to four years for sediment yield to return to pre-fire levels comes from a review that itself notes recovery times of the order of 50 to 500 years in alpine terrain \citep{santi_wildfire_2022}. In the Alps, recovery after the one dry inner-Alpine burn studied was limited by moisture deficit at montane elevations \citep{moser_potential_2010}, and at the upper margin it has failed entirely in the one massif where this has been studied \citep{degiuli_arrested_2024}. Post-fire susceptibility can therefore persist for decades \citep{gehring_shallow_2019} and, on some subalpine slopes, effectively indefinitely.

Further downstream, the coupling with water resources has no lowland counterpart. The drought that raises fire danger also exhausts the snowpack and shifts summer supply onto glacier melt \citep{leone_20222023_2026}. On one glacier in the Canadian Rockies, soot deposited by fire raised that melt by up to about 10\%, while smoke partly offset the gain by reducing net radiation \citep{aubrywake_fire_2022}. Fire also changes how much snow is stored and when it melts \citep{koshkin_wildfire_2022}, and in one Sierra Nevada basin it enlarged the snow-covered area but brought melt-out forward \citep{micheletty_application_2014}. This coupling is a property of a glaciated range rather than a scaled-up lowland process, and it has not been quantified in the Alps in either direction.

% =======================================================
\section{Forest: Change, Ecology and Management}
\label{sec:forest}

\subsection{Forest change, composition and fuel connectivity}

The Alpine forest that now burns is largely a recent product of land use, and this determines which policy levers can change it. For centuries, grazing, litter collection and fuelwood extraction kept Alpine forests open in a mosaic that limited fire spread whatever the weather. As that regime retreated, the fuel bed filled in again \citep{bebi_changes_2017}. Forest expansion is not uniform. It is fastest where recent agricultural abandonment coincides with climatic conditions favourable for tree establishment \citep{garbarino_contrasting_2020}. Anticipating future fuel continuity therefore requires knowing when land use changed relative to growing conditions, not only where abandonment occurred. The view that abandonment always increases fire hazard should still be treated as a regionally supported finding rather than a general rule. For disturbances other than fire, the reported effects of abandonment are inconsistent and sometimes opposite in direction \citep{mantero_influence_2020}.

Expansion also has a vertical component. In one intensively surveyed Central Swiss valley, the treeline advanced by an average of 10\,m per decade between the early 1970s and 2012, with local maxima of 42\,m per decade. Of the 765 cells of a 50\,m grid in which trees had established since 1972, only 21\% marked an advance beyond the former treeline, while 79\% were infilling below it \citep{frei_evidence_2023}. Fuel is therefore densifying as well as moving upslope, and it is this infilling below the treeline that turns scattered stands into continuous fuel. At Alpine scale, 24\% of the 10.5 million hectares of forest is in an open-canopy state, concentrated at high elevations, forest edges and warm dry sites \citep{stritih_alternative_2023}. These are the same positions where fuel is now filling in.

Where regrowth happens depends on economics and varies locally. A spatially explicit analysis of the Swiss mountains found forest regrowth concentrated where cultivation costs were high and yield potential low, with non-linear relationships to slope steepness and soil stoniness \citep{gellrich_agricultural_2007}. The study concluded that general policy measures applied across the whole mountain area cannot prevent abandonment and regrowth, which instead require attention to local characteristics. Swiss forest area has risen by roughly a third over the last 150 years, mostly through natural regrowth rather than planting. Fuel accumulation in the Alps therefore follows agricultural marginality, which is also why it concentrates on the steep ground where fire runs longest.

A multi-millennial record from the Apennines shows that the present composition of a European mountain forest can reflect land use rather than climatic equilibrium. At 1{,}390\,m in the Northern Apennines, species-rich mixed silver fir forests dominated from about 10{,}500 to 5{,}500 years ago, with fir making up to half the pollen assemblage. Beech then expanded and eventually replaced them entirely. Charcoal influx and temperature each explained about a tenth of the variation in the vegetation record. Fir, elm, ash and lime all declined strongly with burning, while beech tolerated light to moderate fire and declined under heavy fire \citep{moralesmolino_longterm_2021}. The study concluded that today's monospecific beech forests result from multi-millennial land-use intensification combined with late Holocene cooling, and that the displaced mixed fir forests were the more drought-tolerant system. The Alpine fuel bed has the same kind of origin. In both ranges the stand that now carries fire reflects how people used the land, so the correction must also come through land use.

The denser fuel bed has a direct effect on fire behaviour. Across 21 Swiss pyroregions, potential fire-path length, which combines slope length with fuel connectivity, explains between 56 and 65\% of the variance in large-fire statistics \citep{conedera_linking_2024}. Forest expansion onto the steep marginal ground that was abandoned first therefore increases exactly the quantity that matters. In this way a landscape can become more fire-prone without any change in weather.

The amount of fuel has been measured destructively in one Alpine country. Across 93 plots in Austrian Norway spruce, Scots pine and black pine forests, fuel load depended on region and forest type rather than on stand age \citep{neumann_fuel_2022}. Dead surface fuel increased and live understorey fuel decreased with basal area, and spruce stands carried the highest loads. Black pine on the dry eastern lowland margin, the forest type that burns most often in Austria, carried the lowest load. It combined open canopies with more shrub and grass cover and the deepest fuel bed, and the authors take this to indicate that fuel structure, together with the drier eastern climate, contributes to its fire frequency. The measurements leave out the layers that matter most once the fuel bed dries. The unsampled duff layer is estimated, from the literature and from Austrian forest-floor inventories, to add between 15 and 70\,t\,ha$^{-1}$ that can burn and smoulder under drought. Canopy fuel, roughly 10\,t\,ha$^{-1}$ in spruce stands, is consumed only in crown fires, which are currently rare in central Europe. A drier climate would bring both reserves into play.

Fire in the Alps also enters forests that are already shaped by other disturbances, and these disturbances interact rather than simply add up. Wind, snow break, bark beetles and avalanches shape Alpine stands continuously \citep{bebi_changes_2017}. Roughly half of disturbed stands passed into the open-canopy state described above \citep{stritih_alternative_2023}. A European baseline of disturbance regimes, resolved by causal agent, is now available. Classification of more than 3.8 million fire, wind and bark beetle events in European forests between 1985 and 2023 separates fire regimes along a north-south gradient \citep{miguel_unravelling_2026}. Fire is severe and rare in boreal and temperate forest, and moderate, large and frequent in Mediterranean and transitional forest. Wind and beetle disturbance fall into four regimes, of which the frequent and small regime predominates in Central Europe.

For Alpine assessment, the most important result is that these regimes are unstable. About 15\% of grid cells changed regime within the record \citep{miguel_unravelling_2026}. Fire frequency and severity fell in temperate and Mediterranean forest between the first and second periods, and wind and beetle disturbance then became more severe in temperate forest between the second and third. A disturbance regime is therefore not a fixed regional property on which a fuel model can be built once.

Susceptibility to disturbance also appears to depend partly on management, not only on climate. The MANFRED programme, a transnational Alpine Space collaboration of fifteen partners in six countries, noted, drawing on earlier European work, that recorded storm damage is not driven by climate alone \citep{jandl_management_2013}. The prevailing practice of letting increasingly dense forests mature has itself produced stand types that are highly vulnerable to storms. Fuel build-up and storm vulnerability thus share a common origin in forest management.

Fire differs from the other disturbance agents in what it leaves behind. Windthrow appears to leave biological legacies that continue to protect the slope. In the one Alpine stand where this has been simulated explicitly, rockfall modelling over a windthrown site in the Dolomites found that the barrier effect index improved after the storm \citep{costa_biological_2021}. The fallen stand was more protective than the forest before it, because lying logs and stumps act as physical obstacles rather than as energy absorbers. This legacy can persist for up to thirty years, until decay and soil contact degrade it. Fire does not leave such a legacy. Burnt stems and rootstocks can fall during and after the fire and carry rock with them, as observed at a single Austrian site \citep{melzner_brief_2019}, and in fire-injured beech stands most damaged stems collapse within about two decades \citep{maringer_temporal_2016}. Simulations that remove the vegetation entirely show the corresponding loss of the barrier effect \citep{asteriou_assessing_2025}. Two disturbances with similar canopy loss can therefore have opposite short-term effects on slope protection, so fire should not be treated as interchangeable with the other Alpine disturbances.

Salvage decisions after fire therefore cannot simply adopt the consensus developed after windthrow. That consensus advises against salvage in order to retain protective legacies \citep{costa_biological_2021}, but its evidence base is less clear-cut than often assumed. Eighteen years after a Tatra windstorm, salvaged areas retained an eighth of the deadwood volume of unsalvaged areas but supported higher regeneration diversity, 5.40 against 3.85 species per plot, while unsalvaged plots were dominated by Norway spruce \citep{seben_impact_2024}. The trade-off between protective legacy, regeneration composition and future fuel is real and differs between disturbance types, and it has not been evaluated for Alpine post-fire sites.

A comparative case from outside Europe shows both the parallel and its limits. In the southern Appalachian Mountains, decades of suppression preceded a regional outbreak of more than a dozen large wildfires during one of the warmest and driest droughts of the last century \citep{reilly_drivers_2022}. The historical mosaic of that landscape had been maintained by frequent low- and mixed-severity fire, and land-use change had already reorganised its composition. The Alps never had a frequent-fire regime to suppress, so the trajectory differs. The comparison shows only that a landscape whose fuel structure no longer matches its disturbance history can support a regional outbreak once drought provides the trigger.

Which species grow back matters as much as how much forest grows back, because species composition both shapes fire and is shaped by it. Alpine tree species differ sharply in their capacity to survive fire, and Alpine forest composition is changing on a timescale comparable to that of the fire regime. Simulation of Swiss Alpine landscapes indicates that climate-driven vegetation shifts will influence biomass and species distribution more than management will \citep{schumacher_relative_2006}. In those simulations, fire emerges as a property of the landscape and then feeds back on vegetation.

Fire vulnerability depends on tree traits, and this now has independent experimental support. Simulation across 4438 French Alpine forest plots showed that cold-season fires, about half of all fires there, had limited impact except on low-elevation southern stands \citep{dupire_vulnerability_2019}. Average summer conditions produced the highest mortality in mixed and broadleaved low-elevation forests, with up to 75\% mortality in coppice. Extremely dry summers produced high mortality in all communities, with the lowest mortality in the thick-barked, high-canopied larch stands of the inner Alps. Controlled heating of stem discs from six tree species common in Central Europe reproduces this ranking experimentally \citep{mutterer_bark_2026}. Bark thickness appears to be the most robust predictor of the time the cambium takes to reach the lethal threshold. It places larch, Douglas-fir and Scots pine above silver fir and Norway spruce, with beech lowest, and thermal resistance declines with bark moisture and with height on the stem. Because the two methods are independent, their agreement is a stronger basis than either alone.

The relationship also runs from fire to composition. After a stand-replacing fire in the Rhone valley, recruitment favoured pioneer birch, aspen and willow over the pre-fire larch, spruce and Scots pine \citep{moser_potential_2010}. Seedling density tracked moisture availability rather than fire intensity, which led to the conclusion that drought limits post-fire recruitment in the Central Alps. Palaeoecological records extend the same logic over millennia. Fire and grazing, rather than temperature alone, determined whether the timberline reached its climatic potential \citep{colombaroli_species_2010}. At two forest sites in the central Dinaric Alps of Montenegro, fire and land use likely reduced old-growth fir-spruce-beech forest during the Middle Ages \citep{cagliero_effect_2022}.

\subsection{Protective function, biodiversity and the recovery window}

The consequences of fire on an Alpine slope depend on what fire removes, how long it stays removed and what else changes in the meantime. Among the post-fire processes involved, rockfall is the most specific to the Alps and the least studied internationally. As recently as 2019, no work on post-wildfire rockfall activity had been published for the eastern Alps, and even in the better-studied western United States the process had received little attention compared with debris flow \citep{melzner_brief_2019}.

The thermal effects on rock begin at known temperatures. A review of three Spanish cases shows fire acting at the same time on the source area, the propagation zone and the affected area \citep{sarro_effects_2021}. In the source area, mineral expansion, dehydration and decarbonation open microcracks. Granites lose strength appreciably above 500\degree{}C as quartz approaches its phase transition at 573\degree{}C, and limestones alter from about 300\degree{}C. In the propagation zone, the loss of vegetation removes the energy absorption that reduces both the reach and the frequency of falling blocks. In the affected area, heat degrades the cables and barriers installed to stop rockfall. The same review notes that the restitution and friction coefficients used in rockfall trajectory models have not been standardised for post-fire conditions. The parameters needed to model a burnt slope are therefore the ones least constrained by measurement.

Vegetation affects rockfall on two timescales. The litter and ground vegetation that stop small blocks are lost immediately, while the standing stems that intercept falling material are lost more slowly. Exposure can therefore rise for years as fire-killed trees decay and fall. Trajectory modelling isolates the effect of vegetation and shows that it depends strongly on block size \citep{asteriou_assessing_2025}. Small blocks of around 0.01\,m$^3$, which mostly came to rest on a vegetated slope, travel much further once cover is removed. For blocks an order of magnitude larger, kinetic energy overwhelms the resistance of vegetation. Protection is thus lost first against the frequent, small events that dominate long-run exposure.

The rate and duration of this loss have been quantified for beech on the southern side of the range, and the results are more severe than the general disturbance-window literature implies. Across 244 plots in burnt beech stands of Ticino and Piedmont, where beech grows mainly between 600 and 1{,}700\,m, root reinforcement fell by roughly 70\% within five years of high-severity fire and by 98\% by year fifteen \citep{gehring_shallow_2019}. Moderate severity produced about 15\% loss in the first five years and 60\% by year twenty-three, while low-severity burns remained stable. After high-severity fire, modelled protective capacity against shallow landslides fell below a quarter of the unburnt value regardless of soil and slope class, and landslide probability stayed elevated for at least forty years. Rockfall protection follows a similar course on a different timescale. Across 34 burnt beech stands in the southern Alps, low-severity burns retained close to unburnt protective capacity, while moderate and high severity reduced it significantly, most strongly between ten and thirty years after the fire \citep{maringer_temporal_2016}.

These results change how the consequence window should be understood. Protective capacity reaches its lowest point years after the fire rather than immediately, because standing dead stems keep working until they fall and regeneration has not yet replaced them. An assessment made in the first year after a fire therefore measures the wrong part of the recovery curve. For risk assessment, we take the key parameter to be how long post-fire hazard remains elevated, rather than whether it is elevated. The geomorphological literature gives a general disturbance window of three to four years, occasionally extending to thirty \citep{santi_wildfire_2022}. The forty-year figure for shallow landslide disposition exceeds the usual three to four years by an order of magnitude. This is the clearest available evidence that the short end of the general range is optimistic for the Alps, and that the Alpine window is set by the recovery of biological structure rather than by the decay of a soil property.

Independent lines of evidence support this conclusion. Soil properties remained altered three years after an Alpine mixed-severity fire, including properties that usually recover quickly \citep{negri_wildfire_2026}. Post-fire seedling establishment, the first step towards a protective canopy, was drought-limited at the low and medium altitudes of one Valais burn, where successful establishment required precipitation far above the long-term average \citep{moser_potential_2010}. At Alpine scale, 72\% of disturbed forests return to a closed-canopy state within 35 years, but submediterranean forests remain open for decades \citep{stritih_alternative_2023}. At the upper margin the window may not close at all, although this observation comes from a single massif. On burnt subalpine slopes in the Northern Limestone Alps at 1800--1950\,m, tree seedlings were nearly absent and no meaningful mountain pine re-establishment occurred in about eighty years \citep{degiuli_arrested_2024}. This state is better described as arrested succession than as delayed recovery. The length of the window is therefore partly a management outcome, and current Alpine restoration practice is poorly matched to it (Section~\ref{sec:fuelbreaks}).

The protective function has been valued at one site, and it is itself exposed to climate change. For a protection forest above Martigny in the Valais, coupled forest-dynamics and three-dimensional rockfall simulation places annual rockfall risk roughly five times higher without the forest than with it \citep{moos_climate_2021}. Under climate scenarios, the risk rises within 150 years as composition shifts from beech to pubescent oak, and small blocks reaching the railway move from once per decade to once every three years. The monetary values are given in Section~\ref{sec:economics}. Simulation along an elevation gradient in South Tyrol shows that the loss of protection depends on elevation \citep{hillebrand_contrasting_2023}. Under severe warming, montane forests decline drastically in avalanche protection, with drought-induced biomass reductions reaching 97\% in dry inner valleys, while subalpine stands remain stable or improve.

Fire in the Alps is not, however, uniformly destructive, and the review literature openly warns against assuming otherwise. A systematic assessment of mountain wildfire against fifty ecosystem services, imported from the Colorado Front Range of the Rocky Mountains, found that current fire regimes mostly reduce service provision and that effects on cultural services are mixed \citep{vukomanovic_systematic_2019}. The dominant mechanism is change in vegetation composition and structure. The same assessment suggests that a bias towards negative characterisation may be built into how the research questions are posed.

Alpine forest can now be placed in a European ranking of ecological vulnerability. A continental assessment combining ecological value, the resistance of vegetation to fire and post-fire recovery time ranks temperate conifer forest second among European biomes, behind tundra and ahead of boreal forest \citep{arrogantefunes_mapping_2024}. Its high rank is driven by recovery time rather than by resistance. Non-forested vegetation, including the grassland, heathland and shrubland of the subalpine and alpine belts, is assigned no coping capacity in that framework. This is a modelling convention rather than a measurement, but it follows from the absence of the bark and crown traits on which resistance is defined. Taken together, these results place the Alpine problem in recovery. What fire removes is not resisted at the time of the fire and returns slowly afterwards, so the total ecological value lost over the recovery period can be large even where the burnt area is small.

Abundance and composition respond differently to fire, and confusing them leads to wrong conclusions in either direction. At two burnt sites in the inner-Alpine Valais, high-severity fire roughly halved bird abundance (21.68 against 41.59 territories per 10\,ha) without significantly changing species richness \citep{rey_effects_2019}. Fourteen species were more abundant in burnt stands and five in unburnt ones. Red-listed and priority species were substantially more common in burnt stands, and this effect strengthened with time since fire, with many species peaking eight to twelve years after the event and declining as the vegetation became denser. Fire at these sites removed individuals and rearranged the community rather than eliminating species, and the species it favoured were disproportionately those of conservation concern. Vegetation shows the same pattern. The burnt subalpine slopes in the Northern Limestone Alps described above carried significantly higher vascular plant richness than unburnt controls, including alpine specialists absent from the unburnt plots, even though tree regeneration failed entirely \citep{degiuli_arrested_2024}.

Neither measurement captures animal mortality. Individual animals are killed by burning, smoke inhalation and heat stress \citep{demiguel_fires_2021}. The most exposed taxa are those that cannot outrun a fire front or that depend on structures the fire removes, such as small mammals, cavity-nesting birds, reptiles, and aquatic organisms affected by post-fire sediment delivery. Hollows used for denning and nesting are lost with the stems that contain them, and poor dispersers are slowest to recolonise. These mechanisms are general rather than Alpine, and the only order-of-magnitude figures available come from a single non-European event. No Alpine estimate of faunal mortality exists in the literature assembled here, so faunal losses in the Alps are unmeasured rather than known to be small.

These findings refine the risk framing of this review rather than weaken it, because they show that Alpine losses have several components that behave differently. The protective function of forest on slopes above assets is spatially concentrated, can be valued in money \citep{moos_climate_2021} and is slow to recover \citep{stritih_alternative_2023, degiuli_arrested_2024}. The ecological value lost over that recovery period is high for the forest type that dominates the Alpine belt \citep{arrogantefunes_mapping_2024}. For biodiversity, fire is negative for abundance and for closed-forest specialists and positive for open-habitat and disturbance-dependent species \citep{rey_effects_2019, degiuli_arrested_2024}. Faunal mortality, by contrast, remains unquantified. Combining these components into a single measure leads either to treating fire as uniformly destructive or to treating it as ecologically benign because some species benefit, and both errors appear in the wider literature.

\subsection{Managing for protection: silviculture, species and fuel treatment}

In Alpine forestry the principal product is often not wood. Mountain forests serve a protective function against natural hazards \citep{dinca_timber_2025}, and across large parts of the range management aims to maintain a stand structure that can stop rockfall, anchor the snowpack against avalanche release and hold soil on steep ground. Timber is a secondary output, since in difficult terrain harvesting costs are often not covered by timber revenue \citep{dinca_timber_2025}. The priority of protection is long established. The same programme that linked storm vulnerability to management practice described protection as by far the most important ecosystem service of mountain forests \citep{jandl_management_2013}. It identified eco-engineering as the response where climate change puts that function at risk by increasing vulnerability to rockfall, avalanches and landslides.

The protective service can be given a monetary value. A replacement-cost framework harmonised for the Alpine Space values it by the cost of the engineered structures that would otherwise be required, designed to European standards and net of the management expenditure needed to sustain protective capacity \citep{accastello_asforesee_2019}. In a single demonstration case, a ten-hectare Scots pine stand in the Italian Alps assessed over a twenty-five-year horizon, the forest met local safety requirements without supplementary artificial measures. The avoided-risk valuation of the Valais protection forest described above \citep{moos_climate_2021} points the same way. Under either accounting convention the protective value is substantial, whereas the thinning needed to sustain it often has a negative stumpage value on steep slopes \citep{accastello_asforesee_2019}, and it is this protective value that fire puts at risk. The monetary figures from both studies are given in Section~\ref{sec:economics}.

This creates a management problem. Protective capacity is not a fixed attribute of a stand but depends on a structure that must be actively maintained, and under climate change the prospects for maintaining it depend on elevation. The South Tyrol simulations described above project the steepest losses in montane forests, while subalpine stands hold or improve their rockfall protection \citep{hillebrand_contrasting_2023}. Silvicultural effort therefore needs to be concentrated where the protective function is most exposed, which is not where the forest is most productive.

Fire also sharpens conflicts between management objectives. Scenario modelling of Austrian forests over the century finds stem biomass declining by between 9\% and 29\% under harvest-oriented strategies and rising by 34\% under a reduced-harvest, conservation-oriented strategy \citep{jandl_strategies_2018}. Long-lived wood products emerge as the option that best balances carbon sequestration, economic return and disturbance risk. The reduced-harvest option maximises stored carbon but also maximises standing fuel, and increasing forest cover extends the area over which fire is possible \citep{gerberding_mapping_2025}. Carbon policy and fire policy therefore pull in opposite directions in Alpine forests, and no Alpine study reviewed here evaluates that trade-off explicitly.

If protective structure is the management objective and fire is an emerging threat to it, species choice becomes a fire management decision. Alpine forestry is already making that decision for other reasons, and its direction appears to be set by climate rather than by fire. The Swiss simulations cited above place climate-driven change in composition ahead of management as the main driver of future biomass and species distribution \citep{schumacher_relative_2006}. Upward shifts of altitudinal vegetation belts would require site types to be reassessed, and the treeline advance and infilling recorded in one Swiss valley show that this movement is already under way \citep{frei_evidence_2023}. Species suitable for a site today are not necessarily those suitable in fifty years, so Alpine silviculture is moving towards mixtures and towards provenances from warmer analogue sites.

Fire interferes with this trajectory in a specific way. The species favoured by warming and drought at low and mid elevations, and those that colonise after fire, are largely not the species that currently provide protective structure. Post-fire recruitment in the Rhone valley favoured pioneer broadleaves over the pre-fire conifers, as described above \citep{moser_potential_2010}, and the Valais protection forest is projected to shift from beech to pubescent oak under warming, with rockfall risk rising as it does \citep{moos_climate_2021}. Pioneer broadleaves establish faster than the conifers they replace but hold rock less effectively at the relevant stem diameters, and the transition period coincides with the consequence window described above.

Where fire resistance can be built into species choice, the evidence is reasonably firm, because stand-level mortality simulation and stem-disc heating give the same ranking \citep{dupire_vulnerability_2019, mutterer_bark_2026}. Larch and Scots pine, two of the three species ranked highest, are already major components of Alpine protection forest. This is a more favourable starting position than the fire intolerance of the range as a whole would suggest.

The same programme also evaluated what sits alongside species choice. Provenance is a lever in its own right. Long-term trials for Norway spruce indicate that matching provenance to anticipated future site conditions, rather than to present ones, can substantially raise potential productivity and reduce risk over the stand lifetime \citep{jandl_management_2013}. Browsing, by contrast, is a constraint. Adaptive silvicultural concepts can only be implemented where browsing damage from high ungulate densities is held to acceptable levels, and at the densities then prevailing, deer were inhibiting the establishment of species-rich forest. Species diversification is the central recommendation of Alpine adaptive management, and ungulate management is its precondition. A wildlife-management decision therefore sits upstream of fire resilience.

The unresolved tension lies with the other objectives that silviculture must serve at the same time. Reduced-harvest strategies maximise carbon storage while building standing biomass, as noted above \citep{jandl_strategies_2018}. Regeneration-oriented approaches to mountain forest management call for actively restoring degraded systems rather than maintaining existing conditions \citep{majlingova_sustainability_2025}. The harvesting operations needed to implement any of this are technically and economically constrained in steep terrain \citep{dinca_timber_2025}. No Alpine study reviewed here optimises species selection jointly for protective function, carbon, timber and fire resistance. Until one does, the trade-off is made implicitly at every planting decision.

Alpine fuel load can be reduced mechanically, by grazing or by fire, and in this setting the evidence favours mechanical treatment and grazing. Mechanical treatment has a justification independent of fire. One review of drought-induced Scots pine mortality suggests that thinning and shrub removal could ease drought stress by reducing competition for water, improving carbon uptake and lowering mortality risk, although whether thinning mitigates drought in the most sensitive areas remains unclear \citep{giuggiola_droughtinduced_2010}. The same operations reduce surface and ladder fuel and interrupt canopy continuity. In inner Alpine valleys, drought-driven Scots pine mortality is the main forest health problem and fuel accumulation is the emerging fire problem. Treating one largely treats the other, so the case for thinning does not depend on accepting a projection of future fire.

Grazing offers the broadest policy leverage. A systematic review of the herbivore literature concludes that fuel reduction by grazing animals is a promising way to mitigate wildfire where land abandonment has led to shrub and tree encroachment, with effectiveness depending on herbivore type, density and feeding behaviour \citep{rouetleduc_effects_2021}. The same review finds that agricultural and forest policy could support grazing far more effectively than they do. Policy could regulate and promote herbivory in forest areas from which animals have traditionally been excluded to prevent damage, and could direct agri-environmental subsidies towards extensive grazing and pastoralism. Landscape simulation quantifies the gain from embedding fuel treatment in land-use policy rather than applying it alone \citep{pais_optimizing_2023}. Prescribed burning alone reduced subsequent wildfire by up to 36\% over thirty years, but the leverage of treatment reached unity (one hectare of wildfire avoided per hectare treated) only when combined with restoration of high nature value farmland and fire-smart forest conversion.

The available effect sizes for grazing are large, and all of them come from outside the Alps. Managed grazing has reduced rate of spread by more than 60\% in pure grass communities and by over 50\% in mixed grass and shrub, and Mediterranean studies report burnt-area reductions between 25.9\% and 60.9\% \citep{xu_how_2026}. Grazing works by creating patchiness rather than by removing fuel uniformly. Because it interrupts horizontal continuity, it behaves more like a distributed set of small breaks than like a thinning. Its effect is non-linear and reverses at high intensity. Removing less than about half the photosynthetic leaf area reduces fuel and favours fire-resistant species. Heavier use, particularly during drought, increases erosion, admits flammable annuals and can leave a continuous, homogeneous fuel bed more flammable than the one it replaced. On erodible Alpine slopes above assets, overgrazing may carry a second cost by removing low, dense ground cover, which retained small blocks in post-fire rockfall simulations from Mediterranean shrubland in Greece \citep{asteriou_assessing_2025}. The threshold between the two regimes has not been located for any Alpine pasture. Timing adds a further complication, since grazing in high-altitude systems may be most effective in the drier spring and autumn rather than in the humid summer, when Alpine stock is at its highest elevation \citep{xu_how_2026}.

The only Alpine evidence on grazing and fire is historical and correlational. In the Valais, livestock density in forests was associated with fewer fires before the late 1950s, while wood pasturing remained widespread \citep{zumbrunnen_human_2012}. After 1956, high livestock densities coincided with more fires. Those densities lay in the driest eastern part of the canton, where dryness rather than fuel removal is the more likely explanation.

Where treatment is placed in the landscape matters as much as whether it happens. Landscape modelling of a transboundary Iberian biosphere reserve over 2010--2050 found that restoring high nature value farmland saved roughly €3.9 million in suppression costs in present-value terms \citep{lecinadiaz_incorporating_2023}. Combining it with fire-smart forestry cut total societal cost (suppression plus ecosystem service damage) from €25.7 million under business as usual to €7.1 million. Fire-smart forestry applied alone, by contrast, \emph{raised} suppression costs, so the saving came from the agricultural-landscape component. In the Alps, abandonment of agricultural land is the proximate cause of fuel accumulation. This is therefore the most directly relevant economic result available, and it suggests that the effective lever is agricultural policy rather than forest policy alone.

Empirical support for this argument now exists at national scale, though not for the Alps alone. An analysis of about 49{,}000 Italian fires larger than a hectare between 2007 and 2017 covered 602 agro-forest districts spanning Alpine, temperate and Mediterranean climate regions \citep{spadoni_active_2023}. Districts with more active land governance showed smaller burnt shares of ecosystem services and of the wildland-urban interface, and lower forest burn severity. The association held after accounting for climate, fire weather and simulated flammability, while land abandonment increased all three impacts. Governance was represented by proxies, namely sustainable forest management and certification, extensive agro-pastoral activity, protected areas, and rural development and LIFE funding. Its effect was moderate compared with climate and flammability, and it did not always act in the same direction. Severity fell with the governance of rural and marginal areas up to the median value and rose beyond it, a reversal the authors relate to intensive harvesting.

The effectiveness of treatment also depends on the fire weather under which it is tested, and this dependence works against the direction in which the Alps are moving. A randomised, replicated experiment in California that later burned in a wildfire found that combined thinning and prescribed burning reduced the probability of tree mortality by 0.60, crown volume consumed by 86\% and bole char height by more than seven metres relative to untreated controls \citep{brodie_forest_2024}. The mortality benefit was greatest at intermediate values of the burning index and declined under the most extreme fire weather, although crown-fire metrics remained responsive. Fuel treatment therefore achieves most where danger is moderate. Land governance shows the same dependence at landscape scale, with its mitigating effect strongest in warm, dry and flammable districts and weakening as fire weather during the events became more severe \citep{spadoni_active_2023}. Because Alpine fire danger is currently low over most of the range but projected to rise towards moderate levels \citep{gerberding_mapping_2025, miller_climate_2024}, the window in which treatment is most cost-effective is now opening.

\subsection{Prescribed fire, fuel breaks and post-fire restoration}
\label{sec:fuelbreaks}

For each of these instruments, the Alpine evidence is thinner than the confidence with which it is discussed. Prescribed burning is also unevenly practised. It is applied at scale in France, where 1{,}000 to 3{,}000\,ha are treated annually in the Alpes-Maritimes by trained technicians working to shared protocols \citep{muller_forest_2020}. This happens in the Mediterranean south-east rather than in the Alpine region proper, and serves grazing recovery as much as hazard reduction. The technique, the training structure and the institutional model therefore exist within an Alpine country but have not been transferred to Alpine conditions.

Prescribed burning has been regarded as a poor fit for the Alps, and the structural reasons are real. The dominant species have little fire adaptation. Silver fir, Norway spruce and especially beech rank lowest in experimentally measured bark insulation \citep{mutterer_bark_2026}, and simulations for the French Alps give thin-barked beech and silver fir the highest post-fire mortality \citep{dupire_vulnerability_2019}. Prescribed burning is little used or prohibited in most Alpine countries, and the early-spring conditions under which it would be carried out are those of the March--April peak of human-caused fires \citep{muller_forest_2020}, so we expect the usable burn window to be narrow. The consequences of an escape are severe, because the slopes that need treatment lie above assets. In Austria the obstacle is also legal, since the Forestry Act prohibits lighting fires in forest and prescribed burning is not considered a management option \citep{neumann_fuel_2022}. Burn severity is also a poor guide to soil effects. After one Alpine wildfire in Scots pine, the most water-repellent surface soils lay in low-severity plots, although fire had on the whole lowered the naturally high repellency of these soils \citep{negri_wildfire_2026}.

An Alpine trial now exists, however, and its results are more favourable than the structural argument would predict. In a south-west Alpine Scots pine stand at 1600\,m on a south-facing slope, eighteen treated and control plots were monitored at four points over a year, with burning applied under manipulated fuel-intensity gradients \citep{gamba_prescribed_2025}. Fuel load fell by roughly half. There was no significant change in branch hydraulic conductivity or stem water potential, no heat-induced anatomical damage, and minimal change in soil organic carbon or water repellency, even though the trial spanned a severe drought. On this evidence, prescribed burning is an ecologically sustainable fuel-reduction option for this forest type.

Evidence from soil organisms points the same way, though from outside Europe. In paired southern Appalachian watersheds, the bacterial and fungal communities of repeatedly prescribed-burnt stands were statistically indistinguishable from those of fire-excluded stands, while wildfire-affected and partially harvested stands differed significantly from both \citep{snyder_taxonomic_2026}. Ectomycorrhizal abundance held at 63 to 66\% under prescribed fire and exclusion but fell to 38\% after wildfire, with saprotrophs rising correspondingly. The comparison is confounded by differences in time since disturbance and in elevation, and the forest type is not Alpine, so the result cannot be transferred directly. It does, however, identify the soil microbial community as a measurable endpoint that clearly separates prescribed fire from wildfire, and the Alpine prescribed-burning question currently lacks such an endpoint.

Meanwhile, one European mountain range has moved beyond the trial stage, in a setting closer to the Alps than the Mediterranean cases usually cited. The Aran valley in the Spanish Pyrenees covers some 634\,km$^2$, with a mean altitude of 1{,}840\,m and a mean slope of 27\degree{} \citep{oliveres_path_2025}. It has an Atlantic rather than a Mediterranean climate and a winter-dominant fire season from January to March. In 2022 the valley unanimously adopted a strategic fire management plan. The plan is built around an explicit target for how much fire the landscape should carry, set at 0.66\% of natural land cover annually against a historical fire rotation of about 350 years. Its first implementation burnt 327\,ha over fourteen working days inside a pre-planned 346\,ha management area. The valley's winter and early spring burning windows, its transhumance history in which shepherds used fire routinely, and its deep organic soils, in which fire can smoulder and reignite, all transfer directly to the Alps. The organic soils matter because smouldering ground fire is the suppression problem that Alpine lightning fires present.

The Pyrenean case also shows that adoption depends mainly on things other than ecological evidence. The plan rests on a community participation process with municipal representatives, on eight formally agreed goals covering severity reduction, emergency resilience, extensive livestock farming and biodiversity, and on landscape fire succession modelling to locate treatments \citep{oliveres_path_2025}. Its authors acknowledge that economic sustainability has not yet been assessed, that managed fire can escape, and that explicit wildfire use under technical and strategic direction remains far from common European practice.

The structural objections and the favourable trial can be reconciled, and doing so defines a research agenda. The favourable result comes from a fire-tolerant species, at a single elevation, on one aspect, at plot scale, over one year. The structural concerns apply to fire-intolerant stands, to operational scale, to escapes and to multi-year soil response. The trial shows that the question can be answered empirically in the Alps, not that it has been answered. The structural argument shows where the next trials should not begin.

Outside mountains, the standard landscape instrument for limiting fire size is linear, and its forms are distinguished by how much vegetation survives treatment. A firebreak is a strip cleared to mineral soil. A fuelbreak retains herbaceous and shrub cover while reducing fuel volume and flammability. A shaded fuelbreak also retains the tree cover, with grass, shrubs and trees all kept at reduced density \citep{ascoli_firebreak_2018}. The distinction matters because it determines whether the instrument is compatible with an Alpine protection forest.

On that criterion, a form suited to mountains already exists and is recommended. Linear and shaded fuelbreaks are considered better suited than cleared firebreaks to mountainous terrain with steep slopes and rough morphology, because retaining vegetation cover limits soil erosion \citep{ascoli_firebreak_2018}. Where terrain is very steep, ridgelines and valley bottoms are the preferred locations. Recommended widths run from 30 to 100\,m for primary networks and from 3\,m to over 50\,m for secondary and tertiary ones, and siting criteria already include maintenance cost, road access and firefighting considerations. European Mediterranean practice adds detailed prescriptions at the interface \citep{rigolot_recommendations_2003}. It distinguishes an inner zone, where trees are removed, from an outer zone, where they are retained but pruned and thinned, with thresholds for shrub phytovolume and crown separation. The Alpine literature reviewed here contains no equivalent design guidance calibrated to Alpine fuels, slopes and protective requirements, which is a gap in application rather than in principle.

What such a network could deliver in the Alps is nonetheless limited, and the limits can be quantified. Linear barriers are less effective than their prominence in policy suggests. A recent synthesis reports that firebreaks stopped only 28\% of large fires in southern California between 2017 and 2020, and that some 73\% of wildfires in western Arnhem Land, northern Australia, breached them fully or partially between 1998 and 2001 \citep{zong_effectiveness_2026}. A proposed national network in Portugal is projected to reduce fires above 500\,ha by 13\% but would still be breached in more than a third of such events. Firebrand transport across the gap is the dominant failure mechanism.

Where fuelbreaks do work, the reason runs against the usual understanding of what they do, and it points to the constraint that matters most in mountains. Analysis of 563 fire-fuelbreak intersections from 220 wildfires in Andalusia found containment in 46.9\% of cases overall \citep{ortega_modeling_2024}. Success depended strongly on whether suppression was operating on the line, at 6.66\% with no suppression, 15.9\% with aerial resources alone, and 76.74\% with ground crews working the line. A fuelbreak is therefore a platform for ground suppression rather than a barrier in its own right. Its value depends on crews reaching and holding it. In Alpine terrain, the slope, access and exposure conditions that motivate building a fuelbreak are the same conditions that prevent ground crews from occupying it. The terrain-conditioned suppression-difficulty framework discussed in Section~\ref{sec:governance} is the natural tool for identifying where this constraint applies, but it has not been applied to Alpine fuelbreak siting.

Geometry adds to the problem. In the Andalusian data, effectiveness depended strongly on the angle at which fire met the line \citep{ortega_modeling_2024}. Near-parallel intersections performed best, and without suppression support perpendicular intersections held in only 5.6\% of cases. Alpine fires run upslope along paths whose length predicts large-fire likelihood \citep{conedera_linking_2024}, so a break aligned with the contour meets the fire run perpendicularly. Siting on ridgelines \citep{ascoli_firebreak_2018} partly resolves this, because a ridge intercepts the run at its end rather than across its flank. It also, however, places the treatment where it is most exposed to wind and least accessible to ground crews.

Even a well-sited network would then face obstacles of ownership and finance. Large, long linear networks are considered poorly suited where the wildland-urban interface is widespread and forest property is largely private, with small cadastral parcels \citep{ascoli_firebreak_2018}. This describes much of the Alpine arc. The Catalan experience quantifies the resulting shortfall. After three funding calls, prevention subsidies had reached under 4\% of the designated strategic area, and funded plots within groups lay an average minimum distance of 287\,m apart, so the spatial continuity the scheme required was achieved only in some cases \citep{gorrizmifsud_current_2026}. The financing model does not transfer either. Optimisation of a 3{,}538\,km network in the north-western United States found that about 97\% of the available network area required thinning with pile burning, with estimated net harvest revenue of US\$63--75 million before planning, road and maintenance costs \citep{ager_optimizing_2023}. Alpine harvesting in steep terrain is technically demanding and economically marginal \citep{dinca_timber_2025}, so such revenue is largely unavailable.

A trade-off specific to protection forests remains, and the shaded form does not remove it. Even a treatment that retains the canopy reduces the density of standing stems that slow falling blocks \citep{moos_climate_2021} and removes the lying material that adds a barrier effect \citep{costa_biological_2021}. Because protective capacity can be valued by its replacement cost \citep{accastello_asforesee_2019}, with figures given in Section~\ref{sec:economics}, the acceptable degree of thinning is an empirical question with a monetary answer that no study has yet computed. On the evidence assembled here, the most tractable unanswered question in Alpine fuel management is how much understorey and stem removal a protection forest can absorb before its protective function declines faster than its fire hazard.

Some support for these instruments comes from vulnerability rather than suppression, and it applies at a different scale. The expert weighting described in Section~\ref{sec:risk} ranked protective strips and firebreaks highest among the building-environment indicators, but that judgement concerns the building rather than the landscape. A break next to a structure serves as a defensible margin, which is not the containment function evaluated in the effectiveness literature. Terrain that makes a landscape-scale line impossible to hold does not prevent a cleared margin around a house. Separating the two uses resolves the apparent contradiction between an instrument that performs poorly at landscape scale and one that practitioners rank first at the scale of the exposed asset.

In our reading, Alpine fuel management is therefore likely to be mainly area-based rather than linear, with linear elements confined to ridgelines, existing road corridors and the interface. The evidence supports grazing that reduces fuel across whole abandoned slopes \citep{rouetleduc_effects_2021}, thinning that removes fuel and may also ease drought stress \citep{giuggiola_droughtinduced_2010}, prescribed burning where the stand tolerates it \citep{gamba_prescribed_2025}, and prioritisation of the terrain positions from which large fires actually develop \citep{conedera_linking_2024}. Multi-criteria zoning of the kind developed for Catalonia offers one way to set such priorities \citep{krsnik_spatial_2024}.

Post-fire management determines how long the consequence window described earlier in this section stays open, and it is the least evidenced phase of the Alpine fire cycle. The deficit applies to fire science as a whole, not only to the Alps. Bibliometric analysis counts roughly 3{,}976 post-fire publications against 5{,}802 for prevention and 7{,}041 for active fire, so the recovery phase is around 46\% less studied than the pre-fire phase \citep{goncalves_fostering_2025}.

The deficit is also institutional, and it has now been described at European level. Post-fire management is the least developed component of the fire management cycle across the European Union \citep{vieira_postfire_2026}. The obstacle is a gap between available science and operational practice, not a lack of tools, since stabilisation measures, restoration techniques and monitoring protocols all exist. What is missing is a common assessment procedure, multidisciplinary operational teams, long-term funding and knowledge exchange between member states. That diagnosis draws on Greek, Portuguese and Spanish experience, so its transfer to the Alps is an open question. The Alpine case is the harder one, because the intervention must prevent gravitational rather than agronomic consequences, within a window set by the recovery of protective structure.

The existing Alpine evidence is critical of current practice. Analysis of short-term recovery after a large mixed-severity fire in the Western Alps concluded that the restoration techniques most often applied are poorly matched to the ecological context and to the consequences of global change \citep{mantero_short-term_2023}. Intervention can either help recovery or cause further degradation. Restoration must also work within the biological limits described above, since drought limits recruitment at low and mid elevations \citep{moser_potential_2010} and succession at the upper margin may be arrested altogether \citep{degiuli_arrested_2024}.

Salvage logging after fire is the decision on which the Alpine evidence is thinnest and the temptation to import evidence from windthrow is strongest. The windthrow literature advises retaining fallen stems for their barrier effect \citep{costa_biological_2021}. As discussed above, however, fire does not leave that legacy, and the salvage consensus is itself less settled than often assumed \citep{seben_impact_2024}. Post-fire salvage in an Alpine protection forest is therefore currently decided without directly relevant evidence, even though the decision determines both the immediate barrier effect and the fuel load available to the next fire.

\subsection{What differs in mountains}

Outside mountains, forest fire management generally aims to limit spread and to protect timber and settlements. Where fire has been excluded from systems that once burned frequently, fuel is commonly assumed to accumulate, yet the pathway is not universal. In the southern Appalachians, a range with a frequent-fire history, exclusion shifted forests towards fire-sensitive species and reduced flammability, and the large fires of 2016 followed an acute autumn drought \citep{reilly_drivers_2022}. In the general lowland and Mediterranean baseline, many of the species carrying that fuel have fire adaptations such as thick bark, resprouting or serotiny. Where the stand itself is the asset, thinning and prescribed burning reduce hazard and lower mortality in the trees they treat \citep{brodie_forest_2024}. In the Alps most of these premises do not hold, beginning with the asset.

In the Alps the main asset is often a protective function rather than a product \citep{jandl_management_2013}. Forest above a settlement, road or railway is managed to hold rock, snow and soil. Valued by replacement cost or by avoided risk, that protection has a substantial monetary value \citep{accastello_asforesee_2019, moos_climate_2021}, whereas the interventions that sustain it often yield negative stumpage \citep{accastello_asforesee_2019}. Fuel treatment thins the standing stems that slow falling blocks \citep{moos_climate_2021} and removes the lying material that adds a barrier effect \citep{costa_biological_2021}. Fuel reduction and protective capacity therefore trade against each other directly, and no study has yet established how much treatment a protection forest can absorb. This trade-off does not arise where the protected asset is the stand itself.

The Alpine fuel bed also has a different origin. It reassembled from the mid-nineteenth century as grazing, litter collection and fuelwood extraction retreated, with depopulation, afforestation laws and fossil fuels relieving the pressure on forests \citep{bebi_changes_2017}. Forest expansion is fastest where recent abandonment meets favourable growing conditions \citep{garbarino_contrasting_2020}, it concentrates on marginal steep ground \citep{gellrich_agricultural_2007}, and much of it is infilling below the treeline \citep{frei_evidence_2023}. Fuel connectivity and slope length together predict much of the variation in large-fire statistics across Swiss pyroregions, and if Switzerland were forested up to 2500\,m, 16 of 21 pyroregions would fall in the highest propensity class \citep{conedera_linking_2024}. This expansion is therefore likely to raise hazard independently of climate. Most of the evidence linking abandonment to fire, however, comes from twentieth-century Mediterranean Europe \citep{mantero_influence_2020}. Even so, the evidence suggests that the correction is agricultural and pastoral as much as silvicultural \citep{lecinadiaz_incorporating_2023, spadoni_active_2023}, which places the most effective policy lever outside forest administration.

The species carrying Alpine fuel are mostly fire-intolerant. Silver fir, Norway spruce and beech rank lowest in experimentally measured bark insulation \citep{mutterer_bark_2026}. Larch, Scots and black pine and the deciduous oaks are the thick-barked exceptions, whereas Swiss stone pine and mountain pine of the subalpine inner Alps are among the most fire-sensitive species \citep{dupire_vulnerability_2019}. A fire that a fire-adapted stand would survive can therefore remove an Alpine stand and its protective function together. This narrows the scope for prescribed burning. The usable burn window is likely to be short because it falls in the season of the spring ignition peak \citep{muller_forest_2020}, an escape threatens the assets below, and in Austria lighting fires in forest is prohibited \citep{neumann_fuel_2022}. The single Alpine trial, in fire-tolerant Scots pine, nonetheless found little ecological damage \citep{gamba_prescribed_2025}. In the Pyrenees, the Aran valley has adopted a strategic fire management plan that combines prescribed burns of tens of hectares with managed unplanned ignitions over hundreds of hectares \citep{oliveres_path_2025}. In the Alps, prescribed fire is therefore an open empirical question for tolerant stands rather than a standard tool.

Fire intolerance would matter less if recovery were quick, but in the Alps recovery is slow, depends on elevation and is sometimes absent. Protection against rockfall is lowest a decade or more after the fire, and landslide disposition can stay elevated for at least forty years \citep{maringer_temporal_2016, gehring_shallow_2019}. This is far beyond the three to four years typical of the general geomorphological literature \citep{santi_wildfire_2022}. Submediterranean forest stays open for decades after disturbance \citep{stritih_alternative_2023}, post-fire establishment at low and mid elevations is drought-limited \citep{moser_potential_2010}, and at the upper margin succession has remained arrested for eighty years \citep{degiuli_arrested_2024}. Slow recovery is also what places temperate conifer forest high in the European ranking of ecological vulnerability \citep{arrogantefunes_mapping_2024}.

The same features change how fuel can be managed at landscape scale. A fuelbreak works mainly as a platform for ground crews \citep{ortega_modeling_2024}, and in steep terrain crews often cannot reach or hold it. Linear and shaded fuelbreaks are recommended for mountain terrain, on ridgelines or valley bottoms where slopes are very steep, but small private parcels make continuous networks hard to build \citep{ascoli_firebreak_2018, gorrizmifsud_current_2026}, and in difficult terrain timber revenue often fails to cover harvesting costs \citep{dinca_timber_2025}. Area-based treatment would therefore carry most of the load, relying on grazing of abandoned slopes \citep{rouetleduc_effects_2021}, whose effect sizes come from outside the Alps and reverse under overgrazing \citep{xu_how_2026}, and on thinning that reduces fuel and may also ease drought stress \citep{giuggiola_droughtinduced_2010}.

Post-fire restoration is the least developed part of the fire management cycle across Europe \citep{vieira_postfire_2026}, and the Alpine case adds to the difficulty. Restoration there must re-establish protection against gravitational hazards within a long recovery window, and the restoration techniques most often applied after large fires have been judged poorly suited to the ecological context and to global change \citep{mantero_short-term_2023}. The salvage consensus developed after windthrow does not transfer either. A windthrown stand leaves lying stems that continue to stop rockfall, whereas burnt logs and snags may not provide the same barrier effect \citep{costa_biological_2021}, and the consensus is itself contested \citep{seben_impact_2024}. Where fire is assumed to destroy the vegetation, simulated rockfall run-out lengthens, as shown for burnt Mediterranean slopes in Greece \citep{asteriou_assessing_2025}.
% =======================================================
\section{Risk, Exposure and Vulnerability}
\label{sec:risk}

\subsection{From hazard to risk: the missing terms}

A single-hazard wildfire assessment asks where fire is likely and how large it might be. In the Alps this question stops short of what matters. Burnt area is small and only weakly related to consequence, and the exposed asset usually lies outside the burn perimeter but beneath it. Susceptibility, correctly defined, is the static probability that a place experiences fire given the intrinsic characteristics of the territory \citep{trucchia_defining_2022}. It is a legitimate quantity, but in mountains the gap between susceptibility and risk is unusually wide.

Risk combines hazard, exposure and vulnerability, and Alpine fire science has developed only the first. Hazard, the probability and intensity of fire, is reviewed in Section~\ref{sec:hazard}. Exposure is what lies where fire and its consequences reach. Vulnerability is the susceptibility of those elements to loss once exposed. The natural hazards literature defines vulnerability by the conditions that make an exposed asset susceptible to damage \citep{papathomakohle_towards_2025}, and climate-risk assessments index hazard, exposure and vulnerability with separate indicators \citep{pede_mountain_2022}. Neither has been assembled for the Alpine arc, so Alpine assessments end at a hazard surface that is then read as if it were a risk map.

Both missing terms have been developed in a separate literature, with its own instruments and data requirements. Where damage records are absent, indicator sets are weighted by expert judgement \citep{papathomakohle_towards_2025}. Where such records exist, indices are calibrated on surveyed damage \citep{papathomakohle_wildfire_2022}. For the ecological component, continental mapping integrates lost ecological value over the recovery period \citep{arrogantefunes_mapping_2024}. None of these has yet been joined to an Alpine hazard surface. This is the gap that most directly separates the current Alpine literature from a usable risk assessment. Closing it means combining layers that already exist rather than improving any one of them (Section~\ref{sec:governance}).

\subsection{Exposure: habitats, assets, people and seasonality}

Exposure in mountains is neither static nor confined to the burn perimeter. Its ecological component is the one part already measured at continental scale, and the measured share is large. Of the 1{,}079{,}538\,ha burnt in the European Union in 2025, the largest total in the European record, 424{,}023\,ha lay inside Natura 2000 sites. That is 39\% of the total and the highest share ever recorded, three times the 2024 figure and twice that of 2023. Spain, Portugal and Romania accounted for two thirds of it \citep{jrc_advance_2026}. European burning therefore does not touch the protected-area network only incidentally: roughly two of every five hectares burnt lie inside it.

The Alpine part of the same accounting is informative for a different reason. The four fires mapped in Austria in 2025 totalled 434\,ha. Of this, 420.6\,ha lay inside Natura 2000 sites, and 334\,ha was agricultural land rather than forest \citep{jrc_advance_2026}. Where mapped Alpine fire occurs, it thus coincides almost exactly with the protected-area network. It also falls on abandoned or semi-natural agricultural ground rather than on closed forest, where the fuel accumulation described in Section~\ref{sec:forest} is concentrated. But these four fires are a small fraction of the Austrian annual fire count (Section~\ref{sec:data}). The coincidence is therefore established on the detected tail rather than on the regime. The ecological exposure of the Alpine arc as a whole remains unmeasured, for the same reason that its fire regime does.

Since much of what is exposed lies beyond the burnt area, we argue that exposure analysis must follow the cascade rather than the fire. The vulnerability literature does this better than the hazard literature. A participative multi-criteria approach in south-eastern France distinguished exposure, sensitivity and adaptive capacity across populations, ecosystems and infrastructure \citep{riviere_mapping_2023}. It identified secondary vulnerability hotspots, as well as locations not yet affected by hazard where prevention was nonetheless warranted.

Exposure also changes over time. Human presence in Alpine forests is tourism-driven and seasonal. For one Piedmont zone, projections to 2050 combine shorter high-altitude snow cover with an opportunity to attract summer visitors escaping lowland heat \citep{pede_mountain_2022}. The same assessment expects local fire danger to rise mainly in spring, and recorded burnt area there peaks in January, February and November \citep{pede_mountain_2022}, so how far visitor presence and the fire season will overlap remains open. Across five European pilot sites, site-level ignition models found distance to the wildland--urban interface to be the most relevant human predictor \citep{gelabert_assessing_2025}. The same models, however, treat their human covariates as static, and their authors recommend that they vary with commuting and seasonal recreation. Exposure estimates also depend on institutions. Transboundary simulation in the Saxon--Bohemian borderland found modelled burn probability and flame length at settlements differing across a national border that fire does not recognise \citep{avouris_exposure_2025}.

Compound structure enters the exposure picture at several points, and conflating them has caused confusion. Preconditions compound, since European fire activity typically follows moderate but coincident hydro-meteorological and land-surface anomalies rather than single extremes \citep{miller_compounding_2025}. Disturbances compound, since drought and fresh windthrow deadwood favour bark beetle outbreaks in Alpine forests \citep{bebi_changes_2017}, and the resulting dead material loads the fuel bed before any fire arrives. One Alpine hazard projection acknowledges this pathway but does not represent it \citep{gerberding_mapping_2025}. Consequences compound as well, in the post-fire cascade of Section~\ref{sec:water}. A research roadmap now treats post-fire and frozen-soil dynamics together, and calls for multi-institutional monitoring networks and for governance linking geoscientists, engineers and policymakers \citep{baselt_compound_2026}.

The exposure of linear infrastructure to this compound cascade has been quantified, but not in Europe. Across California, the likelihood of exceeding the historic mean coincidence of intense rainfall over burnt ground rises from 38\% to 68\% and 73\% under moderate and high emissions pathways. For the upper decile of events, it rises from 34\% to 90\% and 91\%. The likelihood that natural gas transmission in California faces this compound hazard is thus projected nearly to double \citep{moftakhari_increasing_2019}. The method transfers even though the coefficients do not. An Alpine equivalent, covering roads, rail and pipelines beneath burnt slopes in confined valleys, has not been attempted.

\subsection{Vulnerability of the built environment}

For the Alps, building vulnerability has so far been constructed from expert judgement, and the exercise shows what is missing. An Austrian index scores buildings in the wildland--urban interface on twenty-one indicators, grouped as building environment, building exterior and building characteristics \citep{papathomakohle_towards_2025}. The indicators were weighted by pairwise comparison across a national and an international expert panel. The index was then applied to twenty-three buildings in a Carinthian settlement on the southern slope of the Dobratsch. Both panels ranked protective strips or firebreaks first among the six environment indicators (weights of 0.286 and 0.198). They ranked the distance of tree crowns from the building first among the seven exterior indicators (0.258 and 0.212), and roof material first among the eight building indicators (0.203 and 0.164). The weights rest on expert judgement because no Austrian wildfire damage record exists from which to derive them.

Where a damage record does exist, the weights differ, and the difference is informative. An index was calibrated on 423 buildings surveyed after a wildland--urban interface fire in Attica that killed 102 people and destroyed 1{,}200 buildings \citep{papathomakohle_wildfire_2022}. It retained eight of thirteen candidate indicators and ranked roof material first, structural type second and terrain slope third, whereas both Austrian panels gave terrain slope a low weight. The two indices nonetheless agree on roof material, and the physical reason is known. Structures ignite by wind-blown embers, radiant heat or direct flame contact, and most buildings are lost to embers and surface fire rather than to a flame front reaching the wall. Roofing and other building materials, and the vents, openings and eaves through which embers enter or lodge, therefore matter alongside proximity to vegetation \citep{naser_vulnerability_2025}. Roof material is stable across the two settings, and it is the one indicator an Alpine programme could adopt without local calibration. Slope, which the Alps have in abundance, is where the transfer fails. Preferring either the judgement-based or the damage-based weighting does not resolve this. The two disagree on the variable that most distinguishes the Alpine setting from the Mediterranean one in which the damage record was collected.

Where fire does reach structures directly, damage is shaped in part by the same wind interaction that terrain amplifies. A review of fire--wind simulation covers the fanning effect of wind, the combination of wind direction with slope angle, and the effects of fire--wind interaction on buildings as fire passes from wildland into the wildland--urban interface \citep{ghodrat_existing_2021}. The same review notes that disturbance interactions affecting structures have been comparatively neglected. Alpine settlements sit in valleys where slope and channelled wind routinely combine.

\subsection{What differs in mountains}

Outside mountains, wildland--urban interface risk assessment rests on a simple geometry. Exposure is the set of structures inside or next to the potential perimeter. Vulnerability is a property of the building fabric and its immediate surroundings, which is what damage-calibrated indices and ignition-pathway analysis describe \citep{papathomakohle_wildfire_2022, naser_vulnerability_2025}. Loss is generally assumed to scale with burnt area, because what burns lies inside the area that burns. In the Alps this geometry fails at each step, beginning with the location of the asset.

In mountains, loss is displaced from the fire in both space and time. Fire removes a protective function on the slope, and the damage falls on the settlement or transport corridor below, often months later and through a different process such as debris flow \citep{tiranti_wildfires_2021, laporta_postwildfire_2025}. A susceptibility map shows where fire is likely, not where this loss occurs \citep{trucchia_defining_2022}, and burnt area predicts loss poorly for the same reason. Methods that reach beyond the burnt area do exist. A participatory assessment in south-eastern France mapped localised secondary vulnerability hotspots and vulnerable places not yet confronted with fire \citep{riviere_mapping_2023}. In California, the likelihood of wildfire followed by intense rainfall has been quantified alongside the exposure of gas pipelines to each hazard, although nothing comparable has been attempted for Alpine valleys \citep{moftakhari_increasing_2019}.

Exposure also shifts with the season in ways that lowland registers do not capture. Visitor numbers in Alpine forests follow the tourist calendar, which projections for one Piedmont zone suggest will tilt towards summer as snow cover shortens and visitors seek relief from lowland heat \citep{pede_mountain_2022}. In Mediterranean destinations, tourism raises ignition risk through crowding and careless behaviour \citep{abrham_wildfires_2024}. Even a recent ignition model covering five European pilot sites uses static human covariates \citep{gelabert_assessing_2025}. A static population or building layer therefore misstates exposure, most of all where the visitor season and the fire season coincide.

Once the asset is located, vulnerability weights still do not transfer on the variable that most distinguishes the setting. The Austrian expert-judgement index and the Attica damage-calibrated index agree on roof material, for which there is a physical mechanism \citep{papathomakohle_towards_2025, papathomakohle_wildfire_2022, naser_vulnerability_2025}. They disagree on terrain slope, the defining Alpine variable. Adopting either weighting wholesale imports an assumption about slope that neither evidence base supports for the other setting.

Beyond buildings, the ecological share of exposure is large and the loss has several components. In Austria, mapped fire fell almost entirely inside the Natura 2000 network, although this rests on four detected fires \citep{jrc_advance_2026}. The loss includes a protective-function component \citep{moos_climate_2021} and an ecological-value component integrated over a long recovery \citep{arrogantefunes_mapping_2024}. It also includes a biodiversity component whose sign differs between guilds \citep{rey_effects_2019, degiuli_arrested_2024}, and a faunal component that has never been measured for the Alps (Section~\ref{sec:forest}). A single-hazard assessment that ends at a burn probability leaves all of these unquantified \citep{trucchia_defining_2022}.

% =======================================================
\section{Governance, Policy and Economics}
\label{sec:governance}

\subsection{Suppression, preparedness and the practitioner constraint}

Alpine suppression works under several constraints at once. Terrain makes fire behaviour harder to anticipate, a mechanism established mainly outside Europe \citep{sharples_overview_2009} and documented in the Alps for one steep-valley fire \citep{melzner_brief_2019}, and steep, remote ground is hard to reach \citep{muller_forest_2020}. Suppression also depends on shared aerial resources. Drought and vapour pressure deficit are the dominant drivers of European wildfire activity, and fire weather is shaped by large-scale weather patterns \citep{miller_compounding_2025}, so we expect demand for these resources to peak across the arc at the same moment.

In practice, these constraints produce a suppression system organised around water logistics and vertical distance, although ground and surface fires are still attacked directly at the front or flanks \citep{muller_forest_2020}. Where vehicles can reach, crews attack from tank tenders using ponds, hydrants and long hose lines, with additional pumps inserted to overcome large differences in height. Where vehicles cannot reach, initial attack on ground fire uses rakes and beaters, and subsurface hot spots are reached with backpack pumps and high-pressure lances. Vehicles are chosen to suit the road rather than the fire. Italian forest-fire units, for example, use all-wheel-drive pickups sized for narrow Alpine forest roads instead of the heavier tankers used on open ground. In the southern Alps, sprinkler lines wet vegetation along fuel breaks, particularly at night. Tactical or counter-firing during an incident, as distinct from prescribed burning for prevention, is practised mainly in France and Italy and is uncommon elsewhere in the range.

Aerial support is a structural part of this system rather than an add-on, and it serves a different function than in the Mediterranean. Instead of laying retardant along a long front, helicopters in Austria and Bavaria increasingly ferry water to mobile basins near the fire, from which ground crews lay hose. Every Alpine country maintains units trained for bucket work and winch insertion. Road density is not the main limitation, since forest road density in the Alpine region is high compared with Mediterranean or Scandinavian Europe, although road practicability and barriers to vehicles repeatedly cause problems \citep{muller_forest_2020}. The constraint is vertical rather than horizontal, which corrects an assumption that is easily imported from Mediterranean practice. What limits access is the slope, exposure and vertical separation between the road and the fire. A terrain difficulty index calibrated on Californian initial attack approximates this through slope, fuel penetrability and travel time along the road network \citep{cardil_assessing_2025}.

Alpine crews also face different hazards. Rockfall and falls are ordinary features of the working environment, and burning material rolling downslope ignites new fires below crews working above it. In sectors of the range where the two world wars were fought, buried munition remains a recognised hazard. Where terrain is steep enough, mountain rescue teams join the action force to secure the firefighters themselves \citep{muller_forest_2020}. Confined topography adds a behavioural hazard, since canyons can produce sharp increases in fire intensity without any increase in wind \citep{muller_forest_2020, hayajneh_wind_2026}. A suppression doctrine written for terrain where the main danger is the fire front does not describe this working environment.

The Alpine system also differs in organisation, and this is the part least represented in the modelling literature. In Austria and Switzerland, dense networks of volunteer brigades and a population that reliably reports smoke by telephone mean that most fires are attacked soon after ignition \citep{muller_forest_2020}. We take this to be part of the explanation for a fire-size distribution dominated by events under a hectare, although the inference has not been tested. The same fire weather in a range without that response density would plausibly produce a different distribution, and no Alpine hazard model contains a term for it. This capability also carries a demographic exposure of its own, because it rests on voluntary participation in valley communities that are ageing and shrinking (Section~\ref{sec:economics}).

A quantitative framework for the accessibility constraint now exists, although it has not yet been applied in the Alps. It combines a terrain difficulty index with a fire behaviour index to assess initial attack \citep{cardil_assessing_2025}. The terrain index is built from accessibility, penetrability, line-construction difficulty, mean slope and road proximity. In its highest class, the odds of successful initial attack are 90\% lower than in the lowest, and mean burnt area rises from 18.5 to 1731\,ha across that range. Response time interacts with terrain difficulty to reduce containment probability further. Because the framework has been validated outside its region of origin, the Alpine gap is one of application rather than of method, and it can be closed. Where suppression cannot be fast everywhere, it must be fast in the right places. These places can be identified in advance from terrain and fuel connectivity, which is the operational case for targeted pre-suppression infrastructure such as water points for aerial firefighting \citep{conedera_linking_2024}.

How suppression capability and fuel treatment interact matters in a range whose danger distribution is shifting. Fuel treatment reduces the severity that suppression crews have to confront. Its benefit for tree mortality, however, is greatest under intermediate fire weather and declines under the most severe conditions \citep{brodie_forest_2024}. Initial-attack success, by contrast, falls steeply with terrain difficulty, an effect separate from that of fire behaviour, in Californian data \citep{cardil_assessing_2025}. In accessible Alpine terrain under moderate danger, the two are substitutes, and the cheaper one can be chosen. On steep slopes with long, continuous forested upslope runs, where most large Swiss fires have spread rapidly and proved hard to suppress \citep{conedera_linking_2024}, suppression is weak and treatment is the only lever that acts before the fire. These are also the sites where mechanised treatment costs most \citep{dinca_timber_2025}, which is why the pastoral and silvicultural instruments of Section~\ref{sec:forest} matter most there.

Alpine fire management currently spends several times more on response than an integrated prevention framework would cost. Section~\ref{sec:economics} gives the ratio, together with econometric evidence from Italian regions that points the same way. What no Alpine study provides is the quantity that would settle the allocation: the marginal effectiveness of a euro spent on suppression capability against a euro spent on fuel treatment, conditioned on terrain.

The strategic question is how much of the projected increase management can absorb. The continent-scale projections introduced in Section~\ref{sec:intro} set limits on the answer. Sustained investment in capacity offsets most of the climate-driven increase in burnt area but may not offset all of it under strong warming \citep{billing_future_2026}. Even with rising capacity, one of the two models still projects increases of 200 to 300\% along the Alpine foothills under high emissions \citep{billing_future_2026}. For the Alps, investment in capacity is therefore the largest single lever within the region's control, not a marginal adjustment. It is unlikely, however, to be sufficient on its own in the foothills.

\subsection{Prevention: placement, early warning and integrated fire management}

Section~\ref{sec:forest} set out the instruments of Alpine prevention; the governance question is whether they are used where they are needed. Fire-smart management is already the operational focus in Mediterranean regions. In the currently less affected Alps, the case for building a pre-fire capability rests on anticipation rather than experience \citep{gerberding_mapping_2025}. Allocation is therefore a political as well as a technical problem.

The Alpine evidence on allocation is discouraging. A mapping of 97 nature-based solutions for climate adaptation across the European Alps found that most addressed drought or soil instability and aimed to deliver several benefits at once \citep{dubo_nature-based_2023}. Yet they were generally not located within hotspots of the hazards they claimed to address, nor within hotspots of the ecosystem contributions they aimed to secure, nor where the two coincided. Measures are being placed where they are institutionally feasible rather than where they are physically needed. Closing that gap requires local climate projections and quantified ecosystem contributions to be integrated into planning, through a macro-regional strategy combined with local stakeholder engagement.

A low-cost correction is available. Large-fire likelihood in mountain terrain can be mapped from slope length and fuel connectivity using only a digital elevation model and a forest map \citep{conedera_linking_2024}. Protective function can be valued per hectare by replacement cost \citep{accastello_asforesee_2019}. The two layers needed to rank candidate treatment sites by expected avoided loss therefore already exist for most of the Alpine arc. What is missing is the institutional mandate to combine them, because the hazard layer sits with fire and forest services and the value layer with the authorities responsible for the assets below.

A transferable method for this prioritisation already exists. A multi-criteria analysis of some 238{,}000 forest polygons across Catalonia ranked management zones on fuel hazard, simulated fire behaviour, exposure of settlements and infrastructure, and landscape continuity \citep{krsnik_spatial_2024}. It also ranked them on accessibility, the criterion that matters most in mountains, expressed through terrain aspect, road density and distance to water points and strategic management areas. The method is directly portable to Alpine conditions, where accessibility is the binding constraint on suppression. Its sensitivity is also instructive. Expert weighting designated 20.69\% of the territory as high priority, while equal weighting designated 7.98\%, so the participatory step shapes the result as much as the data do.

In our view, the most tractable near-term gain is not a new Alpine model but the extension of early warning systems that already exist. The Alps have dense hydro-meteorological monitoring, mature torrent and avalanche warning services and long institutional experience of gravitational hazards. What those systems lack is a fire term. The concrete proposal in the Alpine literature is to adjust catchment runoff parameters automatically for burnt catchments, so that a debris-flow early warning system lowers its rainfall thresholds without waiting for a field survey \citep{tiranti_wildfires_2021}. The thematic layers required (catchment and fan delimitation, geology, bedrock outcrop, land use and satellite-derived fire perimeters) are already available in geographic information systems. The Bussoleno event suggests why this is needed, because its radar-estimated triggering intensity appears to have fallen below the threshold that a system designed for unburnt basins would have been watching for \citep{laporta_postwildfire_2025}.

The modelling evidence also bears on how such a system should be designed. Its outputs are most useful as decision thresholds rather than probability surfaces, with the asymmetry between missed detections and false alarms built into the specification rather than considered afterwards. Alpine ignition modelling has already been built this way, deriving dynamic thresholds from combinations of true-positive and false-alarm rates \citep{moreno_spacetime_2025}. Integrated danger assessment has likewise reported false-alert performance as a main result \citep{muller_towards_2020}. This design choice matters where a missed detection means an unwarned slope above a village.

For the same reason, integrated fire management is best specified around indicators that different institutions can each observe and act on. Alpine post-fire inventory work proposes exactly this \citep{melzner_wildfire-induced_2025}. It identifies indicators for burn severity and for subsequent geohydrological processes, sets out systematic documentation procedures, and is framed for use by policymakers within an integrated fire management strategy rather than by researchers. The costed Alpine framework already allocates between a tenth and a fifth of a proposed annual budget to post-fire management and a further share to knowledge transfer \citep{muller_forest_2020}. The compound-hazard roadmap outlines the monitoring options from which networks supplying the indicators could be built \citep{baselt_compound_2026}. Each component has been proposed separately; what is missing is their assembly.

\subsection{Fragmentation, harmonisation and cross-border coordination}

The Alpine wildfire cascade crosses an institutional boundary at every link. Prevention sits with forest administrations, and suppression with fire services and civil protection. Post-fire rockfall and torrent control sit with geological surveys and infrastructure operators, and the land-use decisions that govern fuel continuity sit with agricultural and planning authorities. The sequence described in Section~\ref{sec:water} passes through all four.

The drivers of fire also differ between countries, which is one reason a single Alpine prescription is unlikely to fit. Rural abandonment dominates in France and Italy, and tourist activity in Austria, Germany and Switzerland, while negligence is the leading proximate cause throughout \citep{muller_forest_2020}. Prevention that works through agricultural policy in the western and southern Alps must work through visitor management in the north and east.

Fragmentation extends to basic facts. A survey of three stakeholder groups across seven Alpine countries found significant differences in perceptions of fire numbers, burnt area and future trend \citep{hoang_stakeholders_2026}. These differences persisted even though the groups agreed that precipitation, temperature and wind are the dominant natural drivers, and that recreation, rural abandonment and owner awareness are the dominant socio-economic ones. A practitioner community that does not share a factual picture of the hazard is poorly prepared for its growth.

How this plays out where citizens have formal standing has been surveyed in one Alpine setting. In the Locarnese region of Ticino, residents valued local amenity highly and strongly supported firefighting services \citep{bardsley_wildfire_2021}. They also judged that the main driver of risk at the forest and settlement interface is the current level of urban expansion rather than the hazard itself. The planning system is built on deliberative instruments, yet younger, locally rooted and less formally educated respondents were the least satisfied with their influence over it. The availability of popular referenda was judged to foster a disproportionate sense of entitlement over local decisions. In a range where land-use decisions govern fuel continuity, public support for suppression capacity is therefore not evidence of support for the land-use restraint that would reduce exposure. Formal participation rights, moreover, do not distribute influence evenly among those who carry the risk.

Practitioners within a single Alpine country give the same diagnosis when asked directly. A structured survey of 79 wildfire experts across seven Italian regions, three of them Alpine, found the institutional structure highly fragmented, with responsibility increasingly spread across separate organisations \citep{kirschner_governance_2024}. Emergency response persisted as the default strategy despite formal commitments to precaution. Respondents reported that local stakeholder participation had been displaced by a shift towards extra-local actors, and judged state-level resources disproportionate to regional capacity. Prevention dominated both halves of the assessment: 46\% of the coded shortcomings concerned prevention deficits, and 58\% of the identified opportunities concerned enhanced prevention. That one theme dominates both what is missing and what is possible is the clearest available evidence that the constraint is institutional rather than technical.

Fragmentation is also visible one level up, in the European policy framework itself. A review of forty EU sectoral policies and initiatives across seven clusters found integrated fire management, or a homologous term, named seven times in the whole corpus, and only once, in passing, in a law or policy text \citep{plana_framing_2024}. The term extreme wildfire event appeared four times in total, and law and policy texts treated fire only as a threat, without its beneficial role. Nature-based solutions appeared only marginally in the forestry and disaster-risk clusters, despite their prominence in climate and biodiversity policy. A comparison of European and United States arrangements for critical infrastructure reached a parallel conclusion, identifying insufficient integration of scientific evidence into policy and recommending standardised cross-border risk assessment \citep{kalapodis_towards_2025}. The Alpine coordination problem is therefore not only that eight countries fail to align with one another. The supranational framework they would align to barely names the practice.

Harmonisation has been attempted, but it has been driven by projects. The first harmonised Alpine fire database and regime characterisation came from a transnational project \citep{valese_wildfires_2011}. Its data layers were later reused in national danger assessment \citep{muller_towards_2020}. A later EUSALP project, drawing on an expert panel, a multilingual survey and a workshop, produced an Alpine-scale white paper that sets out prevention, suppression, post-fire and knowledge-transfer priorities with costed estimates \citep{muller_forest_2020}. What does not yet exist at Alpine scale is standing infrastructure, such as the monitoring networks and post-fire documentation procedures discussed above. Compound hazard research requires those networks to combine remote sensing, in situ instrumentation and geophysical methods \citep{baselt_compound_2026}. The documentation procedures are what would make indicator-based integrated fire management operable \citep{melzner_wildfire-induced_2025}. That fire crosses national borders has been shown directly at a central European transboundary interface (Section~\ref{sec:risk}) \citep{avouris_exposure_2025}.

Standing coordination across the jurisdictions of a mountain range does exist elsewhere, and the one documented case shows both what it can achieve and where it stops. The protected areas of the Australian Alps fall under the legislation of three states and territories. Since 1986, a memorandum of understanding between four governments has supported a cooperative management programme coordinated by the Australian Alps Liaison Committee \citep{lockwood_biodiversity_2014}. The programme treats the range as one biogeographic region for fire management, water and catchments, invasive species and Aboriginal engagement, and it has survived political and funding cycles and major fire events. Coordination did not insulate fire policy from politics, however. After the 2009 Victorian fires, a hazard-reduction burning target of 5\% of public land per year became state policy, although such burning is not undertaken in the treeless alpine landscape. In 2010, cattle were returned to the Victorian Alps for a trial of their contribution to reducing fire risk, despite published evidence suggesting that alpine grazing has no significant effect on it. The federal government later removed them on the strength of the range's national heritage listing. The case is single and imported, but it suggests that a standing body can align agencies without settling the land-use disputes that fuel management raises.

\subsection{The economics of Alpine wildfire}
\label{sec:economics}

The economics of Alpine wildfire run against intuition. What the forest protects is worth more than what it produces, and prevention costs little beside either.

Protective function dominates the value at risk. Replacement-cost valuation of a ten-hectare Alpine Scots pine protection stand returns €30{,}440 per hectare, or €950 per hectare per year \citep{accastello_asforesee_2019}. Avoided-risk valuation of a Valais protection forest places annual rockfall risk at 19{,}100\,CHF with the forest and 95{,}500\,CHF without it. Under moderate and high emissions scenarios, the risk rises to 59{,}250 and 70{,}100\,CHF per year within 150 years, as forest composition shifts (Section~\ref{sec:forest}) \citep{moos_climate_2021}. Neither figure is a timber value. In terrain where extraction is technically demanding and economically marginal \citep{dinca_timber_2025}, the timber value is frequently the smaller number.

On the expenditure side, response dominates. Direct firefighting and post-fire management across the Alpine arc are estimated at around €75 million a year. A proposed integrated fire management framework is costed at roughly €10 million a year, of which €5--7 million would go to prevention, €1.5--2.5 million to suppression capability, €1--2 million to post-fire management and €0.5 million to knowledge transfer \citep{muller_forest_2020}. Whatever the precision of these estimates, the ratio frames prevention as inexpensive relative to the response it would displace. An econometric analysis of Italian regions, which projects more frequent fires under climate change, led its authors to the same recommendation from a different direction, namely that investing in prevention would be more efficient than suppression \citep{michetti_forest_2019}.

The same analysis gives the clearest quantitative statement of how socio-economic and climatic drivers combine in a European mountain setting. Across nineteen Italian regions, a ten-millimetre increase in monthly precipitation reduced fire frequency by 8.9\% in the same month, with effects persisting longer in the north than in the south. A one-degree temperature increase raised frequency by up to 15\% in the south and by between 0.9\% and 5.7\% elsewhere. Socio-economic terms were of comparable size. A one-percentage-point rise in tertiary education reduced fire events by 11\% in the north. In the south, a single additional extortion per hundred thousand inhabitants raised fire frequency by 4.6\% and burnt area by 7\%. Projected to 2016--2035, the northern cluster gains 11--12\% in annual fire events and 23--27\% in annual burnt area \citep{michetti_forest_2019}. Education and institutional quality are thus fire management variables, although they are not usually treated as such.

Delivery is a further problem. Fragmented Alpine forest ownership shares it with the one case where it has been measured, in Catalonia. There, wildfire prevention subsidies funded by a vehicle carbon tax, offering €5{,}000 to €200{,}000 per proposal, reached only 3.87\% of the 151{,}196\,ha designated as strategic \citep{gorrizmifsud_current_2026}. The average minimum distance between subsidised parcels was 287\,m, so the scheme did not achieve the spatial continuity it was designed to create. The barriers identified were non-adjacent parcels, uneven social capital between owner associations, and reimbursement rules that required owners to advance costs. Fuel treatment that must be contiguous to work is thus being funded through instruments that allocate to whoever applies.

The economics literature also identifies structural incentive problems that Alpine studies have not yet examined. A private parcel owner deciding on mitigation weighs private costs and benefits \citep{bayham_economics_2022}. The owner ignores spatial spillover externalities, such as fire spreading across parcel boundaries, and conventional externalities such as smoke. Liability rules and insurance markets then shape mitigation incentives in ways that vary by jurisdiction. In an Alpine context of small private parcels and eight national legal systems, these are first-order determinants of whether fuel treatment happens at all. Catalan subsidies designed to capture this spillover illustrate how hard it is to secure adjacent treatment even when coordination is paid for \citep{gorrizmifsud_current_2026}.

Pastoral systems sit on both sides of the Alpine ledger, and neither side has been priced. Grazing is the fuel-treatment instrument with the broadest policy leverage (Section~\ref{sec:forest}). It is cost-effective relative to mechanical clearance because it relies on natural animal behaviour \citep{xu_how_2026}, and, in our reading, because the work is done by animals kept for other reasons. Its marginal cost to a fire budget is therefore largely the cost of sustaining a pastoral economy that already exists and is contracting. On the other side, alpine pasture, its access infrastructure and the stock on it are exposed assets, and transhumant animals are at their highest elevation during the summer fire season. We found no Alpine study that values either side. There is no estimate of the fuel-management service that extensive grazing supplies without payment, and no record of livestock or pasture loss to fire against which to set it. In the one European landscape scenario analysis available, from an Iberian transboundary reserve, reversing agricultural abandonment supplied almost all of the suppression-cost saving \citep{lecinadiaz_incorporating_2023}. This quantitative gap therefore bears directly on how agri-environmental payment is justified.

What the economics cannot yet resolve is distribution. Alpine valley communities are small, ageing and economically specialised. They bear the costs of both the hazard and its prevention, while a large share of the protective benefit accrues to through-traffic, downstream water users and regions outside the mountains. The services at stake propagate across scales, because the vegetation changes through which fire acts on them operate locally but register regionally \citep{vukomanovic_systematic_2019}. Risk assessment frameworks that make adaptive capacity an explicit component of vulnerability come closest to representing this asymmetry \citep{riviere_mapping_2023, lecinadiaz_assessing_2021}. Tourism also sits on both sides of the ledger, since in Mediterranean evidence fires depress visits while visitors raise fire risk \citep{abrham_wildfires_2024}. In one Italian Alpine area, warming is projected to draw more summer tourists to high-altitude municipalities while fire hazard moves towards them (Section~\ref{sec:risk}) \citep{pede_mountain_2022}. We found no Alpine study that quantifies who bears the costs and who receives the benefits.

\subsection{What differs in mountains}

Most fire policy assumes a suppression model, in which the dominant institutional response to fire is to fight it \citep{moritz_learning_2014}. Italian practitioners describe exactly this bias towards emergency response \citep{kirschner_governance_2024}. The model takes for granted that capability scales with the fire, that aerial and ground resources reach the line, and that prevention subsidies buy contiguous fuel treatment. It also assumes that timber revenue helps fund forest management and that one national administration covers the fire environment. Each of these assumptions is weaker in the Alps, and the weaknesses add up rather than cancel out.

The difference starts with what counts as an extreme event, and so with when a fire service is expected to escalate. Outside mountains, an extreme fire is generally understood as one that burns thousands of hectares and stops when fuel runs out or the weather changes rather than because crews acted, that is, a fire beyond suppression capacity. In the Alps, a fire of a few hectares is extreme if its severity removes the protective function above a settlement \citep{muller_forest_2020}. Severity there is defined by what lies beneath the fire, not by its size or by the limits of the suppression system. An operational doctrine calibrated on the lowland definition will therefore not escalate on the events that matter in mountains.

Once a fire is burning, Alpine suppression differs in kind and not only in difficulty. It is organised around water logistics and vertical lift, aerial support is a structural part of it, and crews face rockfall, falls, burning material rolling downslope and, in places, buried munition. Dense volunteer brigade networks reach most ignitions quickly, a capability that fire-prone lowland regions often lack and that no hazard or scenario model represents. Road density is high by European standards and is not the main limit on access, although road practicability and barriers to vehicles repeatedly cause problems \citep{muller_forest_2020}. The main limit is slope and the vertical separation between road and fire, which terrain difficulty indices, so far calibrated on Californian fires, approximate through slope, fuel penetrability and road accessibility \citep{cardil_assessing_2025}.

This weakness is located where hazard concentrates. In Californian data, containment probability falls steeply as terrain difficulty rises \citep{cardil_assessing_2025}, and in Switzerland large-fire statistics are highest in the regions with the longest continuous forested upslope runs \citep{conedera_linking_2024}. Andalusian data show that a fuel break works as a platform for ground crews rather than as a barrier (Section~\ref{sec:fuelbreaks}) \citep{ortega_modeling_2024}, so it fails where crews cannot reach. The lever on which most fire budgets rely therefore has the least purchase on steep ground.

Above the level of the incident, the hazard domain and the jurisdiction do not coincide. Eight national systems, with different legal frameworks, reporting conventions and fire services, share a single range across which fire weather, fuel continuity and smoke move freely, and fire has been shown to cross borders in central Europe \citep{avouris_exposure_2025}. Surveyed authorities, scientists and fire-service representatives from seven Alpine countries differ significantly in how they perceive fire numbers, burnt area and future trends \citep{hoang_stakeholders_2026}. The European framework to which they might align barely names integrated fire management \citep{plana_framing_2024}. The fire record itself was harmonised only recently, through a transnational project \citep{valese_wildfires_2011}, so arc-wide statements rest on a short series. Standing coordination across a mountain range is possible, but the one documented case is imported from the Australian Alps, and there a standing body aligned agencies without resolving land-use disputes over fuel \citep{lockwood_biodiversity_2014}.

Below the national level, forest ownership is fragmented at a scale that defeats instruments requiring contiguity. Owners of small private parcels have little private incentive to provide the spillover benefit of treating their land \citep{bayham_economics_2022}. Even subsidies that require grouped applications can then produce scattered treatment that falls short of the continuity the scheme was designed to create. This shortfall has been documented in Catalonia rather than in the Alps \citep{gorrizmifsud_current_2026}, but Alpine ownership is fragmented in a similar way.

The economics also differ, and with them the question of who should pay. In our reading, the main value at risk is a public protective service rather than private timber. Rockfall protection by forest has been valued in two Alpine case studies \citep{accastello_asforesee_2019, moos_climate_2021}, and timber extraction in steep terrain is often marginal \citep{dinca_timber_2025}. Prevention would cost a fraction of current response spending \citep{muller_forest_2020}. Yet much of the benefit of protection flows out of the valleys, while the costs stay with the communities that live in them (Section~\ref{sec:economics}). Shifting budget from response to prevention would therefore need a distributional accounting, and no Alpine study has yet performed one.
% =======================================================
\section{Discussion}
\label{sec:discussion}

The comparisons that close Sections~\ref{sec:hazard} to~\ref{sec:governance}, collected in Table~\ref{tab:contrast}, show that mountain fire departs from the non-mountain baseline along the whole chain from ignition to loss. Because these differences arise from different causes and do not cluster at one point, a single methodological correction is unlikely to reach them. Nor are they specific to the Alps. Where the comparative ranges of Table~\ref{tab:ranges} provide evidence on the same process, they mostly agree with the Alpine record. The anthropogenic spring season recurs in the Carpathians, the gap between satellite and national fire records recurs in Fennoscandia, and the land-use origin of the fuel bed recurs in the Apennines. A burnt catchment in Colorado also began producing debris flows at a rainfall intensity within a few millimetres per hour of the Alpine value.

\begin{table}[H]
\centering
\small
\setlength{\tabcolsep}{4pt}
\renewcommand{\arraystretch}{1.15}
\caption{The mountain difference by domain. The baseline column states the position that non-mountain, predominantly Mediterranean and lowland, fire science supports; the departure column states what mountain evidence establishes instead, drawing on the Alps as the focal case and on the comparative ranges of Table~\ref{tab:ranges}.}
\label{tab:contrast}
\begin{tabularx}{\textwidth}{@{}p{2.1cm} X X@{}}
\toprule
\textbf{Domain} & \textbf{Non-mountain baseline} & \textbf{Alpine departure} \\
\midrule
Ignition & Predominantly anthropogenic, clustered near settlement; a single summer season & Mixed, with lightning a leading hazard predictor igniting higher and steeper; bimodal, with a spring surface-fire season having no Mediterranean equivalent \\
Spread & Wind-driven, slope entering as a separable correction & Slope-driven and non-linear, with wind and slope interacting; the driving wind field is not the interpolated synoptic field \\
Hazard models & Calibrated on long records of comparable fires; burnt area is the terminal quantity & Transferred indices lose skill, local refitting has a short record, and burnt area is a poor predictor of loss \\
Data and Earth observation & Detection thresholds sit well below the modal fire; severity products index soil condition; detection technology assumes open sightlines & Thresholds sit at or above the modal fire, so the detected population is the tail; severity decouples from soil condition and carries a seasonal-timing bias; camera and drone detection are unresolved in complex relief \\
Climate and weather & A smooth interpolable fire-weather field; one fire season; projection moves one envelope & A discontinuous field, several climatologies within tens of kilometres, a cryospheric season boundary, and terrain-structured lightning \\
Water and cascade & Runoff and sheet erosion rise for one to three years; reservoirs buffer the chemical signal & Debris flow with defined runout onto inhabited fans; the triggering threshold falls below ordinary storm intensity; the window is set by biology and can be effectively permanent \\
Forest and fuel & Fuel accumulates through fire exclusion; species carry fire adaptations; treatment protects the stand & Fuel accumulates through land abandonment; species are largely fire-intolerant; treatment reduces the protective function it is meant to defend \\
Risk and exposure & Exposure lies inside the perimeter; loss scales with burnt area; vulnerability is building fabric & Exposure lies below the perimeter and is tourism-seasonal; loss is decoupled from burnt area; vulnerability weights disagree on slope \\
Governance and economics & Suppression scales with the fire; an extreme event is one beyond suppression capacity; one jurisdiction; timber revenue part-funds management & Suppression collapses with terrain and depends on aerial lift and volunteer brigades; a few hectares is extreme if protective function is lost; eight jurisdictions in one hazard domain; the payer is not the beneficiary \\
\bottomrule
\end{tabularx}
\end{table}

Several entries in the table reverse the direction of a relationship rather than changing its size. Fuel treatment in a protection forest reduces the asset it is meant to defend. Suppression is least effective where hazard is highest. After fire, a catchment responds to ordinary rainfall rather than only to extreme events. The comparative evidence disagrees with the Alpine record in only one place. Elevation predicts fire in opposite directions in the Alps and the Carpathians, which suggests that, across mountains, terrain shapes where fire occurs but not always in the same direction (Section~\ref{sec:hazard}). The differences also reinforce one another. The detection threshold that hides most Alpine fire makes local calibration difficult. That difficulty makes coefficients transferred from lowland settings attractive, and those coefficients carry assumptions that the rest of the table contradicts.

Taken one at a time, the changes documented here would support only a cautious conclusion. Fire danger is rising from a low base, fuel is accumulating slowly, most fires remain small, and documented post-fire debris flows remain rare (Section~\ref{sec:data}). We argue, however, that they should be read together, because they are converging on the same ground. The climatic envelope is expanding fastest, in relative terms, where fire experience is least \citep{gerberding_mapping_2025}, and the climate signal in fire danger emerges earliest where current danger is lowest \citep{miller_climate_2024}. On the same terrain, fuel is becoming continuous as abandoned land regrows where climate allows \citep{garbarino_contrasting_2020}, lengthening the forested upslope runs that Swiss data associate with large fires \citep{conedera_linking_2024}. In one Central Alpine valley, about four fifths of the cells with newly established trees represent infilling below the treeline rather than treeline advance \citep{frei_evidence_2023}. The snow season that has so far limited high-elevation fire is shortening at elevation-dependent rates, as documented mainly for the Tibetan Plateau \citep{rangwala_climate_2012}, and fire activity in European mountain regions has measurably moved upslope \citep{beloiu_climatedriven_2026}. In one Italian Alpine area, the same warming is projected to draw more summer visitors to high-altitude municipalities while fire hazard moves towards them \citep{pede_mountain_2022}.

In a lowland system, a modest increase in fire activity would produce a modest increase in loss. In the Alps, in our reading, the change is qualitative, because fire acts on the protective function of forests standing above fixed, high-value assets. At Bussoleno, fire appears to have lowered the triggering rainfall intensity so that a storm of about five-year return period produced a debris flow \citep{laporta_postwildfire_2025}. Its consequences persist for at least three years in the soil \citep{negri_wildfire_2026} and for at least forty years in shallow-landslide disposition after severe fire \citep{gehring_shallow_2019}. Canopy recovery after disturbance in general takes decades \citep{stritih_alternative_2023}, and at the one subalpine site studied, woody vegetation has not returned some eight decades after fire and may never do so \citep{degiuli_arrested_2024}. Loss is therefore not proportional to burnt area, and what is changing is how often fires reach the conditions under which loss escalates.

This argument does not imply that the geomorphic cascade dominates Alpine fire loss. We have documented the cascade at length because it is the pathway least represented in existing assessment and the one that most clearly separates mountain from lowland fire. It is not the only pathway, nor necessarily the largest in expected loss. Direct destruction of protective structure, ecological change, smoke, damage at the wildland--urban interface, and disruption of tourism and transport all operate without any subsequent rainfall, and several are better evidenced. A regionalisation of fire regimes across the Alpine arc concluded that fire does not currently rival storms and bark beetles for damage and cost in the region \citep{conedera_characterizing_2018}. Nothing in the more recent literature overturns that conclusion for the present. Our argument concerns the direction of change and the disproportionate consequences a fire can have in mountains, not a claim that fire is already the dominant Alpine disturbance.

If this reading is right, the appropriate response is coexistence with fire rather than its exclusion. Fire is unusual among natural hazards because the dominant institutional response is to fight it, whereas for earthquakes, floods and storms the emphasis is on identifying vulnerabilities and adapting to them. A command-and-control approach also neglects the role of fire regimes in sustaining biodiversity and ecosystem services, and the alternative requires place-based approaches linked to land-use planning \citep{moritz_learning_2014}. The Alps are in an unusual position, because they are building their fire management now rather than reforming an entrenched system. They could therefore, in principle, adopt an adaptation-led approach without first passing through a century of suppression.

The implications differ for each community that would have to act on them. For fire science, the Alpine case tests whether the discipline's tools are general or only tuned to the settings where they were built, and the evidence more often supports the second answer. Danger indices require recalibration \citep{de_angelis_modelling_2015, zacharakis_environmental_2023}, and pooled ignition models lose skill at regional level \citep{gelabert_assessing_2025}. In the one Alpine comparison, two spread models that both reproduced the perimeters of three fires disagreed on fireline intensity \citep{malanchini_modeling_2025}, the quantity that in our reading drives the downstream processes. Severity maps failed to identify the most severely burnt soil at the one Alpine fire where this was tested \citep{negri_wildfire_2026}, and continental models are neither calibrated on the small fires that dominate the Alpine record nor evaluated in mountains separately \citep{forrest_understanding_2024}. More constructively, the Alps are a good test bed for tools that rely on well-chosen structure rather than on many fitted parameters, because the data will support little else.

Mountain geomorphology faces the same problem from the other side. Fire is arriving as a new preparatory factor in a field whose frameworks were built around rainfall, snowmelt, freeze--thaw and permafrost degradation. Its Alpine fire literature still consists of first inventories \citep{melzner_wildfire-induced_2025} and individual documented events \citep{tiranti_wildfires_2021, laporta_postwildfire_2025, melzner_brief_2019}. The efficient path is probably to make existing infrastructure fire-aware rather than to build a parallel system. One example is the proposal to adjust curve-number values automatically for burnt catchments in an operating debris-flow warning system \citep{tiranti_wildfires_2021}, and the successful transfer of an empirical post-fire debris-flow likelihood model from the United States to one Greek fire suggests that this is feasible \citep{diakakis_exploring_2023}.

For climate adaptation, the implications are less comfortable. Alpine nature-based solutions are mostly not located where their target hazards concentrate, those addressing drought and floods being the exceptions \citep{dubo_nature-based_2023}, which warns against planning that starts from institutional feasibility rather than from hazard. Several adaptation and mitigation policies, among them carbon sequestration, reduced management intensity and protective forest expansion, also add standing biomass and fuel continuity (Section~\ref{sec:forest}). Expanding forest cover enlarges the area that can burn \citep{gerberding_mapping_2025}. This is not an argument against those policies, but fire hazard should enter their appraisal, and at present it generally does not.

Each of these implications ultimately depends on governance. No single body, whether the fire service, the forest administration or the geological survey, can manage Alpine wildfire risk alone. One workable arrangement divides responsibility over time. Forest and fire authorities would own prevention and suppression, and a post-fire susceptibility phase lasting years would follow, owned jointly with torrent, avalanche and infrastructure authorities. The fire itself triggers the handover between these phases, and at present that handover has no institutional form.

The research needed next follows from the same diagnosis, but not all of it is equally within reach. Some could begin at once, because it applies instruments already validated elsewhere. A terrain-conditioned suppression-difficulty framework could be applied to the Alpine arc immediately, replacing an evidence gap with a measurement \citep{cardil_assessing_2025}. Ranking candidate treatment sites by expected avoided loss requires a large-fire likelihood surface and a per-hectare valuation of protective function. A large-fire likelihood surface exists for Switzerland and can be computed elsewhere from a terrain model and a forest map \citep{conedera_linking_2024}, and a replicable per-hectare valuation of protective function has been demonstrated in one Alpine stand \citep{accastello_asforesee_2019}. What is mainly missing is the mandate to extend and combine them (Sections~\ref{sec:data} and~\ref{sec:governance}). A cascade-explicit assessment, followed from an Alpine fire scenario to the consequences for a named settlement or transport corridor, could also be assembled now. It could build on the proposed template for automatic post-fire threshold modification, derived from two Alpine catchments \citep{tiranti_wildfires_2021}, with the coincidence-likelihood method applied to Californian energy infrastructure supplying the exposure component \citep{moftakhari_increasing_2019}.

Much of the remaining work depends on observations that do not yet exist, and it is slower because it is foundational rather than incremental. The Alps have no long-term monitoring of wildfire and geohydrological processes comparable to that of the western United States \citep{melzner_wildfire-induced_2025}. The instrumentation design proposed for compound mountain hazards offers a reasonable starting specification \citep{baselt_compound_2026}. Without post-fire soil products independent of spectral severity, the post-fire hazard chain continues to rest on a link that failed at the one Alpine fire where it was tested \citep{negri_wildfire_2026}. Without finer detection of small fires, change in the Alpine regime is likely to be underestimated, because the continental analyses that define current understanding consider only fires larger than 30\,ha, whereas most Alpine fires are smaller than 10\,ha \citep{beloiu_climatedriven_2026}. Converting the existing hazard surfaces into risk requires an arc-wide exposure and vulnerability register, validated against Alpine damage as it occurs. The building-scale indices that do exist disagree on terrain slope, the variable that most distinguishes the setting \citep{papathomakohle_towards_2025, papathomakohle_wildfire_2022}.

Once those records exist, the modelling questions become answerable. Under foehn, valley-wind and channelling regimes, valley fire weather departs from interpolated fields. Whether correcting for this improves danger prediction is testable, although the size of the departure has not yet been quantified \citep{ozaki_impact_2022, sharples_overview_2009}. Spread-model evaluation would be most useful if it concentrated on intensity rather than perimeter. Intensity drives the cascade, yet the two simulators compared on Alpine fires agree on propagation and disagree on intensity \citep{malanchini_modeling_2025}. A review of more than 150 experimental and modelling studies also finds that empirical formulations using simple wind and slope adjustment factors degrade outside their calibration envelope, including on steep slopes \citep{hayajneh_wind_2026}. Reconnecting the Alpine tradition of dynamic vegetation modelling to hazard assessment would let fuel condition, rather than fuel type, enter the models \citep{schumacher_relative_2006, hillebrand_contrasting_2023, muller_towards_2020}. The coupling between fire, snow and glaciers belongs to the same programme. It is the one loop in this system that no Alpine study has yet followed through. Drought shifts summer water supply onto glaciers \citep{leone_20222023_2026}. On one Canadian glacier, soot deposited in the years after record fires increased melt by up to 10\%, although smoke itself reduced net radiation \citep{aubrywake_fire_2022}. In the western United States, fire also alters how much snow is stored and when it melts \citep{koshkin_wildfire_2022}.

The questions that matter most to managers are, by contrast, empirical, and none yet has an Alpine answer. The most tractable is probably the acceptable degree of thinning in a protection forest, that is, the point at which protective capacity declines faster than fire hazard falls. The protective side of that trade-off can already be quantified \citep{costa_biological_2021} and monetised \citep{accastello_asforesee_2019, moos_climate_2021} with methods in use. In Andalusian data, fuel breaks contained fire mainly where suppression crews worked on them, and ridgeline breaks performed best where fire met narrow breaks at an angle \citep{ortega_modeling_2024}. Steep Alpine terrain may prevent crews from reaching and holding a break, and whether any siting resolves this tension remains untested in the range. The single Alpine prescribed-burning trial (Section~\ref{sec:fuelbreaks}) was carried out in one stand of a fire-tolerant species. Its authors call for higher intensities, repeated burns, more than one location and observation beyond the first year \citep{gamba_prescribed_2025}, and in our view it also needs extending to fire-intolerant stands. The Pyrenean experience suggests that a participatory, goal-setting process may matter as much as further ecological evidence \citep{oliveres_path_2025}. No Alpine post-fire salvage trial exists, and the evidence for retaining deadwood comes from windthrow \citep{costa_biological_2021, seben_impact_2024}. It may not transfer, because burnt logs and snags may not provide the same barrier effect as windthrown stems \citep{costa_biological_2021}. Nor has anyone estimated, with terrain taken into account, the marginal return on suppression capability compared with fuel treatment. Without that estimate, the allocation implied by current spending cannot be evaluated \citep{cardil_assessing_2025, brodie_forest_2024, muller_forest_2020}.

Some gaps fall outside the usual framing and are easily overlooked for that reason. The faunal component of Alpine fire loss has been measured only in part. Bird communities have been followed for 13 years after fire in one inner-Alpine valley, and earlier Swiss work addressed arthropods \citep{rey_effects_2019}. The response of soil fauna, and the trajectory of habitat quality through the recovery period across the range, have not yet been described \citep{degiuli_arrested_2024, arrogantefunes_mapping_2024}. Extensive grazing reduces fuel at effect sizes measured elsewhere, but its effect reverses above a threshold not yet located for any Alpine pasture. Neither the unpaid fuel-management service that grazing supplies nor the pasture and livestock exposed to fire has been valued in the Alps, and a recent review calls for cost--benefit analyses comparing grazing with other fuel treatments \citep{xu_how_2026}. The distribution of costs and protective benefits across valley communities, through-traffic and downstream water users also remains unquantified, although the valuation methods to address it already exist \citep{accastello_asforesee_2019, moos_climate_2021}.

\label{sec:limitations}
These findings should be read within the limits of the evidence base, the most basic of which concerns the composition of the corpus. We built the corpus from two starting documents, a keyword search and one iteration of citation snowballing, and a single iteration inherits the blind spots of its seed documents. The keyword search of Step~2 is documented by its search strings, because Google Scholar ranks results by relevance and its result sets cannot be exported for exact reproduction. Selection was made with in-house expertise in view, which favours the disciplines represented at one institute. Backward chaining favours works that the seed authors already considered relevant. Forward chaining favours works indexed in open citation infrastructure, which under-represents grey literature, national reports and publications in languages other than English. This gap is material. A substantial part of Alpine fire knowledge exists in Italian, German, French and Slovenian technical reports, and the harmonisation of Alpine fire statistics has itself been driven by individual projects \citep{valese_wildfires_2011}. Restricting the expansion phase to sources whose full text could be retrieved widens the gap, since it biases the added corpus towards open-access and therefore recent publications.

Because the synthesis is thematic, the studies cannot be pooled quantitatively, for the definitional and administrative reasons set out in Section~\ref{sec:scope}. The numbers we report from individual studies therefore illustrate magnitudes rather than estimate Alpine parameters.

A larger limitation is the weight carried by single cases. Several quantities central to this review come from one documented site or event. The Alpine post-fire rockfall record rests on one fire, and the numerically analysed Alpine post-fire debris flow on one catchment \citep{melzner_brief_2019, laporta_postwildfire_2025}. The monetised protective function comes from one protection forest, and the arrested subalpine regeneration of Section~\ref{sec:forest} from one massif \citep{moos_climate_2021, degiuli_arrested_2024}. We have marked these in the text and treated their magnitudes as observed values rather than central estimates. A synthesis built on such a base can show that a mechanism operates in the Alps and indicate an order of magnitude, but it cannot yet support a regional parameter. Confidence in the Alpine cascade argument therefore depends on whether these few observations are representative rather than exceptional. This is why systematic post-fire monitoring forms the foundation of the research programme outlined above. The comparative ranges help, because agreement between an Alpine observation and an independent one from another range is some evidence that the Alpine case is not an outlier. They cannot, however, replace Alpine measurement.

Imported evidence raises a related difficulty. Several mechanisms central to our argument are established mainly outside the Alps, in laboratories or in other mountain ranges. They include terrain effects on the surface wind field, wind and slope effects on rate of spread, post-fire sediment supply and terrain-conditioned suppression difficulty \citep{sharples_overview_2009, hayajneh_wind_2026, morell_sediment_2021, cardil_assessing_2025}. We have used such work to show that a process operates and to describe its mechanism, never to quantify an Alpine magnitude. That distinction is easier to state than to enforce, and the Alpine magnitude of these processes is better treated as unmeasured than as inferred. Comparative mountain research programmes have argued for this kind of cross-range synthesis while being explicit about its limits \citep{payne_nature_2020}. The balance struck here is a matter of judgement rather than a rule.

% =======================================================
\section{Conclusions}
\label{sec:conclusions}

This review asked what makes wildfire in mountains different, taking the European Alps as the focal case and other ranges as comparison. Our answer is that mountain wildfire is not a smaller version of the Mediterranean problem but a different kind of problem. The differences appear at every stage, from how fires start to who bears their costs.

The differences begin with the fire regime. Alpine fire is bimodal, with a spring season of human-caused surface fires and a summer season in which lightning causes up to half of ignitions in some areas. Most fires are smaller than a hectare \citep{muller_forest_2020} and burn in stands dominated by species with little fire adaptation, such as silver fir, Norway spruce and beech, whose bark insulates poorly \citep{valese_wildfires_2011, mutterer_bark_2026}. This regime is changing faster than its records can track. Across European mountains, fires larger than 30\,ha have moved upslope by roughly 72\,m per decade since 2000, a shift led by the Balkans, Iberia and Turkey and still weak in the Alps \citep{beloiu_climatedriven_2026}. Fire hazard is projected to rise most in the northern and eastern parts of the Alpine region, which have so far been little exposed to fire \citep{gerberding_mapping_2025}. Land abandonment has also extended forest cover and fuel connectivity \citep{gerberding_mapping_2025}, creating continuous fuel on the long, steep slopes from which large fires develop \citep{conedera_linking_2024}. Because this fuel comes from land use rather than from fire exclusion, agricultural and pastoral policy can influence it as much as silviculture can.

These trends matter because, in mountains, the damage a fire causes does not scale with the area it burns. Many Alpine forests are managed to hold rock, snow and sediment above settlements and infrastructure. The asset at risk therefore usually lies below the fire rather than within it. In the two stands where it has been valued, rockfall protection was worth about 950 euros and 1{,}350 Swiss francs per hectare and year, while the interventions that maintain it ran at a loss in the one case that costed them \citep{accastello_asforesee_2019, moos_climate_2021}. Loss then follows several pathways at once. The pathway that most clearly sets mountain fire apart is the geomorphic cascade, in which a burnt catchment responds to rainfall that an unburnt one would absorb. It is one pathway among several, however, and not necessarily the largest. Evidence from Colorado suggests that the triggering threshold recovers as vegetation returns \citep{rengers_evaluating_2024}, so we expect the duration of the post-fire window to matter as much as how far the threshold first falls; no Alpine record yet tracks that recovery. Protective structure is also lost over decades \citep{gehring_shallow_2019}, ecological change harms closed-forest species while favouring those of open habitats, and the effects on animals remain unmeasured.

The tools available to assess these consequences still describe hazard rather than risk. Danger indices, spread models and severity products were built and calibrated where fires are larger, terrain gentler and records longer, and they consistently lose skill in mountains. Adding complexity is unlikely to help, because the Alpine record is too short and sparse to constrain more parameters. Progress is more likely to come from well-chosen model structure, and from importing mechanisms from elsewhere while refitting their coefficients locally. Assembling exposure and vulnerability data would then allow existing hazard maps to be turned into estimates of risk. Because evidence from other ranges mostly agrees with the Alpine record, these differences appear to belong to mountains in general rather than to the Alps alone.

Responding to mountain wildfire is therefore largely a matter of forestry, land use and institutions. Prevention costs much less than the firefighting and post-fire response it would replace (Section~\ref{sec:economics}). The value it protects, however, is a public service, and those who benefit from it are rarely those who would pay for it. The post-fire period, when that value is most exposed, also falls between institutions rather than within one. Unlike regions with long histories of suppression, the Alps are still shaping their approach to fire and could base it on adaptation and coexistence from the start. Whether they do so will probably depend less on new models than on systematic observation, a defined handover between institutions after fire, and land-use decisions, all of which the evidence shows to be missing.


\section*{Acknowledgements}
\addcontentsline{toc}{section}{Acknowledgements}
% Text here.

\section*{Funding}
\addcontentsline{toc}{section}{Funding}
% Text here.

\section*{Competing Interests}
\addcontentsline{toc}{section}{Competing Interests}
% Text here.

\section*{Data and Code Availability}
\addcontentsline{toc}{section}{Data and Code Availability}
% The screening register underlying this review (all candidate records, tier,
% inclusion decision and reason) is provided as supplementary material.


% =======================================================
\appendix
% =======================================================

\section{Search Protocol and Screening Register}
\label{app:protocol}

\subsection{Counts at each step}

The six steps of Section~\ref{sec:scope} and Figure~\ref{fig:flow} are given below with the numbers each produced. We chose starting documents, a keyword search and citation chaining because the literature is administratively fragmented, multilingual and heavily dependent on national technical reports. Database searches restricted to English, of the kind used in a recent global review of land-abandonment effects, acknowledge that they may miss part of such a literature \citep{mantero_influence_2020}.

\textbf{Step 1, starting documents.} The EUSALP white paper \emph{Forest fires in the Alps: state of knowledge, future challenges and options for an integrated fire management} \citep{muller_forest_2020} covers the whole Alpine Convention area. It combines harmonised fire statistics for the arc, a survey of authorities, scientists and fire services across the Alpine countries, and costed recommendations for prevention, suppression, post-fire management and knowledge transfer. Its reference list contains 120 publications and web sources, and it draws on six national fire databases. The study co-authored at Eurac Research \citep{gerberding_mapping_2025} maps the probability of forest fire hazard across the arc under climate change scenarios.

\textbf{Step 2, keyword search.} The Google Scholar searches combined the terms wildfire, mountain and Alps, and then wildfire and mountains. Publications were kept where they fell within a discipline represented at Eurac Research, namely natural hazards and risk, climate, drought and water, forest science, and economics and governance; these are also the disciplines of the co-authors who review the sections in Step~6. The initial set holds 49 publications, including the Eurac Research study, and 47 of them were available in full text. Of the 49, 31 were published in 2020 or later; three also appear in the white paper's reference list and six cite it.

\textbf{Step 3, citation chaining.} Working backward, we compiled the reference lists of the white paper and of 45 of the initial publications (two had no usable reference list), 3{,}575 references in total, 120 of them from the white paper. After duplicates were removed, 3{,}375 distinct publications remained. Of the white paper's 120 references, 62 were also reached through the initial publications and 58 through the white paper alone. Working forward, we looked up the publications citing the 14 initial publications most central to the Alpine question in the OpenCitations index. This returned 508 citing publications, of which 319 could be identified with full bibliographic details; the other 189 are kept in the register for a later search.

\textbf{Step 4, targeted additions.} Additions were limited to publications whose full text could be obtained and archived. This step therefore weights the literature towards open-access and recent publications, and it is the step that can least be reproduced from the initial set alone.

\textbf{Step 5, screening.} After duplicates across all sources were removed, 3{,}228 publications were screened. Each was assigned to one of four levels, with the reason recorded: \emph{include} (fire content and the European Alps or an Alpine Convention country); \emph{comparative} (fire content and another mountain system, admitted on the strength of the process, following the two-tier rule of Section~\ref{sec:scope}); \emph{context} (mountain processes without fire, or fire without a mountain setting); and \emph{exclude} (general fire science with no mountain anchor). Screening excluded 2{,}332 publications and retained 896: 290 at the include level, 194 comparative and 412 context. The retained publications were then read in order of their relevance to the thematic sections, and 150 of them are cited in the text, each summarised in an evidence note (Section~\ref{app:recorded}).

\textbf{Step 6, discipline review.} Climatologists review the climate and meteorology evidence of Section~\ref{sec:climate}, drought and water researchers Section~\ref{sec:water}, and forest scientists Section~\ref{sec:forest}. Specialists in natural hazards and risk review Sections~\ref{sec:hazard} and~\ref{sec:risk}, and economists and governance researchers review Section~\ref{sec:governance}. Publications added by the reviewers go through Steps~3 and~5 like the rest of the literature.

\subsection{What was recorded for each source}
\label{app:recorded}

Each screened record carries a tier, a decision, a reason and a mapping to one or more sections of this review. For every source cited in the text, a structured evidence note recording study region, method, quantitative results and stated limitations was captured and archived alongside the screening register. Both are provided as supplementary material. Where a claim in the text rests on fewer than three independent cases, the number is stated in the text, and mechanisms imported from outside the Alps are identified as imported at the point of use.

Synthesis is thematic rather than quantitative, for the reasons of definitional and administrative heterogeneity set out in Section~\ref{sec:scope}.

\subsection{What the method determines about the corpus}

The method shapes what the corpus contains, and its effects should be read as constraints on coverage rather than on the reliability of individual sources. Citation chaining reproduces the coverage of its starting point, so topics well represented in the white paper and among the seeds are recovered densely and topics absent from both are recovered weakly. Forward chaining operates on open citation infrastructure, which indexes journal literature comprehensively but national technical reporting and non-English publication sparsely, a material consideration for a range whose fire statistics were harmonised only recently \citep{valese_wildfires_2011}. The expansion phase was confined to material whose full text could be retrieved, which weights the added corpus towards open-access and recent publication.

A subsequent iteration should add database searching with recorded strings and dates, resolve the 189 forward records held over, and target the national fire and forest reporting of the Alpine Convention countries directly, where much of the operational and grey-literature knowledge base resides.

\section{Glossary of Terms}

\begin{description}
    \item[Fire danger] The meteorological and fuel-condition potential for ignition and spread, expressed by an index; a statement about conditions, not about realised fire.
    \item[Fire regime] The realised distribution of fire frequency, size, severity, seasonality and cause over a region and period.
    \item[Burn severity] The degree of fire-induced change to vegetation and soil, commonly estimated remotely; see Section~\ref{sec:water} for the limits of this proxy.
    \item[Fuel connectivity] The extent and spatial aggregation of contiguous burnable vegetation, distinct from fuel load.
    \item[Protection forest] Forest whose primary managed function is mitigation of gravitational natural hazards affecting assets below.
    \item[Geohydrological process] Collective term used in the Alpine literature for erosion, rockfall, slides and debris flows driven by the interaction of water and slope material.
    \item[Cascade] A sequence in which one hazard alters the susceptibility of a system to a subsequent, distinct hazard.
    \item[Firebreak] A linear strip from which fuel is removed entirely, intended to stop fire spread.
    \item[Fuel break (\emph{coupure de combustible})] A linear strip in which fuel load and vertical continuity are reduced but the stand is retained; functions primarily as a platform for ground suppression rather than as a barrier.
    \item[Exposure] The elements (people, buildings, infrastructure, habitats, protective forest, livestock and pasture) located where a hazard and its consequences can reach them.
    \item[Vulnerability] The susceptibility of an exposed element to loss given a hazard of stated intensity; distinct from exposure and from hazard, and measured by damage-calibrated or judgement-weighted indicators.
    \item[Ecological vulnerability] Potential loss of ecological value integrated over the post-fire recovery period, combining the value of what burns, the resistance of the vegetation to fire and the time taken to return to pre-fire condition.
\end{description}


% =======================================================
% Bibliography
% =======================================================
\bibliographystyle{apalike}
\bibliography{FireScape_master_v7}

\end{document}
